\RequirePackage[thinlines]{easybmat} %-- muss aufgrund eines Bugs vor etex und tikz geladen werden
\documentclass[a4paper, twoside, headsepline, index=totoc,toc=listof, fontsize=10pt, cleardoublepage=empty, headinclude, DIV=13, BCOR=13mm, titlepage,]{scrartcl}
\usepackage{scrtime} % KOMA, Uhrzeit ermoeglicht
\usepackage{scrpage2} % wie fancyhdr, nur optimiert auf KOMA-Skript, leicht andere Syntax
\usepackage{etoolbox}
\usepackage{letltxmacro}

% Compiler abhaengige IFs
\usepackage{ifxetex,ifluatex}
\newif\ifxetexorluatex
\ifxetex
  \xetexorluatextrue
\else
  \ifluatex
    \xetexorluatextrue
  \else
    \xetexorluatexfalse
  \fi
\fi

%--Farbdefinitionen
%-- muss vor tikz geladen werden
\usepackage[usenames, table, x11names]{xcolor}
\definecolor{dark_gray}{gray}{0.45}
\definecolor{light_gray}{gray}{0.6}

%--Zum Zeichnen
%-- muss vor polyglossia bzw. babel geladen werden
\usepackage{tikz}
\usepackage{tikz-cd}
\usetikzlibrary{external}
\tikzset{>=latex}
\usetikzlibrary{shapes,arrows,intersections}
\usetikzlibrary{calc,3d}
\usetikzlibrary{decorations.pathreplacing,decorations.markings}
\usetikzlibrary{angles}

%-- Konfiguration von tikzexternalize
\tikzexternalize[prefix=tikz/, up to date check=diff]

\ifxetexorluatex
\tikzset{external/system call={xelatex \tikzexternalcheckshellescape %-- verwende LuaLaTeX, wegen dynamischer Speicherallokation
    -halt-on-error -interaction=batchmode -jobname "\image" "\texsource"}}
\else
\tikzset{external/system call={pdflatex \tikzexternalcheckshellescape 
    -halt-on-error -interaction=batchmode -jobname "\image" "\texsource"}}
\fi

%-- tikzexternalize fuer tikzcd deaktivieren, da inkompatibel
\AtBeginEnvironment{tikzcd}{\tikzexternaldisable}
\AtEndEnvironment{tikzcd}{\tikzexternalenable}

%-- um Inkompatibilitaeten von quotes und polyglossia bzw. babel zu vermeiden
\tikzset{
  every picture/.append style={
    execute at begin picture={\shorthandoff{"}},
    execute at end picture={\shorthandon{"}}
  }
}
\usetikzlibrary{quotes}
\usepackage{pgfplots}
\usepgfplotslibrary{colormaps}
\pgfplotsset{compat=1.10}



%-- Mathe
\usepackage{mathtools} % beinhaltet amsmath
\mathtoolsset{showonlyrefs, centercolon}
\newtagform{brackets}[\textbf]{[}{]}
\usetagform{brackets}
\usepackage{wasysym}
\usepackage{amssymb} %zusätzliche Symbole
\usepackage{latexsym} %zusätzliche Symbole
\usepackage{stmaryrd} %für Blitz
\usepackage{nicefrac} %schräge Brüche
\usepackage{cancel} %Befehle zum Durchstreichen
\usepackage{mathdots}
% \usepackage{bm}
\DeclareSymbolFont{bbold}{U}{bbold}{m}{n}
\DeclareSymbolFontAlphabet{\mathbbold}{bbold}
\newcommand{\mathds}[1]{\mathbb{#1}} %Um Kompatibilitaet mit frueheren Benutzung von dsfont herzustellen
\newcommand{\ind}{\mathbbold{1}} % charakteristische-Funktion-Eins
\def\mathul#1#2{\color{#1}\underline{{\color{black}#2}}\color{black}} %farbiges Untersteichen im Mathe-Modus

%-- Alles was mit Schrift und XeteX zu tun hat
\ifxetexorluatex
    % XeLaTeX or LuaTeX
	\ifxetex
		\usepackage{mathspec} %beinhaltet fontspec 
	\else
		\usepackage[no-math]{fontspec}
	\fi
    \usepackage{polyglossia} %babel-ersatz
    \setmainlanguage[spelling=new,babelshorthands=false]{german}
	\newcommand\glqq{"}
	\newcommand\grqq{"}
	\defaultfontfeatures{Mapping=tex-text, WordSpace={1.4}} %
    \setmainfont[Ligatures=Common, BoldFont={* Bold}, ItalicFont={* Light Italic},SmallCapsFont={Linux Libertine O}, SmallCapsFeatures={Letters=SmallCaps}]{Source Sans Pro}
	\setsansfont[Scale=MatchLowercase,Ligatures=Common, BoldFont={* Medium}]{Ubuntu}
	\ifxetex
		\setallmonofonts[Scale=MatchLowercase, ItalicFont={*}]{Consolas} 
	\else
		\setmonofont[Scale=MatchLowercase, ItalicFont={*}]{Consolas}
	\fi
	\usepackage{xltxtra}
	\usepackage{fontawesome}
	\usepackage{microtype}
\else
    % default: pdfLaTeX
    \usepackage[ngerman]{babel}
    \usepackage[T1]{fontenc}
    \usepackage[utf8]{inputenc}
	\usepackage{textcomp} %verhindert ein paar Fehler bei den Fonts
    \usepackage[babel=true, tracking=true]{microtype}
	\usepackage{lmodern}
	\renewcommand{\familydefault}{\sfdefault} 
\fi


%--Mixed
\usepackage[neverdecrease]{paralist}
\usepackage[german=quotes]{csquotes}
\usepackage{makeidx}
\usepackage{booktabs}
\usepackage{wrapfig}
\usepackage{float}
\usepackage[margin=10pt, font=small, labelfont={sf, bf}, format=plain, indention=1em]{caption}
\captionsetup[wrapfigure]{name=Abb. }
\usepackage{stackrel}
\usepackage{ifthen}
\usepackage{multicol}
\usepackage{qrcode}
\flushbottom

%--Hyphenation
\hyphenation{di-men-si-o-na-le}


%--Unterstreichung
\usepackage[normalem]{ulem}
\setlength{\ULdepth}{1.8pt}

%--Indexverarbeitung
\newcommand{\bet}[1]{\textbf{#1}}
\newcommand{\Index}[1]{\textbf{#1}\index{#1}}
\makeindex
\setindexpreamble{{\noindent \itshape Die \emph{Seitenzahlen} sind mit Hyperlinks zu den entsprechenden Seiten versehen, also anklickbar} \ifxetexorluatex \faHandUp\fi \par \bigskip}
\renewcommand{\indexpagestyle}{scrheadings}


%--Marginnote & todonotes
\usepackage{marginnote}
\renewcommand*{\marginfont}{\color{dark_gray} \itshape \footnotesize }
\usepackage[textsize=small]{todonotes}

%--Todonotes mit tikz/externalize kompatibel machen
\LetLtxMacro{\oldtodo}{\todo}
\renewcommand{\todo}[2][]{\tikzexternaldisable\oldtodo[#1]{#2}\tikzexternalenable}
\LetLtxMacro{\oldmissingfigure}{\missingfigure}
\renewcommand{\missingfigure}[2][]{\tikzexternaldisable\oldmissingfigure[{#1}]{#2}\tikzexternalenable}

%--Konfiguration von Hyperref pdfstartview=FitH, 
\usepackage[hidelinks, pdfpagelabels,  bookmarksopen=true, bookmarksnumbered=true, linkcolor=black, urlcolor=SkyBlue2, plainpages=false,pagebackref, citecolor=black, hypertexnames=true, pdfauthor={Jannes Bantje}, pdfborderstyle={/S/U}, linkbordercolor=SkyBlue2, colorlinks=false]{hyperref}

% \let\orighref\href
\ifxetexorluatex
\newcommand{\hrefsym}[2]{\href{#1}{\texttt{#2\,\faExternalLink}}}
\newcommand{\hrefsymmail}[2]{\href{#1}{\texttt{\faEnvelope\,#2}}}
\renewcommand{\url}[1]{\hrefsym{#1}{\nolinkurl{#1}}}
\fi
\providecommand{\hrefsym}{\href}
% \usepackage{cleveref}
% \crefformat{footnote}{#2\footnotemark[#1]#3}

%--Römische Zahlen
\newcommand{\RM}[1]{\MakeUppercase{\romannumeral #1{}}}

%--Punkte (nach hyperref)
\usepackage{ellipsis}

%%--Abkürzungen etc.
\newcommand{\light}{\text{\Large $\lightning$}}

%-- Definitionen von weiteren Mathe-Befehlen
\DeclareMathOperator{\re}{Re} %Realteil
\DeclareMathOperator{\im}{Im} %Imaginaerteil
\DeclareMathOperator{\id}{id} %identische Abbildung
\DeclareMathOperator{\Sp}{Sp} %Spur
\DeclareMathOperator{\sgn}{sgn} %Signum
\DeclareMathOperator{\alt}{Alt} %Alternierende n-linearFormen
\DeclareMathOperator{\End}{End} %Endomorphismen
\DeclareMathOperator{\Vol}{Vol} %Volumen
\DeclareMathOperator{\dom}{dom} %Domain
\DeclareMathOperator{\Hom}{Hom} %Homomorphismen
\DeclareMathOperator{\bild}{Bild} %Bild
\DeclareMathOperator{\Kern}{Kern }
\DeclareMathOperator{\rg}{rg} %Rang
\DeclareMathOperator{\Rg}{Rg} %Rang
\DeclareMathOperator{\diam}{diam} %Durchmesser
\DeclareMathOperator{\dist}{dist} %Distanz
\DeclareMathOperator{\grad}{grad} %Gradient
\DeclareMathOperator{\dive}{div} %Gradient
\DeclareMathOperator{\rot}{rot} %Rotation
\DeclareMathOperator{\hess}{Hess} %Hesse-Matrix
\DeclareMathOperator{\koker}{Koker} %Kokern
\DeclareMathOperator{\aut}{Aut} %Automorphismen
\DeclareMathOperator{\ord}{ord} %Ordnung
\DeclareMathOperator{\ggT}{ggT} %ggT
\DeclareMathOperator{\kgV}{kgV} %kgV
\DeclareMathOperator{\Gr}{Gr} %Gerade
\DeclareMathOperator{\Kr}{Kr} %Kreis
\DeclareMathOperator{\Char}{char} %Charakteristik
\DeclareMathOperator{\Aut}{Aut} %Automorphismen
% \DeclareMathOperator{\D}{D} %formale Ableitung
\newcommand{\D}{\ensuremath{\mathrm{D}\mkern-1.5mu}}
\newcommand{\mathd}{\ensuremath{\mathrm{d}\mkern-1.0mu}}
\newcommand{\Tmap}{\ensuremath{\mathrm{T}\mkern-1.0mu}}
\newcommand{\opL}{\ensuremath{\mathrm{L}\mkern-0.6mu}}
\DeclareMathOperator{\cov}{cov} %Kovarianz
\DeclareMathOperator{\Gal}{Gal} %Galois
\DeclareMathOperator{\supp}{supp} %Träger
\DeclareMathOperator{\Sym}{Sym} %Symmetrische Gruppe
\DeclareMathOperator{\Zyl}{Zyl} %Zylinder
\DeclareMathOperator{\Mod}{Mod} %Moduln
\DeclareMathOperator{\EPK}{EPK} %Einpunktkompaktifizierung
\DeclareMathOperator{\conj}{conj} %Konjugation
\DeclareMathOperator{\ann}{ann} %Annulator
\DeclareMathOperator{\tor}{tor} %Torsionsmodul
\DeclareMathOperator{\Int}{Int} %Interior/Inneres
\DeclareMathOperator{\Br}{Br} %Brauerergruppe
\DeclareMathOperator{\GL}{GL} %allgemeine lineare Gruppe
\DeclareMathOperator{\SL}{SL} %spezielle lineare Gruppe
\DeclareMathOperator{\SU}{SU} %spezielle unitäre Gruppe
\DeclareMathOperator{\Tr}{Tr} %Spur/Trace
\DeclareMathOperator{\conv}{conv} %Konvexe Teilmenge
% \DeclareMathOperator{\tr}{tr}
\newcommand{\tr}{\mathrm{tr}}
\newcommand{\w}{\mathrm{w}}
\newcommand{\weakT}[1]{\ensuremath{\mathcal{T}_{#1}^{\mkern+1.0mu\text{\raisebox{0.4ex}{$\mathrm{w}$}}}}}
\newcommand{\weakTstar}[1]{\ensuremath{\mathcal{T}_{#1}^{\mkern+1.0mu\text{\raisebox{0.4ex}{$\mathrm{w}$}}^*}}}
\newcommand{\rand}[1]{\ensuremath{\partial^{\scriptscriptstyle #1}}}


% -- Kategorien
\DeclareMathOperator{\Ob}{Ob} %Objekte
\DeclareMathOperator{\Mor}{Mor} %Morphismen

% \DeclareMathOperator{\TOP}{TOP} %Topologische Räume
\newcommand{\TOP}{\textsc{Top}}
% \DeclareMathOperator{\HTOP}{HTOP} %Topologische Räume + Homotopieäquivalenz
\newcommand{\HTOP}{\textsc{HTop}}
% \DeclareMathOperator{\VR}{VR} %Vektorräume
\newcommand{\VR}{\textsc{VR}}
% \DeclareMathOperator{\MOD}{MOD} %Moduln
\newcommand{\MOD}{\textsc{Mod}}
% \DeclareMathOperator{\MENGEN}{MENGEN} %Mengen
\newcommand{\MENGEN}{\textsc{Mengen}}
% \DeclareMathOperator{\MAN}{MAN} %Mannigfaltigkeiten
\newcommand{\MAN}{\textsc{Man}}
% \DeclareMathOperator{\GRUPPEN}{GRUPPEN} %Gruppen
\newcommand{\GRUPPEN}{\textsc{Gruppen}}
% \DeclareMathOperator{\ABELGRUPPEN}{ABEL.GRUPPEN} %abelsche Gruppen
\newcommand{\ABELGRUPPEN}{\textsc{Abel.Gruppen}}
% \DeclareMathOperator{\KAT}{KAT} %Kategorien
\newcommand{\KAT}{\textsc{Kat}}
% \DeclareMathOperator{\FUN}{FUN} %Funktoren
\newcommand{\FUN}{\textsc{Fun}}


\newcommand{\op}{\mathrm{op}}

\newcommand{\Ker}{\ensuremath{\Kern \,}}
\newcommand{\zirkel}{\rotatebox[origin = c]{-90}{$\varangle$}}
\newcommand{\wirkung}{\!\curvearrowright\!}

\newcommand{\oE}{\OE}

%--Beweisende
\newcommand{\bewende}{\ifmmode \tag*{$\square$} \else \hfill $\square$ \fi}

%--Volumen
\newcommand{\vol}[1]{\ensuremath{\Vol_{#1}}}

%--Skalarprodukt
\DeclarePairedDelimiterX\sprod[2]{\langle}{\rangle}{#1\,\delimsize\vert\,#2}

%--Betrag, Gaußklammer
\DeclarePairedDelimiter{\abs}{\lvert}{\rvert}
\DeclarePairedDelimiter{\ceil}{\lceil}{\rceil}

%--Norm
\DeclarePairedDelimiter\doppelstrich{\Vert}{\Vert}
\newcommand{\norm}[2][\relax]{
\ifx#1\relax \ensuremath{\doppelstrich*{#2}}
\else \ensuremath{\doppelstrich*{#2}_{#1}}
\fi}


%--Umklammern
\DeclarePairedDelimiter\enbrace{(}{)}
\DeclarePairedDelimiter\benbrace{[}{]}
\DeclarePairedDelimiter\lenbrace{<}{>}

%--Mengen
\DeclarePairedDelimiterX\mengenA[1]{\lbrace}{\rbrace}{#1}
\DeclarePairedDelimiterX\mengenB[2]{\lbrace}{\rbrace}{#1\, \delimsize\vert \, #2}

\makeatletter
\newcommand{\set}[2][\relax]{
\ifx#1\relax \ensuremath{
\mengenA*{#2}}
\else \ensuremath{%
  \mengenB*{#1}{#2}}
\fi}
\makeatother


%--offen und abgeschlossen
\newcommand{\off}{\! \stackrel[\text{offen}]{}{\subset} \!}
\newcommand{\abg}{\! \stackrel[\text{abg.}]{}{\subset} \!}

%--Differential
\newcommand{\diff}[2]{\ensuremath{\frac{{\partial #1}}{{\partial #2}} }}
\newcommand{\diffd}[2]{\ensuremath{\frac{\mathd #1}{\mathd #2} }}

%--Diagonale Linien für Matrizen
\newcommand\tikzmark[1]{%
  \tikzexternaldisable\tikz[overlay,remember picture,baseline] \node [anchor=base] (#1) {};\tikzexternalenable}

\newcommand\MyLine[3][]{%
  \tikzexternaldisable\begin{tikzpicture}[overlay,remember picture]
    \draw[#1] (#2.south) -- (#3.north);
  \end{tikzpicture}\tikzexternalenable}
  
%--Jordankasten
\newcommand{\Jord}{ \tikzexternaldisable \begin{BMAT}{|ccc|}{|ccc|} %
				0 & & \\ %
				1\tikzmark{z} &  \diagdown & \\ %
				& \tikzmark{u}1 \MyLine[thick]{z}{u} &  0 %
			\end{BMAT} \tikzexternalenable}
			
%--Jordankasten mit Lambda
\newcommand{\JordLa}[1][\empty]{%
	\tikzexternaldisable\ifthenelse{\equal{#1}{\empty}}
		{\begin{BMAT}{|ccc|}{|ccc|} %
				\lambda & & \\ %
				1\tikzmark{z} &  \diagdown & \\ %
				& \tikzmark{u}1 \MyLine[thick]{z}{u} &  \lambda  %
			\end{BMAT}}
		{\begin{BMAT}{|ccc|}{|ccc|} %
				\lambda_{#1}  & & \\ %
				1\tikzmark{z} &  \diagdown & \\ %
				& \tikzmark{u}1 \MyLine[thick]{z}{u} &  \lambda_{#1}  %
			\end{BMAT}}\tikzexternalenable}

%--Inhaltsverzeichnis
\usepackage[tocindentauto]{tocstyle}
\usetocstyle{KOMAlike}			

\newcommand{\fach}{Einführung in die Funktionalanalysis}
\newcommand{\semester}{WiSe 2014}
\newcommand{\homepage}{https://wwwmath.uni-muenster.de/u/wilhelm.winter/wwinter/funktionalanalysis.html}
\newcommand{\prof}{Prof.\,Dr.\,Wilhelm Winter}
\providecommand{\verfasser}{Jannes Bantje}

%--Konfiguration von scrheadings
\setheadsepline{1pt}[\color{light_gray}]
\pagestyle{scrheadings}
%\fancyhf[H, F]{}
\clearscrheadfoot

\providecommand{\shortFach}{\fach}
\lehead{ \includegraphics[height=0.5 cm,keepaspectratio]{../!config/Bilder/Logo_WWU_Muenster_light_gray.pdf}}
\rehead{\rule{0cm}{0.5cm}\footnotesize \sffamily \color{light_gray} \verfasser -- Mitschrift \shortFach}
\lohead{\rule{0cm}{0.5cm} \footnotesize \sffamily \color{light_gray} Stand: \today \; \thistime[:]}
\rohead{\includegraphics[height=0.5 cm,keepaspectratio]{../!config/Bilder/fb10logo_gray.pdf}}


\ofoot[{ \color{dark_gray} \LARGE \sffamily \thepage}]{{ \color{dark_gray} \LARGE \sffamily \thepage}} %hier wir auch der plain Stil bearbeitet!
\automark{section}
\ifoot{ \color{dark_gray} \small \leftmark}

%--Metadaten
\author{\verfasser}
\titlehead{\includegraphics[height=1.5cm, keepaspectratio]{../!config/Bilder/Logo_WWU_Muenster.pdf}%
\hfill \includegraphics[height=1.3cm, keepaspectratio]{../!config/Bilder/fb10logo.pdf}}
\title{Skript \fach}
\subtitle{Mitschrift der Vorlesung  \enquote{\fach} von \prof}

\publishers{
Aktuelle Version verfügbar bei: \bigskip\\
\normalsize
\begin{minipage}[t]{0.45\textwidth}
	\centering{
	\qrcode[height=3.5cm, version=4]{https://github.com/JaMeZ-B/latex-wwu}
	\medskip\\
	\includegraphics[height=0.4cm, keepaspectratio]{../!config/Bilder/github_octo.pdf} 
	\includegraphics[height=0.4cm, keepaspectratio]{../!config/Bilder/GitHub_Logo.pdf} (inklusive Sourcecode)\\
	{\small\url{https://github.com/JaMeZ-B/latex-wwu}}}
\end{minipage}
\hfill
\begin{minipage}[t]{0.45\textwidth}
	\centering{
	\qrcode[height=3.5cm, version=4]{btsync://B6WH2DISQ5QVYIRYIEZSF4ZR2IDVKPN3I?n=latex_share}
	\medskip\\
	\includegraphics[height=0.4cm, keepaspectratio]{../!config/Bilder/bt_sync_logo.pdf} 
	{\large \textbf{Bittorrent} Sync} \\
	\texttt{B6WH2DISQ5QVYIRYIEZSF4ZR2IDVKPN3I}
	}
\end{minipage}
}



\begin{document}
\pagenumbering{Roman}
\maketitle
\begin{abstract}
\section*{Vorwort --- Mitarbeit am Skript}
Dieses Dokument ist eine Mitschrift aus der Vorlesung \enquote{\fach, \semester}, gelesen von \prof. Der Inhalt entspricht weitestgehend dem Tafelanschrieb. Für die
Korrektheit des Inhalts übernehme ich keinerlei Garantie! Für Bemerkungen und Korrekturen -- und seien es nur Rechtschreibfehler -- bin ich sehr dankbar. 
Korrekturen lassen sich prinzipiell auf drei Wegen einreichen: 
\begin{itemize}
	\item Persönliches Ansprechen in der Uni, Mails an \hrefsymmail{mailto:\mail}{\mail} (gerne auch mit annotieren PDFs) 
	\item \emph{Direktes} Mitarbeiten am Skript: Den Quellcode poste ich auf GitHub (siehe oben), also stehen vielfältige Möglichkeiten der Zusammenarbeit zur Verfügung:
	Zum Beispiel durch Kommentare am Code über die Website und die Kombination Fork + Pull Request. Wer sich verdient macht oder ein Skript zu einer Vorlesung, die 
	ich nicht besuche, beisteuern will, dem gewähre ich gerne auch Schreibzugriff.
	
	Beachten sollte man dabei, dass dazu ein Account bei \url{github.com} notwendig ist, der allerdings ohne Angabe von persönlichen Daten angelegt werden kann. 
	Wer bei GitHub (bzw. dem zugrunde liegenden Open-Source-Programm \enquote{\texttt{git}}) -- verständlicherweise -- Hilfe beim Einstieg braucht, dem helfe ich gerne 
	weiter. Es gibt aber auch zahlreiche empfehlenswerte Tutorials im Internet.\footnote{zB. \url{https://try.github.io/levels/1/challenges/1}, ist auf Englisch, aber dafür 
	interaktives LearningByDoing}
	\item \emph{Indirektes} Mitarbeiten: \TeX-Dateien per Mail verschicken. 
	
	Dies ist nur dann sinnvoll, wenn man einen ganzen Abschnitt ändern möchte (zB. einen alternativen Beweis geben), da ich die Änderungen dann per Hand einbauen muss! Ich freue mich aber auch über solche Beiträge!
\end{itemize}
\section*{Vorlesungshomepage}
{\centering 
\begin{minipage}[c][][c]{\textwidth}
	\centering \qrcode[height=3cm]{\homepage} \medskip\\
	\footnotesize \url{\homepage}
\end{minipage}
\par}

\end{abstract}

\setcounter{page}{1}
\tableofcontents
\cleardoubleoddemptypage
\pagenumbering{arabic}
\setcounter{page}{1}

\section{Metrische Räume und der Satz von Baire} % (fold)
\label{sec:1}

\subsection[Definition: Metrischer Raum]{Definition} % (fold)
\label{sub:11}
Ein \Index{metrischer Raum} ist ein Paar $(X,d)$, wobei $X$ eine Menge und $d : X \times X \to [0,\infty)$ ist, sodass
\begin{enumerate}[1)]
	\item $d(x,y)=0 \iff x=y$ $\forall x,y \in X$
	\item $d(x,y) = d(y,x)$ $\forall x,y \in X$
	\item $d(x,z) \le d(x,y) + d(y,z)$ $\forall x,y,z \in X$
\end{enumerate}
% subsection 11 (end)

\subsection[Definition: Offen, abgeschlossen und Abschluss]{Definition} % (fold)
\label{sub:12}
Sei $(X,d)$ ein metrischer Raum
\begin{itemize}
	\item Eine Teilmenge $U \subseteq X$ heißt \Index{offen}, falls für jedes $x \in U$ ein $\varepsilon > 0$ existiert, so dass 
	\[
		B(x,\varepsilon) := \set[y \in X]{d(x,y) < \varepsilon} \subset U 
	\]
	\item Eine Teilmenge $A \subset X$ heißt \Index{abgeschlossen}, falls $X \setminus A$ offen ist (als Teilmenge von $X$). 
	\item $\mathcal{T}X := \set[U \subset X]{U \text{ offen}} $ ist die Topologie auf $X$ (die von der Metrik $d$ induziert wird)
	\item Falls $W \subset X$ eine Teilmenge ist, dann bezeichnet $\overline{W}$ den \Index{Abschluss} von $W$, d.h. die kleinste abgeschlossene Teilmenge von $X$, die 
	$W$ enthält. Es gilt
	\[
		\overline{W} = \enspace\bigcap_{\mathclap{A \subset X \text{abg.}, W \subset A}} \enspace A 
	\] 
\end{itemize}
Für metrische Räume gilt:
\(
	\overline{W} = \set[\lim_n x_n]{(x_n)_n \subset W \text{ konvergente Folge}}.
\)
Warum gilt $W \subset \overline{W}$?
% subsection 12 (end)

\subsection[Definition: Stetigkeit, gleichmäßige Stetigkeit, Isometrie]{Definition} % (fold)
\label{sub:13}
Eine Abbildung $f : X \to Y$ zwischen zwei metrischen Räumen $(X,d_X)$ und $(Y,d_Y)$ heißt 
\begin{itemize}
	\item \bet{stetig in} $x \in X$, falls
	\(
		\forall \varepsilon >0 : \exists \delta >0 : \forall x' \in X : d_X(x,x') < \delta  \Longrightarrow d_Y \enbrace[\big]{f(x), f(x')} < \varepsilon 
	\)
	\item \Index{stetig}, falls $f$ in jedem Punkt $x \in X$ stetig ist.
	
	(Äquivalent: Für jede offene Menge $V$ in $Y$ ist $f ^{-1}(V)$ offen in $X$) 
	\item \Index{gleichmäßig stetig}, falls
	\(
		\forall \varepsilon>0 : \exists \delta >0 : \forall x,x' \in X : d_X(x,x') < \delta \Longrightarrow d_Y \enbrace[\big]{f(x), f(x')} < \varepsilon .
	\) 
	\item $f$ heißt \Index{Isometrie}, falls \marginnote{automatisch injektiv}
	\(
		\forall x,x' \in X : d_X(x,x') = d_Y \enbrace[\big]{f(x), f(x')}.
	\)
	\item $f$ heißt \Index{isometrischer Isomorphismus}, falls $f$ bijektiv und isometrisch ist. \marginnote{$f ^{-1}$ auch}
\end{itemize}
% subsection 13 (end)

\subsection[Definition: Cauchy-Folge und Vollständigkeit]{Definition} % (fold)
\label{sub:14}
Eine Folge $(x_n)_{n \in \mathds{N}}$ in einem metrischen Raum $(X,d)$ heißt \bet{Cauchy-Folge}\index{Cauchy-Folge}, falls
\[
	\forall \varepsilon>0 : \exists N \in \mathds{N} : \forall n,k >N : d(x_n,x_k) < \varepsilon.
\]
$(X,d)$ heißt \Index{vollständig}, falls jede Cauchy-Folge in $X$ konvergiert. 
% subsection 14 (end)

\subsection[Satz: Existenz einer eindeutigen Vervollständigung metrischer Räume]{Satz} % (fold)
\label{sub:15}
Sei $(X,d)$ ein metrischer Raum. Dann existiert ein vollständiger metrischer Raum $(\tilde{X}, \tilde{d})$ und eine Isometrie $\iota : X \hookrightarrow \tilde{X}$, sodass
$\overline{\iota(X)} = \tilde{X}$ (d.h. $\iota(X)$ ist dicht in $\tilde{X}$). $(\tilde{X}, \tilde{d})$ heißt \Index{Vervollständigung} von $(X,d)$ und ist eindeutig bis auf
isometrische Isomorphie.
\minisec{Beweis}
\begin{description}
	\item[Eindeutigkeit:] 
	Angenommen, $(\hat{X}, \hat{d})$ ist ein weiterer vollständiger metrischer Raum und $\kappa : X \to \hat{X}$ eine Isometrie mit 
	$\overline{\kappa(X)}=\hat{X}$. Definiere $\gamma : \hat{X} \to \tilde{X}$ wie folgt: Falls $y \in \hat{X}$, wähle eine Folge $(x_n)_{n \in \mathds{N}}$ in $X$, sodass
	$y= \lim_{n \to \infty} \kappa(x_n)$. Setze nun
	\[
		\gamma(y) := \lim_{ n \to \infty} \iota(x_n) \in \tilde{X} 
	\]
	Zu zeigen: $\gamma$ ist ein isometrischer Isomorphismus.\todo{Achtung Wortspiel: Beweis \emph{vervollständigen} ;-D}
	\begin{description}
		\item[Injektivität:] Seien $y,y' \in \hat{X}$ mit $\gamma(y)=\gamma(y')$. Dann existieren Folgen $(x_n)_n, (x_n')_n$ in $X$ mit
		\[
			\lim_{ n \to \infty} \iota(x_n) = \lim_{ n \to \infty} \iota(x_n') 
		\]
	\end{description}
	\item[Existenz:] \begin{description}
		\item[Konstruktion von $(\tilde{X}, \tilde{d})$:] Setze $Y := \set[(x_n)_{n \in \mathds{N}}]{(x_n)_n \text{ ist Cauchy-Folge in }X}$. Definiere 
	\[
		(x_n)_n \sim (x_n') :\iff \lim_{ n \to \infty} d(x_n, x_n') = 0
	\]
	$\sim$ ist eine Äquivalenzrelation auf $Y$. Definiere nun $\tilde{X} := \nicefrac{Y}{\sim}$ und $\tilde{d} :  \tilde{X} \times \tilde{X} \to [0, \infty)$ durch
	\[
		\tilde{d} \enbrace*{\benbrace[\big]{(x_n)_n} , \benbrace[\big]{(x_n')_n}  } := \lim_{ n \to \infty} d(x_n , x_n') 
	\]
	$\tilde{d}$ ist eine wohldefiniert, d.h. falls $\benbrace*{(x_n)_n} = \benbrace*{(y_n)_n} $ und $\benbrace*{(x_n')_n} = \benbrace*{(y_n')_n}$, dann ist
	\[
		\lim_{n \to \infty} d(x_n , x_n') = \lim_{n \to \infty} d(y_n, y_n'). \tag*{(leichte Übung)}
	\]
	Weiter ist $\enbrace[\big]{d(x_n x_n')}_n$ eine Cauchy-Folge in dem vollständigen Raum $[0,\infty)$ und somit 
	konvergent: Sei dazu $\varepsilon>0$. Dann existiert ein $N \in \mathds{N}$, sodass $d(x_n,x_m) \le \frac{\varepsilon}{2}$ und $d(x_n', x_m') \le \frac{\varepsilon}{2}$ für alle $n,m \ge N$. Dann gilt nach der \hyperref[sub:vier_ungl]{Vierecksungleichung (siehe Anhang \ref*{sub:vier_ungl})}
	\[
		\abs[\big]{d(x_n,x_n') - d(x_m,x_m')} \le d(x_n, x_m) + d(x_n', x_m') \le \varepsilon
	\]
	\item[Einbettung von $X:$] Definiere nun $\iota : X \to \tilde{X}$ durch $x \mapsto [(x,x,x, \ldots )] \in \tilde{X}$. $\iota$ ist Isometrie, da
	\[
		\tilde{d} \enbrace[\big]{\iota(x),\iota(y)} = \lim_{ n \to \infty} d(x,y) = d(x,y)
	\]
	% Falls $\benbrace*{(x_n)_n} \in \tilde{X}$, finde $y_k \in \iota(X)$, so dass $\benbrace*{(x_n)} = \lim_{ k} y_k$ und $y_k = \benbrace*{(x_k, x_k, \ldots )}= \iota(x_k)$.
% 	Dann gilt $y_k \xrightarrow{k \to \infty} \benbrace*{(x_n)_n} $, d.h.
% 	\[
% 		\lim_{ n \to \infty} d(x_k,x_n) = \tilde{d} \enbrace*{y_k, \benbrace*{(x_n)_n}} \xrightarrow{k \to \infty} 0
% 	\]
% 	$\leadsto \tilde{X}= \overline{\iota(X)}$.
	Sei nun $\benbrace*{(x_n)_n} \in \tilde{X}$ und $\varepsilon>0$. Da $(x_n)_n$ eine Cauchy-Folge ist, gibt es ein $N \in \mathds{N}$, sodass für alle
	$n,m \ge N$ gilt $d(x_n,x_m) <\varepsilon$. Dann gilt
	\[
		\tilde{d}\enbrace[\big]{\iota(x_N), \benbrace*{(x_n)_n} } = \lim_{ n \to \infty} d \enbrace*{x_N, x_n} < \varepsilon
	\]
	\item[Vollständigkeit von $(\tilde{X}, \tilde{d})$:] Sei $\enbrace*{\overline{x}^m}_m$ eine Cauchyfolge in $\tilde{X}$.\marginnote{$\overline{x} \in \tilde{X}$}
	\[
		\Longrightarrow \forall \varepsilon >0 : \exists M(\varepsilon) \in \mathds{N} : \forall m,m' > M : \tilde{d} \enbrace*{\overline{x}^m, \overline{x}^{m'}} < \frac{\varepsilon}{3} 
	\]
	Wenn $\overline{x}^m = \benbrace*{(x^m_n)_n}$ und $\overline{x}^{m'} = \benbrace*{(x^{m'}_n)_n} $, dann gilt also für alle $m,m' >M(\varepsilon)$
	\begin{equation*}
		\lim_{ n \to \infty} d \enbrace*{x^m_n, x^{m'}_n} < \frac{\varepsilon}{3}  \label{eq:15:absch1}\tag{*}
	\end{equation*}
	Für alle $m$ ist $(x^m_n)_n$ eine Cauchyfolge, also gilt
	\begin{equation*}
		\forall m : \exists N(m) : \forall n,n' \ge N(m) : d \enbrace*{x^m_n, x^m_{n'}} < \frac{1}{m}  \label{eq:15:absch2}\tag{**}
	\end{equation*}
	Setze nun $z_n := x^n_{N(n)}$. Behauptung: $(z_n)_n$ ist eine Cauchyfolge. Sei $\varepsilon>0$. Dann gilt für $n,m>N=\max\set{M(\varepsilon), \frac{3}{\varepsilon}}$
	für ein $k > N(m), N(n)$ \marginnote{$d(z_n,z_m)$ ist unabhängig von $k$, also kann man $k$ beliebig groß wählen}
	\begin{align*}
		d(z_n,z_m) = d \enbrace*{x^n_{N(n)}, x^m_{N(m)}} \le 
		\underbrace{d\enbrace*{x^n_{N(n)}, x^n_k}}_{\stackrel{\mbox{\tiny\eqref{eq:15:absch2}}}{<} \frac{1}{n} < \frac{1}{N} \le \frac{\varepsilon}{3}} + 
		\underbrace{d\enbrace[\Big]{x^n_k, x^m_k}}_{\stackrel{\mbox{\tiny\eqref{eq:15:absch1}}}{\le} \frac{\varepsilon}{3}}   + 
		\underbrace{d\enbrace*{x^m_k, x^m_{N(m)}}}_{\stackrel{\mbox{\tiny\eqref{eq:15:absch2}}}{<} \frac{1}{m} < \frac{1}{N} \le \frac{\varepsilon}{3}} 
		< \varepsilon
	\end{align*}
	$\Rightarrow (z_n)_n$ ist eine Cauchyfolge, also 
	\begin{equation*}
		\forall \varepsilon>0 : \exists N_z(\varepsilon) : \forall n,m > N_z(\varepsilon)  : d(z_n,z_m) < \varepsilon  \label{eq:15:absch3} \tag{\#}
	\end{equation*}
	Es bleibt zu zeigen: $\lim_{m \to \infty} \overline{x}^m = \benbrace*{(z_n)_n}$. Sei dazu $\varepsilon>0$. Dann gilt für 
	$m \ge \max \set{\frac{2}{\varepsilon}, N_z(\varepsilon)}$
	\begin{align*}
		\tilde{d} \enbrace[\big]{\overline{x}^m, \benbrace*{(z_n)_n}} = \lim_{ n \to \infty} d \enbrace*{x^m_n, x^n_{N(n)}}\le \lim_{ n \to \infty} 
		\enbrace[\Bigg]{ \underbrace{d\enbrace*{x_n^m, x^m_{N(m)}}}_{\stackrel{\mbox{\tiny\eqref{eq:15:absch2}}}{<} \frac{1}{m} \le \frac{\varepsilon}{2}  } +
		\underbrace{d\enbrace*{x^m_{N(m)}, x^n_{N(n)}}}_{=d(z_m,z_n) \stackrel{\mbox{\tiny\eqref{eq:15:absch3}}}{<} \frac{\varepsilon}{2} } } < \varepsilon
	\end{align*}
	Also gilt $\overline{x}^m \xrightarrow{m \to \infty} \benbrace*{(z_n)_n}$ und $(\tilde{X},\tilde{d})$ ist vollständig. \bewende
	\end{description}
\end{description}
% subsection 15 (end)

\subsection[Definition: Raum der beschränkten, stetigen Abbildungen]{Definition} % (fold)
\label{sub:16}
Sei $(W,\mathcal{T})$ ein topologischer Raum und $(X,d)$ ein metrischer Raum. Sei 
\[
	C_b(W,X) = \set[f : W \to X]{f \text{ stetig und beschränkt}} 
\]
versehen mit der Metrik $d_{W,X}$, definiert durch
\[
	d_{W,X}(f,g) = \sup_{t \in W} d \enbrace[\big]{f(t), g(t)} 
\]
% subsection 16 (end)

\subsection[Bemerkung: $d_{W,X}$ als Metrik auf $C(W,X)$]{Bemerkung} % (fold)
\label{sub:17}
Auf $C(W,X) = \set{f : W \to X \text{ stetig}} $ ist $d_{W,X}$ eine \Index{erweiterte Metrik}, d.h. der Wert $\infty$ ist möglich.
$\tilde{d}_{W,X} := \min \set{1, d_{W,X}} $ ist eine \enquote{echte} Metrik auf $C(W,X)$.
% subsection 17 (end)

\subsection[Proposition: $X$ vollständig $\Rightarrow$ $C-B(W,X)$ und $C8W,X$ vollständig]{Proposition} % (fold)
\label{sub:18}
Falls $X$ vollständig ist, dann sind $C_b(W,X)$ und $C(W,X)$ vollständig (bezüglich $d_{W,X}$ bzw. $\tilde{d}_{W,X}$).
\minisec{Beweis}
Sei $(f_n)_n$ eine Cauchy-Folge in $C_b(W,X)$, also
\[
	\forall \varepsilon>0 : \exists N \in \mathds{N} : \forall n,m >N : \sup_{t \in W} d \enbrace[\big]{f_n(t), f_m(t)}< \varepsilon 
\]
$\Rightarrow$ für alle $t$ ist $\enbrace[\big]{f_n(t)}_n$ eine Cauchyfolge in $X$. Da $X$ vollständig ist, existiert ein $x_t \in X$ sodass 
$f_n(t) \xrightarrow{n \to \infty} x_t$. Definiere $f(t) := x_t$ punktweise.
\begin{description}
	\item[Vorbereitung:] Es gilt $\forall \varepsilon>0 : \exists N_t$ sodass $d(x_t, f_{N_t}(t)) \le \varepsilon$. Weiter ist
	\[
		\sup_{t \in W} d \enbrace[\big]{f_n(t), f_m(t)}\le \varepsilon \enspace \forall n,m \ge N_t
	\]
	$\Rightarrow \forall t : d \enbrace*{f_n(t), f_m(t)}\le \varepsilon \enspace \forall n,m \ge N$
\end{description}
\todo{Beweis zu Ende führen}
% subsection 18 (end)

\subsection[Proposition über eine Folge von Bällen]{Proposition} % (fold)
\label{sub:19}
Sei $(X,d)$ ein vollständiger metrischer Raum und $\enbrace*{\overline{B}(x_n, \varepsilon_n)}_{n \in \mathds{N}}$ mit 
$\overline{B}(x_{n+1}, \varepsilon_{n+1}) \subseteq \overline{B}(x_n,\varepsilon_n)$ und $\varepsilon_n \xrightarrow{n \to \infty} 0$. Dann existiert genau 
ein Punkt in $\bigcap_{n \in \mathds{N}} \overline{B}(x_n, \varepsilon_n)$.
\minisec{Beweis}
\begin{description}
	\item[Existenz:] Behauptung: Die Folge der Mittelpunkte $(x_n)_{n \in\mathds{N}}$ ist eine Cauchy-Folge: Sei $\varepsilon >0$. Finde $N \in \mathds{N}$,
	sodass $\varepsilon_n \le \varepsilon$ für alle $n \ge N$. Dann gilt für alle $n \ge m \ge N$
	\[
		d(x_n, x_m) \le \varepsilon_m \le \varepsilon,
	\]
	da $x_n \in \overline{B}(x_m, \varepsilon)$ ist. Da $X$ vollständig ist, existiert ein $x \in X$ mit $x_n \xrightarrow{n \to \infty} x$.
	
	Behauptung: $x \in \bigcap_{n \in \mathds{N}} \overline{B}(x_n, \varepsilon_n)$. Wähle dazu ein $N \in \mathds{N}$. Dann ist $(x_n)_{n \ge N}$ eine
	Cauchy-Folge in $\overline{B}(x_N,\varepsilon_N)$. $\overline{B}(x_N,\varepsilon_N)$ ist ein abgeschlossener Teilraum von $X$ und somit vollständig. 
	Also ist $x \in \overline{B}(x_N,\varepsilon_N)$. Da $N$ beliebig war, gilt $x \in \overline{B}(x_n,\varepsilon_n)$ für jedes $n \in \mathds{N}$.
	\item[Eindeutigkeit:] Es seien $x,y \in \bigcap_{n \in \mathds{N}} \overline{B}(x_n,\varepsilon_n)$, also 
	$x,y \in \overline{B}(x_n, \varepsilon_n) $ für alle $ n \in \mathds{N}$. Dann gilt
	\[
		d(x,y) \le d(x,x_n) + d(x_n,y) \le 2 \cdot \varepsilon_n \xrightarrow{n \to \infty} 0
	\]
	für jedes $n \in \mathds{N}$ und somit $d(x,y)=0 \Leftrightarrow x=y$. \bewende
\end{description}
% subsection 19 (end)

\subsection[Satz von Baire]{Satz von Baire\protect\footnote{nach René Louis Baire, \url{https://de.wikipedia.org/wiki/René_Louis_Baire}}} % (fold)
\label{sub:110} 
Es gelten folgende äquivalente Formulierungen:
\begin{enumerate}[a)]
	\item Sei $(X,d)$ ein vollständiger metrischer Raum und $A_0, A_1, \ldots $ eine Folge abgeschlossener Teilmengen. Falls $\bigcup_{n \in \mathds{N}} A_n$ eine
	offene Kugel enthält, so auch eines der $A_n$.
	\item In einem vollständigen metrischen Raum hat eine abzählbare Vereinigung von abgeschlossenen Mengen ohne innere Punkte keine inneren Punkte.
	\item In einem vollständigen metrischen Raum ist ein abzählbarer Durchschnitt von dichten offenen Mengen wieder dicht.
\end{enumerate}
\minisec{Beweis}
Zunächst die Äquivalenz: 
\begin{description}
	\item[a)$\Rightarrow$b):] Klar, da b) Kontraposition von a) ist.
	\item[b)$\Rightarrow$c):] Sei $U_n \subseteq X$ offen und dicht $\Rightarrow X \setminus U_n \subseteq X$ abgeschlossen und hat somit keine inneren Punkte. Aus b) folgt
	nun 
	\[
		\bigcup_{n \in \mathds{N}}X \setminus U_n = X \setminus \bigcap_{n \in \mathds{N}} U_n
	\]
	hat keine inneren Punkte. Also ist $\bigcap_{n \in \mathds{N}} U_n$ dicht in $X$.
	\item[c)$\Rightarrow$a):] ebenso.
\end{description}
Wir wollen nun a) durch Widerspruch beweisen, d.h. wir nehmen an, dass gilt: Jede offene Kugel schneidet $X \setminus A_k$ für jedes $k \in \mathds{N}$. Dazu wollen wir 
Folgen $(x_k)_{k \in \mathds{N}} \subset X$, $(\varepsilon_k)_{k \in \mathds{N}} \subset (0,1]$ finden mit 
\begin{enumerate}[(i)]
	\item $\varepsilon_k < \frac{1}{k+1}$,
	\item $\overline{B}(x_{k+1}, \varepsilon_{k+1}) \subset (X \setminus A_k) \cap B(x_k, \varepsilon_k) \subset \overline{B}(x_k, \varepsilon_k)$,
	\item $\overline{B}(x_k, \varepsilon_k) \subset \bigcup_{n \in \mathds{N}} A_n$
\end{enumerate}
für $k \in \mathds{N}$. Dann gilt 
\[
	\bigcup_{n \in \mathds{N}} A_n \stackrel{\text{(iii),(ii)}}{\supset} \bigcap \overline{B}(x_k, \varepsilon_k) \stackrel{\text{(ii)}}{\subset} 
	\bigcap_{k \in \mathds{N}}(X \setminus A_k) \cap B(x_k, \varepsilon_k) \subset \bigcap_{k \in \mathds{N}} X \setminus A_k = X \setminus 
	\enbrace*{\bigcup_{k \in \mathds{N}} A_k} 
\]
Aber wegen (i) und (ii) existiert nach Proposition \ref{sub:19} ein $x \in \bigcap_{k \in \mathds{N}}\overline{B}(x_k,\varepsilon_k)$. \light \medskip \\
Wir suchen also eine Abbildung $\overline{c} : \mathds{N} \to X \times (0,1], k  \mapsto(x_k, \varepsilon_k)$ mit (i),(ii),(iii) für $k \in \mathds{N}$. Setze
\begin{align*}
	P_m &:= \set[{c : \set{0, \ldots ,m} \to X \times (0,1] }]{\begin{array}{l}
		\text{(i),(iii) gilt für }k \in \set{0, \ldots , m}, \\
		\text{(ii) gilt für }k \in \set{0, \ldots , m-1}
	\end{array} }  \\
	P_\infty &:= \set[{c : \mathds{N} \to X \times (0,1]}]{\text{(i),(ii),(iii) gilt für } k \in \mathds{N} \rule{0cm}{0.3cm}} 
\end{align*}
Die Menge $P := \enbrace*{\bigcup_{m \in \mathds{N}} P_m} \cup P_\infty$ ist partiell geordnet bezüglich $\prec$:
\[
	c \prec c' \text{ falls } m \le m' \text{ und } c'\big|_{\set{0,\ldots ,m}} =c , \text{ bzw.} c'=c \text{ falls } m=m'=\infty
\]
$P$ ist nicht leer, denn nach Vorraussetzung existiert $(x_0,\varepsilon_0)$ mit $0 < \varepsilon_0 < 1$ und $\overline{B}(x_0, \varepsilon_0) \subset 
\bigcup_{n \in\mathds{N}} A_n$, d.h. $0 \mapsto (x_0,\varepsilon_0) \in P_0$. Jede total geordnete Teilmenge $\emptyset \not= \Gamma$ von $P$ besitzt eine obere Schranke:
\begin{itemize}
	\item Falls $\Gamma$ ein $c\in P_\infty$ enthält, so ist $c$ obere Schranke.\marginnote{warum?}
	\item Ebenso falls $\Gamma$ ein $c \in P_{\overline{m}}$ enthält und $\Gamma \cap P_{m'}= \emptyset$ für alle $\overline{m} < m' \in \mathds{N} \cup \set{\infty}$.
	\item Falls $\Gamma \subset \bigcup_{\mathds{N}} P_m$, aber $\Gamma \not\subset \bigcup_{m \le \overline{m}} P_m$ für jedes $\overline{m}$, so definieren wir eine
	obere Schranke in $P_\infty$ durch Einschränkung.\marginnote{wie genau?} 
\end{itemize}
Mit dem Lemma von Zorn folgt, dass $P$ ein maximales Element $\overline{c}$ besitzt. Behauptung: $\overline{c}\in P_\infty$ wie gewünscht:

Falls $\overline{c} \in P_m$ für ein $m \in \mathds{N}$, so gilt nach Annahme, dass $\enbrace*{X \setminus A_m} \cap B(x_m, \varepsilon_m) \not= \emptyset$. Dann 
existiert aber $(x_{m+1},\varepsilon_{m+1}) \in X \times (0,1]$ mit $\varepsilon_{m+1} < \frac{1}{m+2}$ und 
\[
	\overline{B}(x_{m+1}, \varepsilon_{m+1}) \subset  \enbrace*{X \setminus A_m} \cap B(x_m,\varepsilon_m) 
\]
Definiere $\overline{\overline{c}} : \set{0, \ldots , m+1} \to X \times (0,1]$ durch
\[
	k \longmapsto \begin{cases}
		(x_k,\varepsilon_k), &\text{ falls } k=m+1\\
		\overline{c}(k), &\text{ falls } k \in \set{0, \ldots , m} 
	\end{cases},
\]
dann gilt $\overline{c} \prec \overline{\overline{c}} \in P_{m+1}$. \light zur Maximalität von $\overline{c}$. \bewende
% subsection satz_baire (end)

\subsection[Bemerkungen zum Satz von Baire]{Bemerkungen} % (fold)
\label{sub:111}
\begin{enumerate}[(i)]
	\item Die Aussage gilt auch für lokalkompakte Hausdorffräume.\index{lokalkompakt} \hfill(Übung)
	\item Tatsächlich genügt eine schwächere Form des Auswahlaxioms(DC); das abzählbare Auswahlaxiom jedoch nicht.
	\item Falls $X$ \Index{separabel} ist (d.h. falls eine abzählbare dichte Teilmenge von $X$ existiert), dann lässt sich der Satz auch ohne (AC) beweisen.\hfill (Übung)
\end{enumerate}
% subsection 111 (end)

\subsection{Korollar: Satz der gleichmäßigen Beschränktheit} % (fold)
\label{sub:112}
Sei $(X,d)$ ein vollständiger metrischer Raum. Sei $F \subset C(X,\mathds{R})$ eine Menge, die \Index{punktweise gleichmäßig beschränkt} ist, d.h. für jedes $x \in X$ 
existiert $K_x \in \mathds{R}$, sodass
\[
	\abs*{f(x)} \le K_x \enspace \forall f \in F. 
\] 
Dann existieren $\emptyset \not= U \subset X$ offen und $K \in \mathds{R}$ so, dass 
\[
	\abs*{f(x)} \le K \enspace \forall x \in U, f \in F 
\]
\minisec{Beweis}
Definiere 
\begin{align*}
	X \supset A_n := \set[x \in X]{\abs*{f(x)} \le n \forall f \in F \rule{0cm}{0.3cm}} = \bigcap_{f \in F} f ^{-1} \enbrace[\big]{[-n,n]} 
\end{align*}
Also sind die $A_n \subset X$ abgeschlossen. Es ist $\bigcup_{n \in \mathds{N}} A_n = X$, wegen der Voraussetzung von punktweiser gleichmäßiger Beschränktheit. Außerdem enthält dies eine nichtleere offene Menge. Also existiert ein $n$ und $U \subset X$ nichtleer und offen mit $U \subset A_n$. \bewende
% subsection 112 (end)
% section 1 (end)
\newpage

\section{Normierte Räume, Hahn-Banach Sätze} % (fold)
\label{sec:2}

\subsection[Definition: Topologischer Vektorraum]{Definition} % (fold)
\label{sub:21}
Sei $X$ ein $\mathds{K}$-Vektorraum mit einer Topologie $\mathcal{T}$. Wir sagen, $X$ ist ein \Index{topologischer Vektorraum}, falls Addition und Skalarmultiplikation stetig sind:
\begin{align*}
	+ : X \times X  \to X \quad &, \quad(x,y) \mapsto x+y \\
	\cdot : \mathds{K} \times X \to X \quad &, \quad (\lambda ,x) \mapsto \lambda \cdot x
\end{align*}
% subsection 21 (end)

\subsection[Proposition: Unterräume topologischer Vektorräume sind topologische Vektorräume]{Proposition} % (fold)
\label{sub:22}
Sei $X$ ein topologischer $\mathds{K}$-Vektorraum und $Y \subset X$ ein Untervektorraum. Dann sind $Y$ und $\overline{Y}$ topologische Vektorräume mit der Unterraumtopologie.
\minisec{Beweis}
\begin{itemize}
	\item Klar für $Y$.
	\item Zu zeigen: $\overline{Y}$ ist ein Untervektorraum. Seien $x_0, y_0 \in \overline{Y}$ und sei $U$ eine offene Menge in $X$ mit $x_0 + y_0 \in U$. Dann existieren
	offene Mengen $V,W$ von $X$ mit $x_0 \in V$, $y_0 \in W$ und $V+W \subset U$, da die Addition stetig ist. Da $x_0 \in \overline{Y}$ und $x_0 \in V$, existiert 
	$x_1 \in Y$ mit $x_1 \in V$. Analog existiert $y_1 \in Y$ mit $y_1 \in W$. Daher ist $x_1 + y_1 \in Y \cap U$. Da $U$ eine beliebige offene Umgebung um $x_0+ y_0$ ist,
	folgt $x_0 +y_0 \in \overline{Y}$.
	
	Skalarmultiplikation genauso. \bewende
\end{itemize}
% subsection 22 (end)

\subsection[Proposition: Normierte Vektorräume sind topologische Vektorräume]{Proposition} % (fold)
\label{sub:23}
Ein normierter Vektorraum $(X, \norm{\cdot})$ ist ein topologischer Vektorraum bezüglich der von $\norm{\cdot}$ induzierten Topologie.
\minisec{Beweis}
Es gilt 
\begin{align*}
	\norm[X]{(x+y) - (x'+y')} \le \norm[X]{x-x'} + \norm[X]{y-y'}   
\end{align*}
Daher ist $+ : X \times X \to X$ gleichmäßig stetig bezüglich der Norm $\norm[X \times X]{(z,z')} := \norm[X]{z} + \norm[X]{z'}$ auf $X \times X$. Ähnlich folgt
\begin{align*}
	\norm[X]{\lambda x - \lambda' x'} = \norm[X]{\lambda (x-x') + (\lambda -\lambda')x'} \le \abs*{\lambda } \cdot \norm[X]{x-x'} + \abs*{\lambda -\lambda'} \cdot 
	\norm[X]{x'}      
\end{align*}
Daher ist $\cdot : \mathds{K} \times X \to X$ stetig, denn falls $(\lambda_i)_{i \in I}$ ein Netz in $\mathds{K}$ und $(x_i)_{i \in I}$ ein Netz in $X$ ist mit 
$\lambda_i \to \lambda $, $x_i\to x$, dann  gilt $\lambda_i x_i \to \lambda x$, also $\norm[X]{\lambda_i x_i - \lambda x} \to 0$. \bewende
% subsection 23 (end)

\subsection[Proposition: Stetigkeit einer linearen Abbildung zwischen topologischen Vektorräumen]{Proposition} % (fold)
\label{sub:24}
Seien $X,Y$ topologische Vektorräume, $T : X \to Y$ linear. Dann sind äquivalent:\index{stetig}
\begin{enumerate}[(1)]
	\item $T$ ist stetig.
	\item $T$ ist stetig in $0$.
	\item $T$ ist stetig in einem Punkt $\overline{x}$.
\end{enumerate}
\minisec{Beweis}
(1) $\Rightarrow$ (2) $\Rightarrow $ (3) ist klar. Wir zeigen (3) $\Rightarrow$ (1):

Für $y \in X$ definiere $L_y : X \to X$, $x \mapsto y+ x$. Dann ist $L_y$ bijektiv und stetig. Da auch $L_{-y}$ stetig ist, ist $L_y$ ein Homöomorphismus. Sei nun $T$ in 
$\overline{x}$ stetig. Sei $x_0 \in X$. Wir zeigen: $T$ ist stetig in $x_0$. Setze $y := \overline{x} - x_0$. Dann ist
\begin{align*}
	T(x_0) = T(x_0+y-y) = T(-y) + T(x_0+y) = L_{T(-y)} \circ T \circ L_y (x_0)
\end{align*}
Da $L_y(x_0) = \overline{x}$, $T$ stetig in $\overline{x}$ ist und $L_y$, $L_{T(-y)}$ stetig sind, ist somit $T$ stetig in $x_0$. \bewende
% subsection 24 (end)

\subsection[Definition: Stetige Funktionale und Operatoren]{Definition} % (fold)
\label{sub:25}
Für topologische Vektorräume $X,Y$ definieren wir 
\[
	\mathcal{L}(X,Y) = \set[T : X \to Y]{T \text{ linear und stetig}}  
\]
Wir schreiben $X^*:= \mathcal{L}(X,\mathds{K})$ für den \Index{Dualraum} von $X$. Die Elemente von $X^*$ heißen (stetige) \Index{Funktionale}.\footnote{Aus der linearen Algebra: $X^* = \Hom(X,\mathds{R}) = \set[T : X \to Y]{T \text{ linear}}$. Die Elemente von $X^*$ heißen Funktionale.}
Die Elemente von $\mathcal{L}(X,X)$ heißen (stetige) \Index{Operatoren} auf $X$. 
% subsection 25 (end)

\subsection[Bemerkung: Vektorraumstruktur auf $\mathcal{L}(X,Y)$, Algebrastruktur auf $\mathcal{L}(X,X)$]{Bemerkung} % (fold)
\label{sub:26}
$\mathcal{L}(X,Y)$ und $\mathcal{L}(X,\mathds{K}) = X^*$ sind Vektorräume, wobei die Vektorraumstruktur punktweise definiert ist. $\mathcal{L}(X,X)$ ist sogar eine Algebra mittels Komposition.\index{Algebra}
% subsection 26 (end)

\subsection[Proposition: Stetigkeit linearer Abbildungen mittels Normabschätzung]{Proposition} % (fold)
\label{sub:27}
Seien $X,Y$ normierte Vektorräume, $T : X \to Y$ linear. Dann ist $T$ stetig genau dann, wenn ein $\mu \ge 0$ existiert, sodass 
\[
	\norm[Y]{T(x)} \le \mu \cdot \norm[X]{x} \quad \forall x \in X  
\]
Ein stetiger Operator zwischen normierten Vektorräumen heißt deswegen auch \bet{beschränkt}\index{Operator!beschränkter Operator}.
\minisec{Beweis}
\begin{description}
	\item[\enquote{$\Leftarrow$}:] Klar: Wenn $x_n \to 0$, dann $T(x_n) \to 0$, denn $\norm{T(x_n)} \le \mu \norm{x_n } \to 0$. Also ist $T$ stetig in $0$ und nach
	\ref{sub:24} überall.
	\item[\enquote{$\Rightarrow $}:] Angenommen $T$ ist stetig. Dann setzen wir
	\[
		\mu := \sup \set[{\frac{1}{\norm[X]{x}} \cdot \norm[Y]{T(x)}  }]{x \in X \setminus \set{0} } 
	\] 
	Falls $\mu = \infty$, dann existieren $x_n \in X$ mit $\frac{1}{\norm[X]{x_n}} \cdot \norm[Y]{T(x_n)} \ge n$. Betrachte 
	$x_n' :=  \frac{1}{n \cdot \norm[X]{x_n}} \cdot x_n $, dann $\norm{x_n'} = \frac{1}{n}$, also $x_n' \to 0$. Aber es gilt $T(x_n')\not\longrightarrow 0$, denn
	\begin{align*}
		\norm[Y]{T(x_n')} = \frac{1}{n \cdot \norm[X]{x_n} } \cdot \norm[Y]{T(x_n)} \ge 1 
	\end{align*}
	für alle $n$, \light zur Stetigkeit von $T$. Also $\mu < \infty$. \bewende 
\end{description}
% subsection 27 (end)

\subsection[Definition: Operatornorm]{Definition} % (fold)
\label{sub:28}
Seien $X,Y$ normierte Vektorräume, $T \in \mathcal{L}(X,Y)$. Die \bet{Norm} (oder \Index{Operatornorm}) von $T$ ist
\[
	\norm{T} = \sup \set[{\frac{1}{\norm[X]{x}} \cdot \norm[Y]{T(x)}}]{x \in X \setminus \set{0} } = \sup \set[{\norm[Y]{T(x)} }]{x \in X \text{ mit } \norm[X]{x} 
	= 1\rule{0cm}{0.4cm}}   
\] 
Nach \ref{sub:27} ist $\norm{T}<  \infty $ und $\norm[Y]{T x} \le \norm{T} \cdot \norm[X]{x}$ für alle $x \in X$.
% subsection 28 (end)

\subsection[Propostion: $\mathcal{L}(X,Y)$ ist ein normierter Raum mit der Operatornorm]{Proposition} % (fold)
\label{sub:29}
Seien $X,Y$ normierte Räume, betrachte $\mathcal{L}(X,Y)$ mit $\norm{\cdot}$. Dann ist $\mathcal{L}(X,Y)$ ein normierter Raum.
\minisec{Beweis}
\begin{itemize}
	\item $\norm{T} \ge 0 $ für alle $T \in \mathcal{L}(X,Y)$.
	\item Es gilt
	\begin{align*}
		\norm{T}=0 \iff \forall x \in X, \norm[X]{x}=1 : \norm[Y]{Tx}=0 &\iff \forall x \in X, \norm[X]{x}=1 : Tx=0 \\ &\iff \forall x \in X : Tx=0 \\ &\iff T=0
	\end{align*}
	\item Sei $\lambda  \in \mathds{K}$, $T \in \mathcal{L}(X,Y)$. Dann gilt
	\[
		\norm{\lambda \cdot T} = \sup_{\norm[X]{x}=1} \norm[Y]{\lambda  \cdot T(x)}  = \sup_{\norm[X]{x}=1 } \abs*{\lambda } \cdot \norm[Y]{Tx} =
		\abs*{\lambda } \cdot \sup_{\norm[X]{x}=1 } \norm[Y]{Tx} = \abs*{\lambda } \cdot \norm{T}      
	\]
	\item Seien $T_1, T_2 \in \mathcal{L}(X,Y)$ und $x \in X$ mit $\norm[X]{x}=1 $. Dann gilt
	\[
		\norm[Y]{(T_1 + T_2)(x)} = \norm[Y]{T_1 x + T_2 x} \le \norm[Y]{T_1 x} + \norm[Y]{T_2 x}
	\]
	Damit folgt weiter
	\begin{align*}
		\norm{T_1 +T_2} = \sup_{\norm[X]{x}=1 } \norm[Y]{(T_1+ T_2) (x)} &\le \sup_{\norm[X]{x}=1} \enbrace*{\norm[Y]{T_1 x} + \norm[Y]{T_2 x}  }   \\
		&\le \sup_{\norm[X]{x}=1} \sup_{\norm[X]{x'}=1} \enbrace*{\norm[Y]{T_1 x} + \norm[Y]{T_2 x'}  }  \\ &= \norm{T_1} + \norm{T_2} \bewende   
	\end{align*}
\end{itemize}
% subsection 29 (end)

\subsection[Definition: Normierte $\mathds{K}$-Algebra]{Definition} % (fold)
\label{sub:210}
Eine $\mathds{K}$-Algebra heißt \bet{normiert}\index{normierte $\mathds{K}$-Algebra}, falls $A$ mit einer Norm $\norm{\cdot}$ versehen ist, so dass $(A,\norm{\cdot} )$ ein normierter Vektorraum und
\[
	\norm{a b} \le \norm{a} \cdot \norm{b} \quad \forall a,b \in A   
\]
Falls $A$ \Index{unital} ist, d.h. es existiert ein Einselement $1_A$, dann gilt $\norm{1_A} \le 1$.
\minisec{Bemerkung}
Wenn $\norm{1_A} < 1$, dann $\norm{1_A} = \norm{1_A \cdot 1_A} \le \norm{1_A} \cdot \norm{1_A} \Rightarrow \norm{1_A} =0 $. Also $1_A=0$ und damit $A= \set{0}$.\todo{Rev Chap 2}
% subsection 210 (end)

\subsection[Beispiele für normierte Algebren]{Beispiele} % (fold)
\label{sub:211}
\begin{enumerate}[(i)]
	\item $C(\Omega,\mathds{K})$, wobei $\Omega$ ein kompakter Hausdorffraum und die Multiplikation punktweise ist. Betrachte $\norm[\infty]{.}$. Für $f,g : \Omega \to K$ gilt dann
	$\norm[\infty]{f \cdot g} \le \norm[\infty]{f} \cdot \norm[\infty]{g}$
	\item $\mathcal{L}(X,X)$ für einen normierten Raum $X$ mit $\norm{.}$. Für $S,T \in \mathcal{L}(X,X)$ gilt
	\[
		\norm{S T} = \sup_{\norm[X]{x}=1} \norm[X]{S(T(x))} \le \sup_{\norm[X]{x}=1} \enbrace[\big]{\norm{S} \cdot \norm[X]{Tx}} = \norm{S} \cdot \norm{T}.     
	\]
\end{enumerate}
% subsection 211 (end)

\subsection[Proposition: Norm auf dem Produkt normierter Räume]{Proposition} % (fold)
\label{sub:212}
Seien $X,Y$ normierte Räume. Dann ist $X \times Y$ ein normierter Raum mit 
\begin{align*}
	 \norm[1]{(x,y)} &:= \norm[X]{x} + \norm[Y]{y}   \quad \text{ für } (x,y) \in X \times Y \\
	\text{oder } \quad \norm[\infty]{(x,y)} &:= \max \set{\norm[X]{x}, \norm[Y]{y}}  
\end{align*}
Beide Normen sind äquivalent.
% subsection 212 (end)

\subsection[Proposition und Definition: Norm auf dem Quotientenraum]{Proposition und Definition} % (fold)
\label{sub:213}
Sei $X$ ein normierter Raum und $Y \subseteq X$ ein abgeschlossener Unterraum. Dann ist $\nicefrac{X}{Y} = \set[x +Y]{x \in X}$ ein normierter Raum mit
\[
	\norm[\nicefrac{X}{Y}]{x +Y} := \inf \set[{\norm[X]{x+y} }]{y \in Y}  
\]
Die Quotientenabbildung $q : X \twoheadrightarrow \nicefrac{X}{Y}$, $x \mapsto x+ Y$ ist stetig, linear und $\norm{q} \le 1$ und offen, d.h. bildet offene Mengen in $X$ auf offene Mengen in $\nicefrac{X}{Y}$ ab.
\minisec{Bemerkung}
Daraus folgt, dass Normtopologie und Quotiententopologie auf $\nicefrac{X}{Y}$ gleich sind.
\minisec{Beweis}
\begin{enumerate}[a)]
	\item $\nicefrac{X}{Y}$ ist ein Vektorraum mit $(x+Y) + (x'+Y) = (x+x') + Y$.
	\item Sei $\overline{x} = x+ Y \in \nicefrac{X}{Y}$. Es gilt
	\begin{itemize}
		\item $\norm{\overline{x}} \ge 0 $ ist klar für alle $\overline{x} \in \nicefrac{X}{Y}$
		\item Angenommen $\norm[\nicefrac{X}{Y}]{\overline{x}} =0$, d.h. $\inf_{y \in Y} \norm[X]{x+y} =0$. Also existiert eine Folge $(y_n)_n \subset Y$ mit
		$\norm[X]{x+y_n} \xrightarrow{n \to \infty} 0 \Rightarrow y_n \xrightarrow{n \to \infty} -x$. Also ist $-y$ und damit auch $y$ in $Y$, da $Y$ abgeschlossen ist.
		Damit ist $\overline{x}=0 $ 
	\end{itemize}
	\item Sei $0 \not= \lambda \in \mathds{K}$, $x \in X$. Dann gilt
	\begin{align*}
		\norm[\nicefrac{X}{Y}]{\lambda \cdot \overline{x}} = \inf_{y \in Y} \norm[X]{\lambda x + y} = \inf_{y \in Y} \norm[X]{\lambda x + \lambda y} 
		= \inf_{y \in Y} \abs*{\lambda } \cdot \norm[X]{x+y}  =   \abs*{\lambda } \cdot \norm[\nicefrac{X}{Y}]{\overline{x}}  
	\end{align*}
	\item Seien $\overline{x}, \overline{y} \in \nicefrac{X}{Y}$. Dann gilt
	\begin{align*}
		\norm[\nicefrac{X}{Y}]{\overline{x} + \overline{y}} = \inf_{z \in Y} \norm[X]{x+y+z} = \inf_{z,z' \in Y} \norm[X]{x+y+z+z'} &\le \inf_{z,z' \in Y}
		\enbrace[\Big]{\norm[X]{x+z} + \norm[X]{y+z'}} \\ &= \norm[\nicefrac{X}{Y}]{\overline{x}} + \norm[\nicefrac{X}{Y}]{\overline{y}}
	\end{align*}
	Damit ist $\norm[\nicefrac{X}{Y}]{.}$ eine Norm auf $\nicefrac{X}{Y}$.
	\item Linearität von $q$ ist klar. Es gilt
	\[
		\norm[\nicefrac{X}{Y}]{\overline{x}} = \inf_{y \in Y} \norm[X]{x+y} \le \norm[X]{x}   
	\]
	Also
	\[
		\norm{q} = \sup_{\norm[X]{x}=1} \norm[\nicefrac{X}{Y}]{q(x)} \le \sup_{\norm[X]{x}=1 } \norm[X]{x} = 1   
	\]
	\item Zu zeigen: $q$ ist offen. Sei $x \in X$, $\varepsilon>0$. Betrachte die offene Kugel $B(x,\varepsilon) \subset X$. Wir zeigen 
	$B(\overline{x}, \varepsilon) \subset q \enbrace*{B(x,\varepsilon)}$. Sei also $\overline{z} = z + Y$ in $B(\overline{x},\varepsilon)$. Dann gilt
	$\norm[\nicefrac{X}{Y}]{\overline{x}- \overline{z}} < \varepsilon$, worauf folgt
	\[
		\inf_{y \in Y} \norm[X]{x-z+y} < \varepsilon 
	\]
	$\Rightarrow$ es existiert $y \in Y : \norm[X]{x-z+y}<\varepsilon $. Es gilt $\overline{z}=\overline{z-y}\in q\enbrace*{B(x,\varepsilon)}$, da $z-y\in B(x,\varepsilon)$.
	
	Sei $V \subset X$ offen, zeige $q(V) \subset \nicefrac{X}{Y}$ offen. Für $x \in V$ finde $\varepsilon>0$, sodass $B(x,\varepsilon) \subset V$. Dann folgt
	$B(\overline{x},\varepsilon) \subseteq q \enbrace[\big]{ B(x,\varepsilon)} \subseteq q(V)$. \bewende
\end{enumerate}
% subsection 213 (end)

\subsection[Definition: Banachraum und Banachalgebra]{Definition} % (fold)
\label{sub:214}
Ein \Index{Banachraum} ist ein vollständiger, normierter Raum. Eine \Index{Banachalgebra} ist eine vollständige normierte Algebra.  
% subsection 214 (end)

\subsection[Beispiel für eine Banachalgebra]{Beispiel} % (fold)
\label{sub:215}
$C(\Omega, \mathds{K} )$, wobei $\Omega$ ein kompakter Hausdorffraum ist, ist eine Banachalgebra.
% subsection 215 (end)

\subsection[Proposition: Produkte und Quotienten von Banachräumen]{Proposition} % (fold)
\label{sub:216}
Produkte und Quotienten (nach abgeschlossenen Unterräumen) von Banachräumen sind wieder Banachräume.
\minisec{Beweis}
\begin{description}
	\item[Produkt $X \times Y$:] Es gilt $\norm[\infty]{(x,y)} = \max \set{\norm{x}, \norm{y}}$. Sei $\enbrace*{(x_n, y_n)}_n \subset X \times Y$ eine Cauchyfolge.
	Dann sind $(x_n)_n$ und $(y_n)_n$ Cauchyfolgen und es gilt $x_n \to x$ und $y_n \to y$ und somit auch $(x_n, y_n) \to (x,y)$.
	\item[Quotient $\nicefrac{X}{Y}$:] Sei $X$ ein Banachraum und $Y \subseteq X$ ein abgeschlossener Unterraum. 
	
	Falls $\overline{x}, \overline{y} \in \nicefrac{X}{Y}$, so existiert $y' \in X$, sodass $\overline{y} = \overline{y'}$ und 
	\[
		\norm[X]{x-y'} \le 2 \cdot \norm[\nicefrac{X}{Y}]{\overline{x} - \overline{y}}  
	\] 
	Sei nun $\enbrace*{\overline{x}_n }_n$ eine Cauchyfolge in $\nicefrac{X}{Y}$. Es gibt eine Teilfolge $\enbrace*{\overline{x}_{n_k}}_k$, sodass
	\[
		\sum_{k=1}^{\infty} \norm[\nicefrac{X}{Y}]{\overline{x}_{n_{k-1}} - \overline{x}_{n_k}} < \infty 
	\]
	Falls $\enbrace*{\overline{x}_{n_k} }_k $ konvergiert, dann auch $\enbrace*{\overline{x}_n}_n$. Wir dürfen also \oE annehmen, dass $\sum_{n=1}^{\infty} \norm{\overline{x}_{n-1} - \overline{x}_n }< \infty$. Wähle $x_1' \in \overline{x}_1$. Wähle induktiv $x_n' \in X$  mit
	\[
		\norm[X]{x_{n-1}' - x_n'} \le 2 \cdot \norm[\nicefrac{X}{Y}]{\overline{x_{n-1} - \overline{x}_n } }  
	\]
	und $x_n' \in \overline{x}_n$. Dann gilt $\sum_{n=1}^{\infty} \norm[X]{x_{n-1}' - x_n'} < \infty $. Daher ist $(x_n')_n$ eine Cauchyfolge in $X$. Also existiert 
	$x \in X$ mit $x_n \to x$. Dann $\overline{x}_n  \to \overline{x}$. \bewende
\end{description}
% subsection 216 (end)

\subsection[Definition: Sublineare Abbildung]{Definition} % (fold)
\label{sub:217}
Sei $X$ ein $\mathds{R}$-Vektorraum. Eine Abbildung $\varphi : X \to \mathds{R}$ heißt \Index{sublinear}, falls gilt:
\begin{enumerate}[(i)]
	\item $\varphi(\lambda \cdot x) = \lambda \cdot \varphi(x)$, für $x \in X$, $\lambda \in [0,\infty)$
	\item $\varphi(x+y) \le \varphi(x) + \varphi(y)$, für $x,y \in X$.
\end{enumerate}
Setze $\mathcal{S}(X) := \set[\varphi : X \to \mathds{R}]{\varphi \text{ sublinear}}$. $\mathcal{S}(X)$ ist partiell geordnet mit 
\[
	\varphi \le \psi :\iff \varphi(x) \le \psi(x) ,\quad  x \in X
\]
% subsection 217 (end)

\subsection[Beispiele für sublineare Abbildungen]{Beispiele} % (fold)
\label{sub:218}
\begin{enumerate}[(i)]
	\item Halbnormen sind sublinear.
	\item \[
		\begin{tikzpicture}[scale=0.5]
			\draw (-3,0) -- (3,0) node[right]{$X$};
			\draw (0,-3) -- (0,3) node[right]{$\mathds{R}$};
			\draw[very thick, dashed,OrangeRed2] (0,0) -- (220:4);
			\draw[very thick,SeaGreen4] (40:4) -- (0,0) -- (205:4);
		\end{tikzpicture}
	\]
	$0 = \varphi(0) = \varphi(x-x) \le \varphi(x) + \varphi(-x)$.
\end{enumerate}
% subsection 218 (end)

\subsection[Proposition: Die sublinearen Abbildungen $\mathcal{S}(X)$ sind nach unten induktiv geordnet]{Proposition} % (fold)
\label{sub:219}
Sei $X$ ein $\mathds{R}$-Vektorraum. $\mathcal{S}(X)$ ist nach unten \Index{induktiv geordnet}, d.h. jede nichtleere total geordnete Teilmenge besitzt eine untere Schranke.
\minisec{Beweis}
Sei $\emptyset \not= (\varphi_i)_{I} \subseteq \mathcal{S}(X)$ total geordnet. Setze $\varphi(x) := \inf_{i \in I} \varphi_i(x)$, dann gilt 
\[
	- \varphi_i(-x) \le \varphi_i(x) \enspace\Longrightarrow \enspace- \varphi(x) \le \varphi(x) \le \varphi_i(x) < \infty
\]
ebenso $-\varphi(x) \le \varphi(-x) \le \varphi_i(-x)$ für $x \in X$, $i \in I$. Insbesondere gilt $- \infty <  \varphi(x) < \infty$, $x \in X$. 
Die Sublinearität von $\varphi$ und $\varphi \le \varphi_i$, $i \in I$ sind klar (Warum?). \bewende
% subsection 219 (end)

\subsection[Propostion: Minimale Elemente in $\mathcal{S}(X)$ sind genau die lineare Abbildungen]{Proposition} % (fold)
\label{sub:220}
Sei $X$ ein $\mathds{R}$-Vektorraum und $\varphi \in \mathcal{S}(X)$. Dann ist $\varphi$ in $\mathcal{S}(X)$ minimal genau dann, wenn $\varphi$ linear ist.
\minisec{Beweis}
\begin{description}
	\item["$\Leftarrow$":] Sei $\psi \le \varphi$ mit $\psi$ sublinear und $\varphi$ linear. Dann folgt $\psi(x) \le \varphi(x)$ und $\psi(-x) \le \varphi(-x)$ für 
	$x \in X$.
	\[
		 \Longrightarrow \enspace- \psi(x) \le \psi(-x) \le \varphi(-x) = - \varphi(x)
	\]
	Daraus folgt $\varphi(x) \le \psi(x)$ und somit muss $\varphi=\psi$ gelten. Also ist $\varphi$ minimal.
	\item["$\Rightarrow $":] Sei $\varphi \in \mathcal{S}(X)$ minimal. Zu $\overline{x} \in X$ definiere $\varphi_{\overline{x}} : X \to \mathds{R}$ durch
	\[
		\varphi_{\overline{x}}(x) := \inf_{\lambda \ge 0} \enbrace[\Big]{\varphi \enbrace*{x+ \lambda \cdot \overline{x} } - \lambda \cdot \varphi(\overline{x})}
	\]
	Es gilt $\varphi_{\overline{x}}(x) \in (-\infty,\infty)$, denn für $x \in X, \lambda \ge 0$ ist \marginnote{mit (ii) aus \ref{sub:217}}
	\[
		-\varphi(-x) \le \varphi(x+ \lambda \cdot \overline{x}) - \lambda  \cdot \varphi(\overline{x}) \le \varphi(x)
	\]
	$\varphi_{\overline{x}}$ ist sublinear:
	\begin{enumerate}[(i)]
		\item Sei $\mu >0$. Dann gilt für $x \in X$.
		\begin{align*}
			\varphi_{\overline{x}}(\mu \cdot x) = \inf_{\lambda \ge 0} \enbrace[\Big]{ \varphi \enbrace*{\mu \cdot x + \lambda  \cdot \overline{x} } 
			- \lambda \cdot \varphi(\overline{x} )} &= \inf_{\lambda \ge 0 } \mu \cdot  \enbrace*{ \varphi \enbrace*{x + \frac{\lambda }{\mu} \cdot \overline{x}}
			- \frac{\lambda }{\mu} \cdot \varphi(\overline{x} )  } \\
			&= \mu \cdot \inf_{\lambda' \ge 0} \enbrace[\Big]{ \varphi \enbrace*{x+ \lambda' \cdot \overline{x} } - \lambda' \varphi(\overline{x})} \\
			&= \mu \cdot \varphi_{\overline{x}}(x) 
		\end{align*}
		$\varphi_{\overline{x}}(0 \cdot x) = 0$ ist klar.
		\item Zu $x,y \in X$, $\varepsilon>0$ wähle $\lambda_x, \lambda_y \ge 0$ mit 
		\begin{align*}
			\varphi_{\overline{x}}(x) &\ge \varphi \enbrace*{x + \lambda_x \cdot \overline{x}} - \lambda_x \cdot \varphi(\overline{x}) - \varepsilon \\
			\varphi_{\overline{x}}(y) &\ge \varphi \enbrace*{y + \lambda_y \cdot \overline{x}} - \lambda_y \cdot \varphi(\overline{x}) - \varepsilon 
		\end{align*}
		Setze $\lambda  := \lambda_x + \lambda_y$, dann gilt
		\begin{align*}
			\varphi_{\overline{x}}(x) + \varphi_{\overline{x}}(y) &\ge \varphi \enbrace*{ x+ \lambda_x \cdot \overline{x}} + 
			\varphi \enbrace*{y+ \lambda_y \cdot \overline{x} } - \lambda  \cdot \varphi(\overline{x}) - 2 \varepsilon \\
			&\ge  \varphi \enbrace*{x +y + \lambda \cdot \overline{x}} - \lambda  \cdot \varphi(\overline{x}) - 2 \varepsilon \\
			&\ge \varphi_{\overline{x}} (x+y) - 2 \varepsilon
		\end{align*}
		Da $\varepsilon>0$ beliebig war, gilt $\varphi_{\overline{x}}(x) + \varphi_{\overline{x}}(y) \ge \varphi_{\overline{x}}(x+y)$. Also ist $\varphi_{\overline{x}}$
		sublinear.
	\end{enumerate}
	$\varphi_{\overline{x}} \le \varphi$ ist klar. Da $\varphi$ minimal ist, folgt $\varphi_{\overline{x}} = \varphi$. Wir erhalten 
	\[
		\varphi(x) + 1 \cdot \varphi(\overline{x}) \le \varphi(x+ 1 \cdot \overline{x}) \le \varphi(x) + 1 \cdot \varphi(\overline{x})
	\]
	für jedes $x, \overline{x}\in X$. Also ist $\varphi$ additiv und somit linear. \bewende
\end{description}
% subsection 220 (end)

\subsection{Satz von Hahn-Banach} % (fold)
\label{sub:221}
Sei $X$ ein $\mathds{R}$-Vektorraum und $\varphi \in \mathcal{S}(X)$.\marginnote{$X$ ist nicht zwangsweise endlich-dimensional!}
Dann existiert $\psi : X \to \mathds{R}$ linear mit $\psi \le \varphi$.
\minisec{Beweis}
Sei $\mathcal{S}_\varphi : =\set[\varphi' \in \mathcal{S}(X)]{\varphi' \le \varphi} \ni \varphi$. Nach Proposition \ref{sub:219} ist $\mathcal{S}(X)$ und damit auch 
$\mathcal{S}_\varphi$ nach unten induktiv geordnet. Nach dem Lemma von Zorn enthält $\mathcal{S}_\varphi$ ein minimales Element $\psi$. $\psi$ ist auch minimal in 
$\mathcal{S}(X)$ \hfill (Warum?)\\
Also ist $\psi$ linear nach Proposition \ref{sub:220}. \bewende \medskip \\
Wichtige Folgerungen: Fortsetzungs- und Trennungssätze.
% subsection 221 (end)

\subsection[Satz (Hahn-Banach): Existenz einer linearen Fortsetzung]{Satz} % (fold)
\label{sub:222}
Sei $X$ ein $\mathds{R}$-Vektorraum und $\varphi : X \to \mathds{R}$ sublinear. Sei $Y \subset X$ ein linearer Unterraum und $\psi : Y \to \mathds{R}$ linear mit
$\psi \le \varphi\big|_{Y}$. Dann existiert $\overline{\psi} : X \to \mathds{R}$ linear mit $\overline{\psi}\big|_Y = \psi$ und $\overline{\psi} \le \varphi$.
\minisec{Beweis}
Definiere $\tilde{\varphi} : X \to \mathds{R}$ durch $\tilde{\varphi}(x) := \inf_{y \in Y} \enbrace[\big]{ \varphi(x-y) + \psi(y)}$, $x \in X$. Es gilt
\[
	\varphi(x-y) + \psi(y) \ge \varphi(-y) - \varphi(-x) - \psi(-y) \ge - \varphi(-x) > -\infty \qquad \forall x \in X, y \in Y
\]
Also ist $\tilde{\varphi}$ wohldefiniert. $\tilde{\varphi}$ ist sublinear: vergleiche \ref{sub:220}. Nach Satz von Hahn-Banach (\ref{sub:221}) existiert ein 
$\overline{\psi} : X \to \mathds{R}$ linear mit $\overline{\psi} \le \tilde{\varphi}$. Aus \ref{sub:220} folgt, dass $\overline{\psi}$ minimal in $\mathcal{S}(X)$ ist.
Ebenso ist $\overline{\psi}\big|_{Y}$ minimal in $\mathcal{S}(Y)$. Also gilt
\[
	\overline{\psi} \le \tilde{\varphi}\big|_{Y} \le \psi \quad \Rightarrow \quad \psi = \overline{\psi}\big|_{Y} \bewende  
\]
% subsection 222 (end)

\subsection{Satz} % (fold)
\label{sub:223}
Sei nun $X$ ein $\mathds{K}$-Vektorraum und $p : X \to \mathds{R}$ eine Halbnorm. Sei $Y \subset X$ ein Untervektorraum und $\psi : Y \to \mathds{K}$ linear mit
$\abs*{\psi(y)} \le p(y)$ für $y \in Y$. Dann existiert $\overline{\psi} : X \to \mathds{K}$ linear mit $\overline{\psi}\big|_{Y} = \psi$ und 
\[
	\abs*{\overline{\psi}(x) } \le p(x) \quad \text{ für } x \in X 
\]
\minisec{Beweis}
Sei zunächst $\mathds{K}=\mathds{R}$. $p$ ist sublinear und es gilt $\psi \le p\big|_{Y}$. Nach \ref{sub:222} existiert eine lineare Fortsetzung 
$\overline{\psi} : X \to \mathds{R}$ mit $\overline{\psi} \le p$. Es gilt auch
\begin{align*}
	- \overline{\psi}(x) =  \overline{\psi}(-x) \le p(-x)  = p(x) \quad \Longrightarrow \quad \abs*{\overline{\psi}(x) } \le p(x) \quad \text{ für } x \in X
\end{align*}
Sei nun $\mathds{K}=\mathds{C}$. Definiere $\psi_1 := \re(\psi) : Y \to \mathds{R}$. $\psi_1$ ist $\mathds{R}$-linear (warum?). Es gilt 
$\abs*{\psi_1(y)} \le \abs*{\psi(y)} \le p(y)$, $y \in Y$. Es existiert also ein $\overline{\psi}_1 : X \to \mathds{R}$ $\mathds{R}$-linear mit
\[
	\overline{\psi}_1 (y) = \psi_1(y), \enspace y \in Y \quad \text{ und } \quad \abs*{\overline{\psi}_1(x)} \le p(x) , \enspace x \in X 
\]
Definiere jetzt $\overline{\psi} : X \to \mathds{C}$ durch $\overline{\psi}(x) := \overline{\psi}_1(x) - i \cdot \overline{\psi}_1(i \cdot x)$, $x \in X$.
Dann ist $\overline{\psi}$ $\mathds{C}$-linear (warum?). Weiter gilt $\overline{\psi}(y) = \psi(y)$ für $y \in Y$.
($\re(\overline{\psi}\big|_{Y})= \re(\psi)$ und $\overline{\psi}\big|_{Y}$ und $\psi$ sind beide $\mathds{C}$-linear).

Zu $x \in X$ wähle $\lambda \in \mathds{C}$ mit $\abs*{\lambda }=1$ und 
\[
	\abs*{\overline{\psi}(x)} = \lambda  \cdot \overline{\psi}(x) = \overline{\psi}(\lambda \cdot x) = \overline{\psi}_1(\lambda \cdot x) \le p(\lambda \cdot x)
	= \abs*{\lambda } \cdot p(x) = p(x) \bewende     
\]
% subsection 223 (end)

\subsection{Satz} % (fold)
\label{sub:224}
Sei $X$ ein normierter $\mathds{K}$-Vektorraum, $Y \subset X$ ein Unterraum und $\psi : Y \to \mathds{K}$ linear und stetig. Dann existiert eine lineare Fortsetzung 
$\overline{\psi} : X \to \mathds{K}$ mit $\norm[\mathcal{L}(X,\mathds{K})]{\overline{\psi}} = \norm[\mathcal{L}(Y,\mathds{K})]{\psi}$.
\minisec{Beweis}
Definiere eine Halbnorm durch $p(x) := \norm{x} \cdot \norm{\psi}$. Nach Satz \ref{sub:223} besitzt $\psi$ eine lineare Fortsetzung $\overline{\psi}$ mit
$\abs*{\overline{\psi(x)}} \le p(x)= \norm{\psi} \cdot \norm{x}$, $x \in X$. Es folgt $\norm{\overline{\psi}} \le \norm{\psi}$. $\norm{\psi} \le \norm{\overline{\psi}}$
ist trivial, da $\psi\big|_{Y} = \overline{\psi}$. \bewende
% subsection 224 (end)

\subsection[Definition: Konvexe Teilmenge eines $\mathds{K}$-Vektorraums]{Definition} % (fold)
\label{sub:225}
Sei $X$ ein $\mathds{K}$-Vektorraum. Sei $M \subset X$ eine Teilmenge. $M$ heißt \Index{konvex}, falls für $a,b \in M$, $\lambda \in [0,1]$ gilt
\[
	(1-\lambda ) \cdot a + \lambda \cdot b \in M
\]
Ist $X$ ein topologischer Vektorraum und $M \subset X$ konvex, so ist auch $\overline{M}$ konvex. \hfill (warum?)\\
Ist $N \subset X$ eine beliebige Teilmenge, so ist 
\[
	\conv(N) := \smashoperator{\bigcap_{\substack{N \subset M \subset X \\ M \text{ konvex} }}} \, M
\]
$\conv(N)$ ist konvex. (warum?)
% subsection 225 (end)

\subsection{Satz} % (fold)
\label{sub:226}
Sei $X$ ein $\mathds{R}$-Vektorraum, $\emptyset \not= M \subset X$ konvex, $\varphi \in \mathcal{S}(X)$. Dann existiert $\psi : X \to \mathds{R}$ linear mit 
$\psi \le \varphi$ und 
\begin{equation*}
	\inf_{y \in M} \varphi(y) = \inf_{y \in M} \psi(y) \label{eq:226:1}\tag{*}
\end{equation*}
\minisec{Beweis}
Setze $\mu := \inf_{y \in M} \varphi(y)$. Falls $\mu = -\infty$, so folgt die Behauptung aus dem Satz von Hahn-Banach \ref{sub:221}. $\eqref{eq:226:1}$ ist trivialerweise
erfüllt. Sei also $\mu \in \mathds{R}$. Definiere $\tilde{\varphi} : X \to \mathds{R}$ durch
\[
	\tilde{\varphi}(x) := \inf_{y \in M, \lambda  \ge 0} \enbrace[\big]{ \varphi(x+ \lambda \cdot y) - \lambda  \cdot \mu} , \quad x \in X
\]
Es gilt $\varphi(x + \lambda \cdot y)- \lambda \cdot \mu \ge - \varphi(-x)$, $x \in X$, $y\in M$, $\lambda \ge 0$. Also ist $\tilde{\varphi}(x) \ge - \varphi(-x) >-\infty$.
$\tilde{\varphi}$ ist sublinear:
\begin{enumerate}[(i)]
	\item $\tilde{\varphi}(\gamma \cdot x) = \gamma \cdot \tilde{\varphi}(x)$, für $x \in X$, $\gamma \in \ge 0$, wie in \ref{sub:222}.
	\item Seien $x,z \in X$ und $\varepsilon>0$. Wähle $y_x, y_z \in M$, $\lambda_x, \lambda_y \ge 0$ mit
	\[
		\tilde{\varphi}(x) \ge \varphi (x +\lambda_x \cdot y_x) - \lambda_x \cdot \mu -\varepsilon \quad \text{ und }\quad 
		\tilde{\varphi}(x) \ge \varphi (z +\lambda_z \cdot y_z) - \lambda_z \cdot \mu -\varepsilon
	\] 
	Es folgt 
	\begin{align*}
		\tilde{\varphi}(x) + \tilde{\varphi}(z) &\ge \varphi \enbrace[\big]{x+z+ \lambda_x \cdot y_x + \lambda_z \cdot y_z} - (\lambda_x + \lambda_z) \cdot \mu 
		- 2 \varepsilon \\
		&= \varphi \enbrace[\Bigg]{x+z + (\lambda_x+ \lambda_z)\cdot \underbrace{\enbrace*{\frac{\lambda_x}{\lambda_x + \lambda_z} \cdot y_x + \frac{\lambda_z}{\lambda_x + \lambda_z} \cdot y_z }}_{\in M} } - (\lambda_x + \lambda_z) \cdot \mu -2 \varepsilon \\
		&\ge \tilde{\varphi}(x+z) - 2 \varepsilon
	\end{align*}
\end{enumerate}
Nach Hahn-Banach (\ref{sub:221}) existiert $\psi : X \to \mathds{R}$ linear mit $\psi \le \tilde{\varphi} \le \varphi$. Für $y \in M$ gilt 
\[
	-\psi(y) = \psi(-y) \le \tilde{\varphi}(-y) \le \varphi(-y + 1 \cdot y)- 1 \cdot \mu= -\mu
\]
also $\mu \le \psi(y) \le \varphi(y)$, woraus \eqref{eq:226:1} folgt. \bewende
% subsection 226 (end)

\subsection{Satz} % (fold)
\label{sub:227}
Sei $X$ ein normierter $\mathds{R}$-Vektorraum und $A,B \subset X$ nichtleere konvexe Teilmengen mit \index{Distanz $\dist$}
\[
	\dist(A,B) := \inf \set[\norm{a-b}]{a \in A, b \in B} > 0
\]
Dann existiert $\psi : X \to \mathds{R}$ stetig und linear mit $\psi(A) \cap \psi(B) = \emptyset$.
\minisec{Beweis}
$A - B := \set[a-b]{a \in A, b \in B} \subset X$ ist konvex. (Warum?) Nach Satz \ref{sub:226} existiert $\psi : X \to \mathds{R}$ linear mit $\psi \le \norm{.}$ und 
\[
	0 < \dist(A,B)=  \inf_{y \in A -B} \norm{y} = \inf_{y \in A-B} \psi(y) = \inf_{a \in A} \psi(a) - \sup_{b \in B} \psi(b). \bewende
\]
% subsection 227 (end)
% section 2 (end)
\newpage

\section[Operatoren zwischen Banachräumen, Satz von der offenen Abbildung]{Operatoren zwischen Banachräumen. Die Sätze von der offenen Abbildung und vom abgeschlossenen Graphen} % (fold)
\label{sec:3}

\subsection[Proposition: $Y$ vollständig $\Rightarrow $ $\mathcal{L}(X,Y)$ ist vollständig]{Proposition} % (fold)
\label{sub:31}
Seien $X,Y$ normierte Vektorräume, $Y$ vollständig. Dann ist $\mathcal{L}(X,Y)$ vollständig bezüglich $\norm[\mathcal{L}(X,Y)]{.}$. \todo{RevChap3}
\minisec{Beweis}
Sei $(T_n)_{n \in \mathds{N}} \subseteq \mathcal{L}(X,Y)$ eine Cauchy-Folge bezüglich $\norm[\mathcal{L}(X,Y)]{.}$. Für $x \in X$ ist dann auch $(T_n x)_{n \in \mathds{N}} \subset Y$ Cauchy
bezüglich $\norm[Y]{.}$. (Warum?)\footnote{mit $\norm[Y]{T_n x- T_m x} \le \norm{T_n -T_m} \cdot \norm{x}$}
Da $Y$ ein Banachraum ist, folgt $T_n x \xrightarrow{n \to \infty} y$ für ein $y \in Y$. Wir definieren $T : X \to Y$ durch $T(x) := \lim_{ n \to \infty} T_n(x)$ für 
$x \in X$. $T$ ist linear: Klar (Warum?).

$T$ ist stetig: $(T_n)_{n \in \mathds{N}}$ ist Cauchy, also auch $\enbrace*{\norm{T_n}}_{n \in \mathds{N}} \subset \mathds{R}$ eine Cauchyfolge (Warum?)\footnote{mit 
$\norm{T_n} = \norm{T_n - T_{n_0}+ T_{n_0}} \le \norm{T_n- T_{n_0}} + \norm{T_{n_0}}$}. Dann existiert $C \ge 0$ mit $\norm[\mathcal{L}(X,Y)]{T_n} \le C $, $n \in \mathds{N}$, also
gilt $\norm[Y]{T x} \le C \cdot \norm{x}$, $x \in X$, woraus $\norm[\mathcal{L}(X,Y)]{T} \le C$ und die Stetigkeit von $T$ folgt. Also ist $T \in \mathcal{L}(X,Y)$.

\noindent Falls $\varepsilon >0$, so existiert $n_0 \in \mathds{N}$ mit $\norm[\mathcal{L}(X,Y)]{T_n - T_m}  < \varepsilon$ für $n,m \ge n_0$. Insbesondere: 
\begin{align*}
	\norm[Y]{(T_n- T_m)x} \le \norm[\mathcal{L}(X,Y)]{T_n -T_m} \cdot \norm[X]{x} < \varepsilon \cdot \norm[X]{x}    
\end{align*}
für $n,m \ge n_0$, $x \in X$. Weiter gilt 
\begin{align*}
	\norm[Y]{(T_n- T)x} \le \norm[Y]{(T_n - T_m)x} + \norm[Y]{(T_m-T)x} &\le \norm[\mathcal{L}(X,Y)]{T_n - T_m} \cdot \norm[X]{x} + \norm[Y]{(T_m-T)x} \\
	&\le \varepsilon \cdot \norm[X]{x} + \varepsilon \cdot \norm[X]{x}  
\end{align*}
für $n,m \ge n_0$ und $m$ groß genug, $x \in X$. Damit folgt nun
$\norm[\mathcal{L}(X,Y)]{T_n -T} \le \varepsilon$ für$n \ge n_0$, also $T_n \to T$ bezüglich $\norm[\mathcal{L}(X,Y)]{.}$. \bewende
% subsection 31 (end)

\subsection[Corollar: Dualraum ist vollständig und $\mathcal{L}(X,X)$ ist Banachalgebra, falls $X$ Banachraum]{Corollar} % (fold)
\label{sub:32}
\begin{enumerate}[(i)]
	\item Der Dualraum $X^*$ eines normierten $\mathds{K}$-Vektorraumes $X$ ist vollständig. 
	\item Falls $\set{0} \not= X$ ein Banachraum ist, so ist $\mathcal{L}(X,X)$ eine Banachalgebra.
\end{enumerate}
\minisec{Beweis}
\begin{enumerate}[(i)]
	\item $X^*=\mathcal{L}(X,\mathds{K})$ und $\mathds{K}$ ist vollständig. Wende \ref{sub:31} an.
	\item $\mathcal{L}(X,X)$ ist ein Banachraum nach \ref{sub:31}. $\mathcal{L}(X,X)$ ist eine normierte Algebra: 
	\begin{itemize}
		\item $\norm[\mathcal{L}(X,X)]{\id_X}=1$ \hfill (Warum?)
		\item Für $S,T \in \mathcal{L}(X,X)$, $x \in X$ gilt
		\begin{align*}
			\norm[X]{(S \circ T)x} \le \norm{S} \cdot \norm[X]{T x} \le \norm{S} \cdot \norm{T} \cdot \norm[X]{x}   
		\end{align*}
		Also ist $\norm{S \circ T} \le \norm{S} \cdot \norm{T}$. \bewende
	\end{itemize}
\end{enumerate}
% subsection 32 (end)

\subsection[Definition und Proposition: lineare Isometrie $\iota_X : X \to X^{**}$]{Definition und Proposition} % (fold)
\label{sub:33}
Sei $X$ ein normierter Raum. Definiere $\iota_X \colon X \to X^{**} \enspace{\color{light_gray}\enbrace[\big]{=(X^*)^*}}$ durch
\[
	x \longmapsto \enbrace[\big]{\varphi \mapsto \varphi(x) } 
\]
für $x \in X$. Die Abbildung $\iota_X$ ist eine lineare Isometrie.\index{kanonische Inklusion}
\minisec{Beweis}
\begin{description}
	\item[Wohldefiniertheit:] Die Linearität von $\iota_X(x)$ ist klar. Mit
	\begin{align*}
		\abs*{\iota_X(x)(\varphi)} = \abs*{\varphi(x)} \le \norm[X^*]{\varphi} \cdot \norm[X]{x} = \norm[X]{x} \cdot \norm[X^*]{\varphi}      
	\end{align*}
	folgt $\varphi \in X^*$ und somit $\norm[X^{**}]{\iota_X(x)} \le \norm[X]{x}$, also ist $\iota_X(x)$ stetig.
	\item[$\iota_X$ ist linear:] Für $\alpha, \beta \in \mathds{K}$, $x,y \in X$, $\varphi \in \mathcal{L}(X,\mathds{K})= X^*$ gilt
	\begin{align*}
		\iota_X(\alpha \cdot x + \beta \cdot y)(\varphi)= \varphi(\alpha \cdot x + \beta \cdot y) = \alpha \circ \varphi(x) + \beta \cdot \varphi(y) 
		&= \alpha \cdot \iota_X(x)(\varphi) + \beta \cdot \iota_X(y)(\varphi) \\
		&= \enbrace[\big]{\alpha \cdot \iota_X(x) + \beta \cdot \iota_X(y)}(\varphi) 
	\end{align*}
	\item[Isometrie:] Es bleibt zu zeigen: $\norm[X^{**}]{\iota_X(x)} \ge \norm[X]{x}$ für $0 \not= x \in X$. Sei $Y := \mathds{K}\cdot x$ und $\psi : Y \to \mathds{K}$
	gegeben durch $\psi(\alpha \cdot x) := \alpha \cdot \norm[X]{x}$. Dann ist $Y \subset X$ ein linearer Unterraum und $\psi(Y) \to \mathds{K}$ linear mit
	$\norm{\psi}=1$. Nach Satz \ref{sub:224} existiert $\overline{\psi}. X \to \mathds{K}$ linear mit $\overline{\psi}(x) = \psi(x) = \norm[X]{x}$ und 
	$\norm{\overline{\psi}}= \norm{\psi} = 1$. Es gilt 
	\begin{align*}
		\norm[X^{**}]{\iota_X(x)} \ge \abs*{\iota_X(x)(\overline{\psi})} = \abs*{\overline{\psi}(x)} = \norm[X]{x}. \bewende
	\end{align*}
\end{description}
% subsection 33 (end)

\subsection[Definition und Proposition: Die transponierte Abbildung]{Definition und Proposition} % (fold)
\label{sub:34}
Seien $X,Y$ normierte Räume. Definiere eine Abbildung $.^{\tr} \colon \mathcal{L}(X,Y) \to \mathcal{L}(Y^*,X^*)$ durch $T \mapsto T^{\tr} = (\varphi \mapsto \varphi \circ T)$.
$.^\tr$ ist eine lineare Isometrie und das Diagramm\index{transponierte Abbildung}
\[
	\begin{tikzcd}
		X \rar["\iota_X"] \dar["T"]& X^{**} \dar["T^{\tr \, \tr}= (T^\tr)^\tr"]\\
		Y \rar["\iota_Y"]& Y^{**}
	\end{tikzcd}
\]
kommutiert für jedes $T \in \mathcal{L}(X,Y)$. Für $\iota_{X^*} \colon X^* \to X^{***}$ und $\iota_{X}^\tr \colon X^{***} \to X^*$ gilt $\iota_X^\tr \circ \iota_{X^*} = \id_{X^*}$.
\minisec{Beweis}
\begin{enumerate}[(i)]
	\item Für $T,S \in \mathcal{L}(X,Y)$, $\alpha, \beta \in \mathds{K}$, $\varphi \in Y^*$ gilt
	\begin{align*}
		\enbrace*{\alpha \cdot T + \beta \cdot S}^\tr (\varphi) = \varphi \circ \enbrace*{\alpha \cdot T + \beta \cdot S} \!\!\stackrel{\text{linear}}{=} \!\!
		\alpha \cdot (\varphi \circ T) + \beta \cdot (\varphi \circ S)  &= \alpha \cdot T^\tr(\varphi) + \beta \cdot S^\tr(\varphi)\\
		&= \enbrace*{\alpha \cdot T^\tr + \beta \cdot S^\tr}(\varphi) 
	\end{align*}
	\item Für $x \in X$, $\psi \in Y^*$ gilt $\enbrace[\big]{\iota_Y T(x)}(\psi) = \psi \enbrace[\big]{T(x)}$; weiter gilt
	\begin{align*}
		\enbrace[\big]{T^{\tr \, \tr} \iota_X(x)} (\psi) = \enbrace[\big]{(T^\tr)^\tr \iota_X(x)}(\psi)  = \enbrace[\big]{\iota_X(x) \circ T^\tr}(\psi)
		&= \iota_X(x) \enbrace*{T^\tr(\psi)}  \\
		&= T^\tr(\psi)(x) = \psi \circ T(x) \\
		&= \iota_Y \enbrace[\big]{T(x)} (\psi) 
	\end{align*}
	$\Rightarrow T^{\tr \, \tr} \iota_X(x) = \iota_Y \enbrace[\big]{T(x)} = \iota_Y \circ T (x)$. Also $T^{\tr \, \tr} \circ \iota_X = \iota_Y  \circ T$. Also kommutiert
	das Diagramm.
	\item Es ist $\abs{(T^\tr \psi)(x)} = \abs*{\psi T(x)} \le \norm[Y^*]{\psi} \cdot \norm[Y]{T x} \le \norm[Y^*]{\psi} \cdot \norm[\mathcal{L}(X,Y)]{T} \cdot \norm[X]{x}$. Also
	ist $\norm[X^*]{T^\tr \psi} \le  \norm[Y^*]{\psi} \cdot  \norm[\mathcal{L}(X,Y)]{T}$ und damit folgt 
	\[
		\norm[\mathcal{L}(Y^*,X^*)]{T^\tr} \le \norm[\mathcal{L}(X,Y)]{T}  
	\]
	Ebenso ist $\norm[\mathcal{L}(X^{**},Y^{**})]{T^{\tr \, \tr}} \le \norm[\mathcal{L}(Y^*,X^*)]{T^\tr} \le \norm[\mathcal{L}(X,Y)]{T}$. Andererseits gilt 
	$\norm[\mathcal{L}(X,Y)]{T} \le \norm[\mathcal{L}(X^{**},Y^{**})]{T^{\tr\, \tr}}$ nach \ref{sub:33} und da das Diagramm kommutiert (Warum?). 
	\[
		\Rightarrow \norm[\mathcal{L}(X,Y)]{T} \le \norm[\mathcal{L}(Y^*,X^*)]{T^\tr}  \le  \norm[\mathcal{L}(X,Y)]{T} 
	\]
	also ist $.^\tr$ eine Isometrie.
	\item Für $\varphi \in X^*$, $x \in X$ gilt
	\begin{align*}
		\enbrace[\big]{\iota_X^\tr \circ \iota_{X^*} (\varphi)} (x) = \iota_{X^*}(\varphi) \circ \iota_X (x) = \iota_X(x)(\varphi) = \varphi(x)
	\end{align*}
	$\Rightarrow \iota_X^\tr \circ \iota_{X^*}(\varphi)= \varphi$, also $\iota_X^\tr \circ  \iota_{X^*} = \id_{X^*}$. \bewende
\end{enumerate}
% subsection 34 (end)

\subsection{Satz: Prinzip der gleichmäßigen Beschränktheit für Banachräume} % (fold)
\label{sub:35}
Sei $X$ ein Banachraum und $Y$ ein normierter Raum. Sei $M \subset \mathcal{L}(X,Y)$, sodass die Menge 
\[
	\set[{x \mapsto \norm[Y]{T x}}]{T \in M \rule{0cm}{1em}} \subset C(X,\mathds{R} )
\]
punktweise \Index{gleichmäßig beschränkt} ist, d.h. für $x \in X$ existiert $C_x \ge 0$ mit $\norm[Y]{T x} \le C_x$, $T \in M$. Dann existiert $C \ge 0$ mit $\norm[\mathcal{L}(X,Y)]{T} \le C$,
$T \in M$.
\minisec{Beweis}
Nach Corollar \ref{sub:112} (Prinzip der gleichmäßigen Beschränktheit für metrische Räume) existieren eine offene Kugel $\emptyset \not= B(x_0,\varepsilon) \subset X$ und 
$K \ge 0$ mit $\norm[Y]{T x} \le K$ für $x \in B(x_0, \varepsilon)$, $T \in M$. Für $x \in X$ mit $\norm{x} \le 1$ gilt 
\begin{align*}
	\norm[X]{T x} = \frac{e}{\varepsilon} \cdot  \norm[Y]{T \enbrace*{\frac{\varepsilon}{2} \cdot x}} &= \frac{2}{\varepsilon} \cdot 
	\norm[Y]{T \enbrace*{\frac{\varepsilon}{2} \cdot x + x_0 - x_0 } }  \\
	&\le \frac{2}{\varepsilon} \cdot   \enbrace*{ \doppelstrich[\Bigg]{T \enbrace[\Bigg]{\underbrace{\frac{\varepsilon}{2} \cdot x + x_0}_{\in B(x_0,\varepsilon)} } }_Y +
	 \norm[Y]{T  \enbrace[\bigg]{\underset{\in B(x_0,\varepsilon)}{x_0}} }} \\
	&\le \frac{2}{\varepsilon} (K+K) = \frac{4K}{\varepsilon}   
\end{align*}
$\Rightarrow \norm[\mathcal{L}(X,Y)]{T} \le \frac{4 K}{\varepsilon} =: C$. \bewende
% subsection 35 (end)

\subsection{Corollar} % (fold)
\label{sub:36}
Sei $Z$ ein metrischer Raum und $N \subset Z$ eine Teilmenge, so dass gilt: Für alle $\varphi \in Z^*$ existiert $C_\varphi \ge 0$ mit $\abs*{\varphi(z)} \le C_\varphi$, 
$z \in N$.  Dann ist $N$ beschränkt.
\minisec{Beweis}
$\iota_Z : Z \to Z^{**}$ ist eine Isometrie nach \ref{sub:33}. Wende Satz \ref{sub:35} an mit $X=Z^{*}$, $Y=\mathds{K}$ und 
$M = \iota_Z(N) \subset Z^{**}= \mathcal{L}(Z^*,\mathds{K})$.
Für $\varphi \in X = Z^{*}$ existiert $C_\varphi \ge 0$ mit $\abs*{T \varphi} = \abs*{\iota_Z(z) (\varphi)} = \abs*{\varphi(z)} \le C_\varphi$, wo $T= \iota_Z(z) \in M$
für ein $z \in N$. Nach \ref{sub:35} folgt: Es existiert ein $C \ge 0$ mit $\norm[Z]{z} = \norm[Z^{**}]{\iota_Z(z)}  \le C$, $z \in N$. \bewende
% subsection 36 (end)

\subsection{Corollar} % (fold)
\label{sub:37}
Sei $X$ ein Banachraum und $Z$ ein normierter Raum. $M \subset \mathcal{L}(X,Z)$ eine Teilmenge, so dass gilt: Für alle $\varphi \in Z^*$ und $x \in X$ existiert $C_{x,\varphi} \ge 0$
mit
\[
	\abs*{\varphi \enbrace[\big]{T(x)}} \le C_{x,\varphi}, \quad T \in M
\]
Dann ist $M$ beschränkt.
\minisec{Beweis}
Für jedes $x \in X$ folgt aus Corollar \ref{sub:36} mit $N := \set[T(x)]{T \in M}$, dass $N$ beschränkt ist. Es folgt, dass $ \set[x \mapsto \norm{T(x)} ]{T \in M}$ ist 
punktweise gleichmäßig beschränkt. Mit Satz \ref{sub:35} folgt: $M$ ist beschränkt. \bewende
% subsection 37 (end)

\subsection{Satz: Prinzip der offenen Abbildung} % (fold)
\label{sub:38}
Seien $X,Y$ Banachräume, $T \in \mathcal{L}(X,Y)$ surjektiv. Dann ist $T$ offen, d.h. für eine offene Menge $U \subset X$ ist $T(U) \subset Y$ offen.
\minisec{Beweis}
\begin{enumerate}[(i)]
	\item Sei $0 \in W \subset X$ offen, dann existiert $\emptyset \not= V \subset Y$ offen mit $V \subset \overline{T(W)}$: $0 \in W \subset X$ $\Rightarrow $ 
	$B(0,\varepsilon) \subset W$ für ein $\varepsilon>0$. Dann gilt
	\[
		X=\bigcup_{n \in \mathds{N}} n \cdot B(0,\varepsilon) = \bigcup_{n \in \mathds{N}} n \cdot W
	\]
	Aus der Surjektivität von $T$ folgt $Y = \bigcup_{n \in \mathds{N} }n \cdot T(W) = \bigcup_{n \in \mathds{N}} n \cdot \overline{T(W)}$. Nach dem Satz von Baire 
	(\ref{sub:110}) folgt: Ein $n \cdot \overline{T(W)}$ enthält eine nichtleere offene Teilmenge, also auch $\overline{T(W)}$.
	\item Sei $= \in M \subset X$ offen, dann ist $0$ ein innerer Punkt von $\overline{T(M)}$: 
	
	Die Abbildung $X \times X \to X$, $(x_1,x_2) \mapsto x_1 -x_2$ ist stetig.\footnote{Banachräume sind topologische Vektorräume} $\Rightarrow$ es gibt $0 \in W \subset X$ 
	offen mit
	\[
		W - W = \set[x_1 -x_2]{x_1,x_2 \in W} \subset M 
	\]
	Es folgt 
	\[
		\overline{T(M)} \supset \overline{T(W-W)} = \overline{T(W)- T(W)} \stackrel{\text{(Warum?)}}{\supset} \overline{T(W)}- \overline{T(W)} \stackrel{\text{(i)}}{\supset}
		V -V
	\]
	für $\emptyset \not= V$ offen in $Y$ wie in (i). Aber $0 \in V-V = \bigcup_{y \in V}\set{y}- V$ ist offen in $Y$. Also ist $0$ ein innerer Punkt von $\overline{T(M)}$.
	\item Ist $0 \in N \subset X$ offen, so existiert $0 \in Z \subset Y$ offen mit $Z \subset T(M)$:
	
	Sei $\varepsilon_0 >0$. Wähle $\varepsilon_i >0$, $1 \le i \in \mathds{N}$ mit $\sum_{i=1}^{\infty} \varepsilon_i < \varepsilon_0$. Nach (ii) gilt: Zu $\varepsilon_i$, 
	$i \in N$ existiert ein $\delta_i >0$ mit $B_Y(0,\delta_i) \subset \overline{T \enbrace*{B_X(0,\varepsilon_i)}}$. Es gilt $\delta_i \xrightarrow{i \to \infty} 0$. Sei 
	nun $y \in B_Y(0,\delta_0) =: Z$. Dann gilt $y \in \overline{T \enbrace*{B_X(0,\varepsilon_0)}}$. $\Rightarrow $ es existiert $x_0 \in B_X(0,\varepsilon)$ mit 
	$\norm[Y]{y- T(x_0)}< \delta_1$. Folglich existiert $x_1 \in B_X(0,\varepsilon_1)$ mit $\norm[Y]{y- T(x_0) - T(x_1)} < \delta_2$. Per Induktion folgt
	\[
		\exists x_i \in B_X(0,\varepsilon_i) \text{ mit } \norm[Y]{y- \sum\nolimits_{j=0}^{i}T(x_i)} < \delta_{i+1}, \qquad i \in \mathds{N} 
	\]
	$\sum_{i \in \mathds{N}}x_i$ konvergiert in $X$ und $\norm[X]{\sum_{i\in \mathds{N}} x_i} < 2 \cdot \varepsilon_0 $, es gilt $T \enbrace*{\sum_{i \in \mathds{N}}x_i}=y$.
	($x_i \in N$, falls $B(0,2 \varepsilon_0) \subset N$)
	\item Sei $U \subset X$ offen. Dann ist $T(U) \subset Y$ offen:
	
	Zu $x \in U$ existiert $0 \in N \subset X$ offen mit $x+ N \subset U$, also $T(x)+ T(N)\subset T(U)$. Nach (iii) existiert $0 \in Z \subset Y$ offen mit $Z \subset T(N)$,
	also $T(x)+ Z \subset T(U)$. Aber $T(x) \in T(x)+ Z$ ist offen in $Y$ und $T(U)$ ist Umgebung von $T(x)$. $x \in U$ war beliebig, also ist $T(U)$ offen. \bewende
\end{enumerate}
% subsection 38 (end)

\subsection{Corollar: Satz von der inversen Abbildung} % (fold)
\label{sub:39}
Seien $X,Y$ Banachräume, $T \in \mathcal{L}(X,Y)$ bijektiv. Dann gilt $T ^{-1} \in \mathcal{L}(Y,X)$, d.h. $T$ ist ein Homöomorphismus\index{Homöomorphismus}.
\minisec{Beweis}
Nach Satz \ref{sub:38} ist $T$ offen und somit ist $T ^{-1}$ stetig. \bewende
% subsection 39 (end)

\subsection{Corollar} % (fold)
\label{sub:310}
Sei $X$ ein Vektorraum und seien $\mathcal{T}_1, \mathcal{T}_2$ zwei Topologien auf $X$, die durch Banachraumnormen $\norm[1]{\cdot }$ und $\norm[2]{\cdot}$ induziert sind.
Falls $\mathcal{T}_1 \subset \mathcal{T}_2$, so gilt bereits $\mathcal{T}_1 = \mathcal{T}_2$.
\minisec{Beweis}
$\mathcal{T}_1 \subset \mathcal{T}_2$ $\Rightarrow $ $(X,\mathcal{T}_2) \xrightarrow{\id_X} (X,\mathcal{T}_1)$ ist stetig und bijektiv, also ist $\id_X$ nach 
\ref{sub:39} ein Homöomorphismus. \bewende
% subsection 310 (end)

\subsection{Corollar} % (fold)
\label{sub:311}
Sei $X$ ein Banachraum mit abgeschlossenen Teilräumen $X_1,X_2$, sodass $X_1 \cap X_2 = \set{0}$ und $X= X_1 + X_2$. Dann sind $X$ und $X_1 \times X_2$ als topologische
Vektorräume isomorph.
\minisec{Beweis}
Die Abbildung $X_1 \times X_2 \to X$, $(x_1,x_2) \mapsto x_1 + x_2$ ist linear und bijektiv. Sie ist stetig, wegen 
\[
	\norm[X]{x_1 + x_2} \le \norm[X]{x_1} + \norm[X]{x_2} = \norm[1,X_1 \times X_2]{(x_1,x_2)} 
\]
Mit Corollar \ref{sub:39} folgt die Behauptung. \bewende
% subsection 311 (end)
% section 3 (end)
\newpage

\section{$L^p$-Räume und der Satz von Riesz-Fischer} % (fold)
\label{sec:4}

\todo{RevChap4}
\subsection[Erinnerung: $L^1(\mu)$ ist ein Banachraum]{Erinnerung} % (fold)
\label{sub:41}
Sei $(X,\Sigma, \mu)$ ein \Index{Maßraum}.\footnote{$X$ Menge, $\Sigma$ eine $\sigma$-Algebra, $\mu$ Maß}  $f : X \to (-\infty,\infty]$ heißt \Index{messbar}, falls 
$f ^{-1}(E) \in \Sigma$ gilt für $E \subset \mathcal{B}(\mathds{R})$ (Borelmengen\todo{Def. hinzufügen}). $s : X \to (-\infty,\infty]$ heißt \Index{einfach}, falls 
$s(X)$ endlich. Für $f : X \to [0,\infty]$ messbar setze
\[
	\int f  \, \mathd \mu := \sup \set[\int s \, \mathd \mu]{s \text{ einfach, messbar }, 0 \le s \le f} 
\] 
Für $f : X \to (-\infty,\infty]$ messbar setze $\int\! f \, \mathd \mu := \int \! f_+ \, \mathd \mu - \int \! f_- \, \mathd \mu$, falls $\int \! f_- \, \mathd \mu < \infty$.
Betrachte
\begin{align*}
	\mathcal{L}^1(\mu) := \set[f : X \to \mathds{R} \text{ messbar}]{\int \abs*{f}\,\mathd \mu < \infty } 
\end{align*}
mit der Halbnorm $\norm[1]{f} := \int \! \abs*{f} \, \mathd \mu $. Setze nun $\mathcal{N}^1 := \set[f : X \to \mathds{R} \text{ messbar}]{\int \abs*{f} \, \mathd \mu =0 }$.
$L^1(\mu)$ ist ein topologischer Vektorraum (Topologie induziert durch $\norm[1]{\cdot}$) und $\mathcal{N}^1 \subset \mathcal{L}^1(\mu)$ ist ein abgeschlossener 
Untervektorraum (Übung). Definiere 
\[
	L^1(\mu) := \nicefrac{\mathcal{L}^1}{\mathcal{N}^1}
\]
$\norm[1]{\cdot }$ induziert eine Norm $\norm[1]{\cdot }$ auf $L^1(\mu)$ via  
\[
	\norm[1]{f+ \mathcal{N}^1} := \norm[1]{f} \quad , \quad f \in \mathcal{L}^1(\mu) \tag*{(Übung)}
\]
$(L^1(\mu),\norm[1]{\cdot})$ ist ein Banachraum (Ana \RM{3}, 9.23).
% subsection 41 (end)

\subsection[Beispiele für $L^1$-Räume]{Beispiele} % (fold)
\label{sub:42}
\begin{enumerate}[(i)]
	\item $L^1(\mathds{R}^n) := L^1(\mathds{R}^n, \Lambda, \lambda)$, wobei $\Lambda$ die Lebesgue-messbaren Mengen sind und $\lambda $ das Lebesgue-Maß ist.
	\item $\ell^1(\mathds{N})= L^1 \enbrace[\big]{\mathds{N}, \mathcal{P}(\mathds{N}), \delta} = \set[f : \mathds{N} \to \mathds{R}]{\sum \abs*{f(n)} < \infty } $, 
	wobei $\delta$ das Zählmaß\footnote{siehe auch \url{http://de.wikipedia.org/wiki/Zählmaß_(Maßtheorie)}} ist.
\end{enumerate}
% subsection 42 (end)

\subsection[Definition und Proposition: Der Raum $L^p(\mu)$]{Definition und Proposition} % (fold)
\label{sub:43}
Für einen Maßraum $(X,\Sigma,\mu)$ definiere 
\begin{align*}
	\mathcal{L}^p(\mu) &:= \set[f : X \to \mathds{R} \text{ messbar}]{\int \abs*{f}^p \mathd \mu < \infty } \\
	\mathcal{N}^p &:= \set[f : X \to \mathds{R} \text{ messbar}]{\int \abs*{f}^p \mathd \mu = 0 } 
\end{align*}
$\mathcal{L}^p(\mu)$ und $\mathcal{N}^p$ sind Vektorräume; wir setzen 
\[
	L^p(\mu) := \nicefrac{\mathcal{L}^p(\mu)}{\mathcal{N}^p}
\]
Definiere $\norm[p]{\cdot } : \mathcal{L}^p(\mu) \to [0,\infty)$ durch $\norm[p]{f} := \enbrace*{\int \abs*{f}^p \mathd \mu }^{1/p}$. $\norm[p]{\cdot }$ induziert eine 
Abbildung $\norm[p]{\cdot } : L^P(\mu) \to [0,\infty) $ via $\norm[p]{f + \mathcal{N}^p} := \norm[p]{f} $, für $f \in \mathcal{L}^p$.
\minisec{Beweis}
\begin{description}
	\item[$\mathcal{L}^p(\mu)$ ist Vektorraum:]  Für $f,g \in \mathcal{L}^p(\mu)$ und $\alpha, \beta \in \mathds{R}$
	\begin{align*}
		\int \abs*{\alpha \cdot f + \beta \cdot g}^p \mathd \mu \le \int \!\enbrace[\big]{\abs*{\alpha \cdot f} + \abs*{\beta \cdot g}}^p \mathd \mu
		&\le \int \!\enbrace[\big]{2 \cdot \sup \set{\abs*{\alpha \cdot f}, \abs*{\beta \cdot g}  } }^p \mathd \mu \\ 
		&\le \int \enbrace*{2^p \cdot \sup \set{\abs*{\alpha \cdot f}^p, \abs*{\beta \cdot g}^p  \rule{0cm}{1em}} } \mathd \mu \\ 
		&\le \int \!2^p \cdot  \enbrace[\big]{\abs*{\alpha \cdot f}^p + \abs*{\beta \cdot g}^p} \mathd \mu < \infty 
	\end{align*}
	\item[$\mathcal{N}^p$ ist Vektorraum:] Übung.
	\item[{$\norm[p]{\cdot}$ wohldefiniert auf $L^p(\mu)$:}] Übung. \bewende
\end{description}
% subsection 43 (end)

\subsection[Proposition: Speziallfall der Youngschen Ungleichung]{Proposition} % (fold)
\label{sub:44}
Für $\alpha, \beta \ge 0$, $p,q >1$ mit $\frac{1}{p} + \frac{1}{q}=1$. Dann gilt 
\begin{equation*}
	\alpha \cdot \beta \le \frac{\alpha^p}{p} + \frac{\beta^q}{q}  \label{eq:44}\tag{\#}
\end{equation*}
Dies ist ein Spezialfall der Youngschen Ungleichung,  \hrefsym{https://de.wikipedia.org/wiki/Youngsche_Ungleichung_(Produkt)}{Wikipedia-Link}.\index{Young-Ungleichung}
\minisec{Beweis}
\begin{align*}
	\frac{1}{p} + \frac{1}{q}=1 \iff p+q= pq \iff (p-1)(q-1)= 1 \iff p (q-1) = q \iff q (p-1) =p
\end{align*}
Falls $\beta = \alpha^{p-1}$, so gilt $\alpha \cdot \beta = \alpha^p$ und 
\[
	\frac{\alpha^p}{p} + \frac{\beta^q}{q} = \frac{\alpha^p}{p} + \frac{(\alpha^{p-1})^q}{q} = \frac{\alpha^p}{p} + \frac{\alpha^p}{q} = \alpha^p = \alpha \cdot \beta   
\]
Also gilt Gleichheit in \eqref{eq:44}. Für festes $\alpha \ge 0$ betrachte $f_\alpha : [0,\infty) \to [0,\infty)$ gegeben durch
$f_\alpha(\beta) := \frac{\alpha^p}{p} + \frac{\beta^q}{q} - \alpha \cdot \beta$. Dann ist $f_\alpha$ stetig auf $[0,\infty)$ und differenzierbar auf $(0,\infty)$. Es gilt
$f_\alpha'(\beta) = \beta^{q-1} -\alpha$ und weiter
\[
	f_\alpha'(\beta)=0 \iff \beta^{q-1} = \alpha \iff \beta= \alpha^{p-1}
\]
$f_\alpha''(\alpha^{p-1}) \ge 0$, also hat $f_\alpha$ bei $\beta = \alpha^{p-1}$ ein lokales Minimum. Es gilt $f_\alpha(\alpha^{p-1})= 0 \le f_\alpha(0)$, also
hat $f_\alpha$ bei $\beta=\alpha^{p-1}$ ein globales Minimum. $\Rightarrow $ \eqref{eq:44}. \bewende
% subsection 44 (end)

\subsection{Satz: Höldersche Ungleichung} % (fold)
\label{sub:45}
Sei $(X,\Sigma,\mu)$ ein Maßraum. Seien $p,q \ge 1$ mit $\frac{1}{p} +\frac{1}{q}=1 $, $f \in \mathcal{L}^p(\mu)$, $g \in \mathcal{L}^q(\mu)$. Dann gilt 
$f \cdot g \in \mathcal{L}^1{\mu}$ und
\[
	\norm[1]{f \cdot g} \le \norm[p]{f} \cdot \norm[q]{g}   
\]
Die entsprechenden Aussagen gelten dann auch für $L^p(\mu)$, $L^q(\mu)$, $L^1(\mu)$.\index{Hölder-Ungleichung}
\minisec{Beweis}
Klar, falls $\norm[p]{f} =0 $ oder $\norm[q]{g} =0$ (Warum?). Sei also $\norm[p]{f}, \norm[q]{g} \not= 0$. Nach Proposition \ref{sub:44} gilt für $x \in X$
\[
	0 \le \underbrace{\underbrace{\frac{\abs*{f(x)} }{\norm[p]{f}}}_{=\alpha} \cdot \underbrace{\frac{\abs*{g(x)} }{\norm[q]{g} }}_{=\beta} }_{\text{messbar}}
	\le \underbrace{\frac{1}{p} \cdot \frac{\abs*{f(x)}^p }{\norm[p]{f}^p } + \frac{1}{q} \cdot \frac{\abs*{g(x)}^q }{\norm[q]{g}^q }}_{\text{integrierbar}}
\]
$\abs*{f}^p , \abs*{g}^q  \in \mathcal{L}^1(\mu)$, also $\abs*{ f \cdot g } \in \mathcal{L}^1(\mu)$. Es gilt
\begin{align*}
	\frac{\int \abs*{f \cdot g} \mathd \mu}{\norm[p]{f} \cdot \norm[q]{g}} \le \frac{1}{p} \cdot \frac{\int \abs*{f}^p \mathd \mu }{\norm[p]{f}^p } +
	\frac{1}{q} \cdot \frac{\int \abs*{g}^q \mathd \mu }{\norm[q]{g}^q} = \frac{1}{p} + \frac{1}{q} = 1       
\end{align*}
$\Rightarrow \inf \abs*{f \cdot g} \mathd \mu  \le \norm[p]{f}  \cdot \norm[q]{g}$. \bewende
% subsection 45 (end)

\subsection{Satz: Minkowskische Ungleichung} % (fold)
\label{sub:46}
Für $1\le p < \infty$, $f,g \in \mathcal{L}^p(\mu)$ (bzw. $L^p(\mu)$) gilt die \Index{Minkowski-Ungleichung}
\[
	\norm[p]{f +g} \le \norm[p]{f} + \norm[p]{g}   
\]
\minisec{Beweis}
Sei $\norm[p]{f+g} >0 $ (sonst trivial). Es gilt 
\begin{align*}
	\abs*{f +g}^p = \abs*{f+g}^{p-1} \cdot \abs*{f+g} \le \abs*{f+g}^{p-1} \enbrace[\big]{\abs*{f} + \abs*{g}}     
\end{align*}
Sei $q>1$ mit $\frac{1}{p} + \frac{1}{q}=1$. Es gilt $\abs*{f+g}^{p-1} \in \mathcal{L}^q(\mu)$, da $\abs*{f+g}^{(p-1)q} = \abs*{f+g}^p \in \mathcal{L}^1(\mu)$. Nach Hölder
\ref{sub:45} gilt
\begin{align*}
	\int \!\abs*{f+g}^p \mathd \mu \le \int \!\abs*{f+g}^{p-1} \abs*{f} \mathd \mu + \int \!\abs*{f+g}^{p-1} \abs*{g}\mathd \mu &\le \norm[p]{f} \cdot \norm[q]{\abs*{f+g}^{p-1}}
	+ \norm[p]{g} \cdot \norm[q]{\abs*{f+g}^{p-1} } \\        
	&= \enbrace*{\norm[p]{f} + \norm[p]{g}} \cdot \enbrace*{\int \abs*{f+g}^{(p-1)q} \mathd \mu }^{\frac{1}{q}} \\
	&= \enbrace*{\norm[p]{f} + \norm[p]{g}} \cdot \enbrace*{\int \abs*{f+g}^{p} \mathd \mu }^{q} \\
\end{align*}
Also gilt 
\[
	\enbrace*{\int \abs*{f+g}^p \mathd \mu}^{1- \frac{1}{q}} = \enbrace*{\int \abs*{f+g}^{p} \mathd \mu }^{\frac{1}{p}} \le \norm[p]{f} + \norm[p]{g} \bewende    
\]
% subsection 46 (end)

\subsection[Corollar: $L^p(\mu)$ ist ein normierter Vektorraum für $1 \le p < \infty$]{Corollar} % (fold)
\label{sub:47}
$L^p(\mu)$ ist ein normierter Vektorraum für $1 \le p < \infty$.
% subsection 47 (end)

\subsection[Definition und Proposition: Der normierte Raum $L^\infty(\mu)$ und ]{Definition und Proposition} % (fold)
\label{sub:48}
\begin{align*}
	\mathcal{L}^\infty(\mu) &:= \set[f : X \to \mathds{R} \text{ messbar}]{\exists C \ge 0 : \abs*{f} \le C \text{ fast überall} \rule{0cm}{0.7em} }  \\
	\mathcal{N}^\infty &:= \set[f : X \to \mathds{R}]{f=0 \text{ fast überall}} 
\end{align*}
Wir definieren $L^\infty(\mu) := \nicefrac{\mathcal{L}^\infty(\mu)}{\mathcal{N}^\infty}$. Wir definieren eine Norm auf $\mathcal{L}^\infty(\mu)$ durch
\[
	\norm[\infty]{f} := \inf \set[C]{\abs*{f} \le C \text{ fast überall} }\quad , \qquad f \in \mathcal{L}^\infty(\mu)  
\]
Dies induziert wieder eine Norm auf $L^\infty(\mu)$ durch $\norm[\infty]{f + \mathcal{N}} := \norm[\infty]{f}$. Damit ist $\enbrace*{L^\infty(\mu), \norm[\infty]{\cdot}}$
ein normierter Vektorraum. Für $f \in L^1(\mu)$, $g \in L^\infty(\mu)$ gilt $f \cdot g \in L^1(\mu)$ und $\norm[1]{f \cdot g} \le \norm[1]{f} \cdot \norm[\infty]{g}$
\minisec{Beweis}
Übung. \bewende
% subsection 48 (end)

\subsection[Beispiele für $L^p$-Räume]{Beispiele} % (fold)
\label{sub:49}
\begin{enumerate}[(i)]
	\item $(\mathds{R}^n, \norm[p]{\cdot})$, $1 \le p < \infty$, also $X=\set{1,\ldots,n}, \Sigma=\mathcal{P}(X), \mu =$ Zählmaß.
	\item $\enbrace[\big]{\ell^p(\mathds{N}), \norm[p]{\cdot }}$, $1 \le p \le \infty$, also $X=\mathds{N}, \Sigma=\mathcal{P}(\mathds{N}), \mu=$ Zählmaß.
	Dabei ist 
	\[
		\ell^p(\mathds{N}) = \set[(a_n)_n]{a_n \in \mathds{R}, \sum_{n \in \mathds{N}} \abs*{a_n}^p < \infty} 
	\]
	für $p < \infty$ bzw. $\ell^p(\mathds{N})= \set{\text{beschränkte Folgen}}$ für $p=\infty$.
\end{enumerate}
% subsection 49 (end)

\subsection[Proposition: Lineare Isometrie $\iota_p : L^p(\mu) \to L^q(\mu)^*$]{Proposition} % (fold)
\label{sub:410}
Sei $(X,\Sigma,\mu)$ ein Maßraum, $p,q >1$ mit $\frac{1}{p} + \frac{1}{q}=1$, bzw. $p=1$, $q=\infty$. Die Abbildung $\iota_p : L^p(\mu) \to L^q(\mu)^*$, gegeben durch
\[
	\iota_p(f)(g) := \int \! f \cdot g \, \mathd \mu
\]
ist eine lineare Isometrie.
\minisec{Beweis}
\begin{description}
	\item[$\iota_p(f) : L^q(\mu) \to \mathds{R}$ ist linear und wohldefiniert:] Übung!
	\item[$\iota_p$ linear und wohldefiniert:] Übung!
\end{description}
$\iota_p(f)$ ist beschränkt durch $\norm[p]{f}$ nach Hölder (\ref{sub:45}), also ist $\iota_p$ beschränkt durch $1$. Noch zu zeigen: Zu $f \in L^p(\mu)$ existiert 
$g \in L^q(\mu)$ mit $\abs*{\iota_p(f)(g)} = \norm[p]{f} \cdot \norm[q]{g}$. (Denn dann $\norm[L^q(\mu)^*]{\iota_p(f)}  \ge \norm[p]{f}$)

\noindent Gegeben $f \in L^p(\mu)$, setze $g := \sgn(f) \cdot \abs*{f}^{p-1}$. Dann ist $\abs*{g}^q= \abs*{f}^p$, also $g \in L^q(\mu)$. Weiter gilt
\begin{align*}
	\int \!f \cdot g \,\mathd \mu = \int \! \abs*{f \cdot g} \mathd \mu = \int \!\abs*{f}^p \mathd \mu  = \enbrace*{\int\! \abs*{f}^p \mathd \mu }^{\frac{1}{p}}
	\cdot  \enbrace*{\int \!  \abs*{f}^p \mathd\mu}^{\frac{1}{q}} = \norm[p]{f} \cdot \norm[q]{g}    
\end{align*}
Für $p=1$ und $q=\infty$ Übung. \bewende
% subsection 410 (end)

\subsection{Satz von Riesz-Fischer} % (fold)
\label{sub:411}
Sei $(X,\Sigma,\mu)$ ein Maßraum. $L^p(\mu)$ ist vollständig für $p \in [1,\infty]$.\index{Riesz-Fischer}
\minisec{Beweis}
Für $p=\infty$: Übung. Sei also $p <\infty$ und sei $(f_n)_{n \in \mathds{N}} \subset L^p(\mu)$ Cauchy.
\begin{description}
	\item[Behauptung:] Es gibt $E_l \subset X$, $l \in \mathds{N}$, messbare Teilmenge mit $\mu(E_l) <\infty$ und sodass für $E := \bigcup_{l \in \mathds{N}} E_l$ gilt 
	$\chi_{X \setminus E} \cdot f_n = 0$ in $L^p(\mu)$ für $n \in \mathds{N}$. Beweis: Übung!
\end{description}
Wähle $n_0 < n_1 < \ldots $ mit $\norm[p]{f_n - f_{n_k}} < \frac{1}{2^k}$ für $n > n_k$, $k \in \mathds{N}$. Für jedes $l,k \in \mathds{N}$ gilt
$\abs*{f_{n_{k+1}} - f_{n_k}} \in L^p(\mu)$ und $\chi_{E_l} \in L^q(\mu)$ (wobei $\frac{1}{p} + \frac{1}{q}=1$), also ist nach Hölder (\ref{sub:45}) auch
\(
	\abs*{f_{n_{k+1}} - f_{n_k}} \cdot \chi_{E_l} \in L^1(\mu)
\).
Weiter gilt für $j \in \mathds{N}$, $l \in \mathds{N}$
\[
	\sum_{k=0}^{j}  \doppelstrich[\Big]{\abs*{f_{n_{jk+1}}- f_{n_k}} \cdot \chi_{E_l}}_{1}
	\le \sum_{k=0}^{j} \norm[p]{f_{n_{k+1}}- f_{n_k}} \cdot \norm[q]{\chi_{E_l}} \le \sum_{k=0}^{j} \frac{1}{2^k} \cdot \norm[q]{\chi_{E_l}}
\]
$\enbrace*{\sum_{k=0}^{j} \abs*{f_{n_{k+1}} - f_{n_k}} \cdot \chi_{E_l}}_{j \in \mathds{N}}$ ist aufsteigende Folge von integrierbaren Funktionen. Nach dem Satz von
Lebesgue über monotone Konvergenz gilt: $\sum_{k=0}^{\infty} \abs*{f_{n_{k+1}} - f_{n_k}} $ konvergiert auf $E_l$ fast überall. Da abzählbare Vereinigungen von Nullmengen
Nullmengen sind, folgt: $\sum_{k=0}^{\infty} \abs*{f_{n_{k+1}}- f_{n_k}}$ konvergiert auf $E$ fast überall. Also konvergiert $\sum_{k=0}^{\infty} \abs*{f_{n_{k+1}}- f_{n_k}}$
auf $X$ fast überall. $\Rightarrow \sum_{k=0}^{\infty} f_{n_{k+1}}- f_{n_k}$ konvergiert auf $X$ fast überall. Also konvergiert auch $(f_{n_k})_{k \in \mathds{N}}$ auf
$X$ fast überall.\todo{vlt strukturiert übereinander?} Setze
\[
	f(x) := \begin{cases}
		\lim_{k\to \infty} f_{n_k}, &\text{ falls der Limes existiert}\\
		0, &\text{ sonst}
	\end{cases}
\]
$f$ ist messbar (Limes messbarer Funktionen). Weiter gilt $\abs*{f_{n_k}}^p \to \abs*{f}^p$ fast überall. 

Noch zu zeigen: $f \in L^p(\mu)$: Die Folge $\norm[p]{f_{n_k}}^p$ ist beschränkt, da $(f_n)_{n \in \mathds{N}}$ eine Cauchyfolge bezüglich $\norm[p]{\cdot}$ ist. Es gilt
\marginnote{$\abs*{f}^p \stackrel{\text{f.ü.}}{=} \liminf \abs*{f_{n_k}}^p$}
\[
	\int \abs*{f}^p \mathd \mu \le \liminf\limits_{k \to \infty} \int \abs*{f_{n_k}}^p \mathd \mu < \infty
\]
nach dem Lemma von Fatou. Also gilt $f \in L^p(\mu)$.
Noch zu zeigen: $f_n \xrightarrow{n \to \infty} f$ bezüglich $\norm[p]{\cdot }$. Sei dazu $j \in \mathds{N}$. Sei $n \ge n_j$. Für $i>j$ gilt
\[
	\norm[p]{f_n - f_{n_i}} \le \norm[p]{f_n - f_{n_j}} + \norm[p]{f_{n_j}- f_{n_i}} < 2 \cdot \frac{1}{2^j}    
\]
Weiter gilt $\abs*{f_{n}- f_{n_i}}^p \xrightarrow{i \to \infty} \abs*{f_n -f}^p$ fast überall. Mit dem Lemma von Fatou folgt 
\[
	\int\!\abs*{f_n- f}^p \mathd\mu \le \liminf \int \! \abs*{f_n - f_{n_i}}^p \mathd\mu  \le \enbrace*{2 \cdot \frac{1}{2^j}}^p
\]
$\Rightarrow \norm[p]{f_n -f} \le 2 \cdot \frac{1}{2^j}$, falls $n \ge n_j$, also $\norm[p]{f_n -f} \xrightarrow{n \to \infty} 0 $. \bewende
% subsection 411 (end)
% section 4 (end)
\newpage

\section{Schwache Topologien, Reflexivität} % (fold)
\label{sec:5}
\subsection[Definition: Schwache Topologie]{Definition} % (fold)
\label{sub:51}
Sei $X$ ein normierter $\mathds{K}$-Vektorraum. Die \Index{schwache Topologie} $\weakT{X}$ ist die gröbste\footnote{auch: "dollste"} Topologie auf $X$, sodass alle 
$\varphi \in X^*$ stetig sind. 
\todo{RevChap5}
% subsection 51 (end)

\subsection[Bemerkungen zu schwachen Topologie]{Bemerkung} % (fold)
\label{sub:52}
\begin{enumerate}[(i)]
	\item Jedes $\varphi \in X^*$ ist (nach Definition von $X^*=\mathcal{L}(X,\mathds{K})$) stetig bezüglich $\mathcal{T}^{\norm{\cdot}}_X$, d.h. 
	$\weakT{X} \subset \mathcal{T}_X^{\norm{\cdot}}$. Mit anderen Worten: Die Abbildung 
	$\id_X: \enbrace*{X,\mathcal{T}_X^{\norm{\cdot}}}\to\enbrace[\Big]{X,\weakT{X}}$ ist stetig.
	\item Mengen der Form $\varphi ^{-1}(U)$ für $U \subset K$ offen, $\varphi \in X^*$ bilden eine Subbasis für $\weakT{X}$. Mengen der Form 
	\[
		\bigcap_{i=1}^n \varphi_i ^{-1} (U_i) \quad, \qquad U_i \subset \mathds{K} \text{ offen}, \varphi_i \in X^*
	\]
	bilden eine Basis für $\weakT{X}$.
\end{enumerate}
% subsection 52 (end)

\subsection[Proposition: Die schwache Topologie ist Hausdorffsch]{Proposition} % (fold)
\label{sub:53}
$\enbrace[\big]{X,\weakT{X}}$ ist Hausdorffsch.
\minisec{Beweis}
Seien $x \not= y \in X$. Nach \hyperref[sub:227]{Hahn-Banach} existiert $\varphi \in X^*$ mit $\varphi(x) \not= \varphi(y)$. Setze 
$\varepsilon := \frac{\abs*{\varphi(x)- \varphi(y)} }{2}>0$, dann sind $\varphi ^{-1} \enbrace[\big]{B_\mathds{K} \enbrace*{\varphi(x),\varepsilon} } $ und
$\varphi ^{-1} \enbrace[\big]{B_\mathds{K} \enbrace*{\varphi(y), \varepsilon}}$ disjunkte offene Umgebungen von $x$ bzw. $y$. \bewende
% subsection 53 (end)

\subsection[Proposition: Äquivalenz zu schwacher Konvergenz]{Proposition} % (fold)
\label{sub:54}
Sei $X$ ein normierter Raum, $x \in X$, $(x_\lambda )_\Lambda \subset X$ ein Netz. Dann sind äquivalent:
\begin{enumerate}[(i)]
	\item $x_\lambda \xrightarrow{\w}  x$, d.h. $x_\lambda \to x$ in $\weakT{X}$
	\item $\varphi(x_\lambda ) \to \varphi(x)$ für alle $\varphi \in X^*$.
\end{enumerate}
\minisec{Beweis}
\begin{description}
	\item[(i)$\Rightarrow$(ii):] Sei $\varphi \in X^*$, $\varepsilon>0$. $\varphi(x) \in B_\mathds{K}\enbrace[\big]{\varphi(x),\varepsilon} \subset \mathds{K}$ offen.
	$\varphi$ stetig bezüglich $\weakT{X}$
	\[
		\Rightarrow  x \in \varphi ^{-1} \enbrace[\Big]{ B_\mathds{K}\enbrace[\big]{\varphi(x), \varepsilon} } \underset{\text{$\w$-offen}}{\subset} X 
	\] 
	$x_\lambda \xrightarrow{\w} x  \Rightarrow  \exists \overline{\lambda } \in \Lambda$ mit 
	$x_\lambda \in \varphi^{-1}\enbrace[\big]{B_\mathds{K}\enbrace*{\varphi(x), \varepsilon}}$, falls $\lambda \ge \overline{\lambda }$. Dann folgt
	$\varphi(x_\lambda ) \in B_\mathds{K}\enbrace*{\varphi(x),\varepsilon}$ falls $\lambda \ge \overline{\lambda}$, also $\varphi(x_\lambda ) \to \varphi(x)$ in $\mathds{K}$.
	\item[(ii)$\Rightarrow$(i):] Sei $x \in V \subset X$ $\w$-offen. Nach \ref{sub:52}(ii) existieren $\varphi_1, \ldots , \varphi_m \in X^*$, $U_1, \ldots , U_m$ offen in
	$\mathds{K}$ mit 
	\[
		x \in \phi_1 ^{-1}(U_1) \cap \ldots  \cap \varphi_m ^{-1}(U_m) \subset V
	\]
	Wegen (ii) existieren $\lambda_1, \ldots , \lambda _m \in \Lambda$ mit $\varphi_i(x_\lambda ) \in U_i$, falls $\lambda  \ge \lambda_i$, $i=1,\ldots,m$.
	Wähle $\overline{\lambda} \in \Lambda$ mit $\overline{\lambda} \ge \lambda_1, \ldots , \lambda_m$. Dann gilt $\varphi_i(x_\lambda) \in U_i$, falls 
	$\lambda \ge \overline{\lambda}$, $i=1,\ldots ,m$. Damit ist $x_\lambda \in \varphi_i ^{-1}(U_i)$, falls $\lambda \ge \overline{\lambda}$, $i=1,\ldots,m$. Also gilt
	$x_\lambda \in \bigcap_{i=1}^m \varphi_i ^{-1}(U_i) \subset V$, falls $\lambda \ge \overline{\lambda}$. \bewende
\end{description}
% subsection 54 (end)

\subsection{Corollar} % (fold)
\label{sub:55}
Wenn $x_\lambda \xrightarrow{\w} x$, dann gilt: $(x_\lambda)_\Lambda \subset X$ ist beschränkt.
\minisec{Beweis}
Prinzip der gleichmäßigen Beschränktheit (siehe \ref{sub:35}). \bewende
% subsection 55 (end)

\subsection[Beispiel: Konvergenz impliziert schwache Konvergenz, aber nicht umgekehrt]{Beispiel} % (fold)
\label{sub:56}
Im Allgemeinen gilt
\[
	x_\lambda \xrightarrow{\norm{\cdot}} x \quad \begin{matrix}
		\Rightarrow \\
		\nLeftarrow 
	\end{matrix} \quad  x_\lambda \xrightarrow{\w} x 
\]
Betrachte $X=\ell^2(\mathds{N})$ mit $\norm[2]{\cdot }$. 
Sei $e_n \in \ell^2(\mathds{N})$ gegeben durch $e_n(m) := \delta_{n,m}$\marginnote{Kronecker-Delta}. Dann gilt $e_n  \not\to 0$ für $n \to \infty$ in 
$\mathcal{T}^{\norm[2]{\cdot }}$, denn $\norm[2]{e_n -0} = \norm[2]{e_n} \equiv 1 \not\to 0$. Aber $e_n \xrightarrow{\w} 0$ für $n \to	\infty$: 

Wir werden später sehen, dass $\ell^2(\mathds{N})^* \cong \ell^2(\mathds{N})$ via $\sprod*{a}{b} = \sum \overline{a}(k) b(k)$.
Es gilt aber 
\[
	\sprod*{a}{e_n} = \sum \overline{a(k)} e_n(k) = \overline{a(n)} \xrightarrow{n \to \infty} 0.
\] 
$\Rightarrow e_n \xrightarrow{\w}0$ für $n \to \infty$. \smallskip\\
\emph{Direktes Argument: Übung.}
% subsection 56 (end)

\subsection{Proposition} % (fold)
\label{sub:57}
Seien $X,Y$ metrische Räume, $T \in \mathcal{L}(X,Y)$ ein beschränkter (d.h. normstetiger) linearer Operator. Dann ist $T$ auch stetig bezüglich der schwachen Topologie auf
$X$ und $Y$.
\minisec{Beweis}
Zu zeigen: Falls $U \subset Y$ schwach-offen ist, so ist $T ^{-1}(U)$ schwach-offen in $X$. Wir können annehmen, dass $U$ von der Form
$\varphi ^{-1} \enbrace*{B_\mathds{K}(\lambda,\varepsilon)}$, $\varphi \in Y^*, \lambda \in \mathds{K}, \varepsilon>0$ ist, da diese Mengen eine Subbasis bilden.
Es gilt nun
\begin{align*}
	T ^{-1} \enbrace[\Big]{ \varphi ^{-1} \enbrace[\big]{ B_\mathds{K}(\lambda,\varepsilon)} } = \enbrace[\Big]{\underbrace{\varphi \circ T}_{\in X^*}} ^{-1} 
	\enbrace[\big]{B_\mathds{K}(\lambda,\varepsilon)} \underset{\text{$\w$-offen}}{\subset} X   \bewende
\end{align*}
% subsection 57 (end)

\subsection[Definition: Die $\w^*$-Topologie auf $X^*$]{Definition} % (fold)
\label{sub:58}
Sei $X$ ein normierter Raum. Die \bet{$\mathbf{w}^{\mathbf{*}}$-Topologie} $\weakTstar{ }$\index{w-Topologie@$\w^*$-Topologie} ist die gröbste Topologie auf $X^*$, sodass alle $\iota_X(x) \in X^{**}$ stetig sind.
% subsection 58 (end)

\subsection[Bemerkungen zur $\w^*$-Topologie auf $X^*$]{Bemerkung} % (fold)
\label{sub:59}
\begin{enumerate}[(i)]
	\item Es gilt $\weakTstar{X^*} \subset \weakT{X^*} \subset \mathcal{T}^{\norm[X^*]{\cdot}}_{X^*}$
	\item Mengen der Form $W(\varphi,x \varepsilon) := \set[\psi \in X^*]{\abs*{\psi(x)- \varphi(x)}< \varepsilon}$ bilden eine Subbasis für $\weakTstar{X^*}$.
	Mengen der Form $\set[\psi]{ \psi(x) \in U}$, $U \subseteq \mathds{K}$ offen, bilden ebenfalls eine Subbasis für $\weakTstar{X^*}$.
	\item Ein Netz $(\varphi_\lambda)_{\lambda \in \Lambda} \subset X^*$ konvergiert $\w^*$ gegen $\varphi \in X^*$ genau dann, wenn 
	$\iota_X(x)(\varphi_\lambda)=\varphi_\lambda(x) \xrightarrow{\lambda} \varphi(x) = \iota_X(x)(\varphi)$ für jedes $x \in X$.
	\item $\weakTstar{X^*}$ ist Hausdorff. Beweis: Übung.
\end{enumerate}
% subsection 59 (end)

\subsection{Satz (Banach-Alaoglu)} % (fold)
\label{sub:510}
Sei $X$ ein normierter Raum. Dann ist die abgeschlossene Einheitskugel in $X^*$
\[
	\overline{B_{X^*}(0,1)} = \set[\varphi \in X^*]{\norm[X^*]{\varphi} \le 1}  
\]
$\w^*$-kompakt.
\minisec{Beweis}
Für $x \in X$ setze $F_x := \mathds{K}$ und $F := \prod_{x \in X} F_x$. Wir versehen $F$ mit der Produkttopologie $\mathcal{T}_F^\Pi$. Definiere eine Abbildung
$\kappa : X^* \to F$ durch $\kappa(\varphi)_x := \varphi(x)$, $x \in X$. Dann ist $\kappa$ injektiv (warum?) und $\weakTstar{ }$ stimmt überein mit der durch $\kappa$ und
$\mathcal{T}_F^\Pi$ induzierten Topologie auf $X^*$. Diese besteht aus den Mengen $\set[\kappa ^{-1}(U)]{U \in \mathcal{T}_F^\Pi}$. Eine Basis für $\mathcal{T}_F^\Pi$ sind
Mengen der Form 
\[
	\prod_{x \in X} U_x, \quad U_x \subset \mathds{K} \text{ offen} 
\]
mit $U_x = \mathds{K}$ für alle bis auf endlich viele $x \in X$. Eine Basis, für die durch
$\kappa$ und $\mathcal{T}_F^\Pi$ induzierte Topologie auf $X^*$ sind Mengen der Form
\[
	\kappa ^{-1} \enbrace*{ \prod_{x \in X} U_x} = \underbrace{\bigcap_{i=1}^n \set[\varphi]{\varphi(x_i) \in U_{x_i}} }_{
	\text{bilden Basis für $\weakTstar{X^*}$, \ref{sub:59}(ii)}}
\]
mit $U_x$ wie oben. Dabei erhalten wir die $x_i$ durch $\kappa ^{-1} \enbrace*{\prod_{x \in X} V_x} = \set[\varphi]{\varphi(x_i)\in U_{x_i}}$, wobei 
\[
	V_x = \begin{cases}
		\mathds{K}, &\text{ falls } x \not= x_i\\
		U_{x_i}, &\text{ falls } x=x_i
	\end{cases}
\]
Falls $\norm{\varphi} \le 1$, so gilt $\abs*{\varphi(x)} \le \norm[X]{x}$, also 
\[
	\kappa \enbrace*{\overline{B_{X^*}(0,1)}} \subset \underbrace{\prod_{x \in X} \overline{B_{\mathds{K}}(0, \norm{x})}}_{\text{kompakt nach Tychonov}} \subset F
\]
Bleibt zu zeigen: $\kappa \enbrace*{\overline{B_{X^*}(0,1)}} \subset F$ ist abgeschlossen. Sei $f \in \overline{\kappa \enbrace*{\overline{B_{X^*}(0,1)}}}$. Dies definiert
eine Abbildung $\varphi_f : X \to \mathds{K}$ durch $\varphi_f(x) := f_x$. Falls nun $\varphi_f$ linear und stetig ist mit $\norm{\varphi_f} \le 1$, so gilt
$\kappa(\varphi_f) = f$, denn $\kappa(\varphi_f)(x) = \varphi_f(x) = f_x$. Damit folgt $f \in \kappa \enbrace*{ \overline{B_{X^*}(0,1)}}$ und damit 
$\overline{B_{X^*}(0,1)}$ kompakt.
\begin{description}
	\item[Additivität:] Seien $y,z \in X$, $\varepsilon>0$. Dann ist 
	\[
		W_\varepsilon := \set[g \in \prod_{x \in X} \overline{B_\mathds{K}(0, \norm{x} )} ]{\begin{array}{r}
			\abs*{g_{y+z}- f_{y+z}} < \varepsilon \\ \enspace\abs*{g_y-f_y} < \varepsilon \\ \enspace\abs*{g_z -f_z}  < \varepsilon
		\end{array}  }  
		\underset{\text{offen}}{\subset} \prod_{x \in X} \overline{B_\mathds{K}(0,\norm{x})} 
	\]
	$f \in \overline{\kappa \enbrace*{\overline{B_{X^*}(0,1)}}} \Rightarrow \exists g \in W_\varepsilon \cap \kappa \enbrace*{\overline{B_{X^*}(0,1)}}$. Sei 
	$g' \in \overline{B_{X^*}(0,1)}$ mit $\kappa(g')=g$. $g' \in X^*$ ist linear, also gilt auch $g_{y+z}= g_y + g_z$ und daher 
	\[
		\abs*{f_{y+z} - (f_y + f_z)} < \varepsilon + \varepsilon + \varepsilon= 3 \varepsilon
	\]
	$\varepsilon>0$ war beliebig, also $f_{y+z}= f_y+f_z$ und $\varphi_f(y+z)= \varphi_f(y)+ \varphi_f(z)$. Also ist $\varphi$ additiv;
	\item[Skalarmultiplikation:] analog.
\end{description}
Es gilt $\abs*{\varphi_f(x)} = \abs*{f_x} \le \norm{x}$, $x \in X$, also $\norm[X^*]{\varphi_f} \le 1$. \bewende
% subsection 510 (end)

\subsection[Erinnerung: separabel, 1. abzählbar, 2. abzählbar]{Erinnerung} % (fold)
\label{sub:511}
\begin{enumerate}[(a)]
	\item Ein topologischer Raum heißt
	\begin{enumerate}[(i)]
		\item \Index{separabel}, falls $X$ eine abzählbare dichte Teilmenge enthält.
		\item \Index{1. abzählbar}, falls gilt: Jedes $x \in X$ besitzt abzählbar viele Umgebungen $U_i$, $i \in \mathds{N}$, sodass jede Umgebung $V$ von $x$ wenigstens ein
		$U_i$ enthält.  
		\item \Index{2. abzählbar}, falls gilt: $\mathcal{T}_X$ besitzt eine abzählbare Basis, d.h. es gibt offene Mengen $W_i$, $i \in N$, sodass für jedes $V \subset X$ 
		offen gilt: $V= \bigcup_{W_i \subset V} W_i$.
	\end{enumerate}
	\item 2. abzählbar $\Rightarrow $ 1. abzählbar und separabel.
	\item Ist $X$ kompakt und 1. abzählbar, so besitzt jede Folge in $X$ eine konvergente Teilfolge.
	
	Beweisskizze: $(x_n)_{\mathds{N}} \subset X$, setze $A_n := \overline{\set{x_n, x_{n+1}, \ldots }}$. Dann ist $\bigcap_\mathds{N} A_n$ nichtleer wegen Kompaktheit, also
	hat $(x_n)_\mathds{N}$ einen Häufungspunkt. Benutze nun 1. abzählbar.
	\item Sei $X$ kompakt und Hausdorff. Dann ist $X$ metrisierbar $\iff$ $X$ ist 2. abzählbar.
	
	"$\Longrightarrow$" ist eine einfache Übung. "$\Longleftarrow$" Urysohn, Metrisierungssätze. Den Satz von Urysohn hatten wir auch in "Grundlagen der Analysis, Topologie 
	und Geometrie" im letzen Semester bewiesen.
\end{enumerate}
% subsection 511 (end)

\subsection[Satz: Metrisierbarkeit von $\overline{B_{X^*}(0,1)}$ bezüglich der $\w^*$-Topologie]{Satz} % (fold)
\label{sub:512}
Sei $X$ ein separabler normierter Raum. Dann ist $\overline{B_{X^*}(0,1)}$ bezüglich $\weakTstar{X^*}$ metrisierbar. Insbesondere ist $\overline{B_{X^*}(0,1)}$ bezüglich
$\weakTstar{X^*}$ 2. abzählbar und kompakt, und jede Folge besitzt eine konvergente Teilfolge.
\minisec{Beweis}
Sei $\set[x_n]{n \in \mathds{N}} \subset B_X(0,1)$ dicht. Definiere eine Abbildung $d: X^* \times X^* \to [0,\infty)$ durch
\[
	d(\varphi, \psi) := \sum_{n \in \mathds{N}} \frac{1}{2^n} \cdot \abs[\big]{\varphi(x_n) - \psi(x_n)}
\]
Dann ist $d$ eine Metrik: $d$ ist offensichtlich symmetrisch und erfüllt die Dreiecksungleichung. Zu zeigen: $d$ ist definit:
\(
	\varphi \not= \psi \Rightarrow \exists x \in B_X(0,1) \text{ mit } \varphi(x) \not= \psi(x).
\)
Also existiert $x_n$ mit $\varphi(x_n) \not= \psi(x_n)$. Dann ist $d(\varphi,\psi) \not= 0$.

Sei nun $(\varphi_\lambda)_{\lambda \in \Lambda} \subset \overline{B_{X^*}(0,1)}$ ein Netz und $\varphi \in \overline{B_{X^*}(0,1)}$. Es genügt zu zeigen
\[
	\varphi_\lambda \xrightarrow{\w^*} \varphi \enspace \iff \enspace  d(\varphi_\lambda, \varphi) \longrightarrow 0
\]
denn dann stimmen $\weakTstar{X^*}$ und $\mathcal{T}_{X^*}^d$ überein. (Warum?)
\begin{description}
	\item["$\Rightarrow $":] Es gelte $\varphi_\lambda \xrightarrow{\w^*} \varphi$, also $\varphi_\lambda(x_n) \to \varphi(x_n)$ für $n \in \mathds{N}$ (nach 
	\ref{sub:59}(iii)). Sei $\varepsilon>0$. Wähle $n_0 \in \mathds{N}$ mit $\sum_{n \ge n_0} \frac{1}{2^n} < \frac{\varepsilon}{4}$. Wähle $\overline{\lambda} \in \Lambda$
	mit $\abs*{\varphi_\lambda(x_n)- \varphi(x_n)} < \frac{\varepsilon}{4}$ falls $n < n_0$, $\lambda \ge \overline{\lambda}$. Dann gilt für $\lambda \ge\overline{\lambda}$
	\begin{align*}
		d(\varphi_\lambda,\varphi) &\le \sum_{n < n_0} \frac{1}{2^n} \cdot \abs[\big]{\varphi_\lambda(x_n)- \varphi(x_n)} + \sum_{n \ge n_0} \frac{1}{2^n} \cdot 
		\underbrace{\abs[\big]{\varphi_\lambda(x_n) - \varphi(x_n)}}_{\le 2} \\
		&< 2 \cdot \frac{\varepsilon}{4} + 2 \cdot \frac{\varepsilon}{4}  = \varepsilon
	\end{align*}
	\item["$\Leftarrow$":] Übung mit: $\forall n . \varphi_\lambda(x_n) \to \varphi(x_n)$ und $ \set{x_n}$ dicht $\Rightarrow \varphi_\lambda(x) \to \varphi(x)$ für jedes
	$x \in X$. \bewende
\end{description}
% subsection 512 (end)

\subsection*{Nachtrag (zum Beweis von \ref{sub:512})} % (fold)
\label{sub:nachtrag_512}
Topologien sind durch Konvergenz von Netzen bestimmt: Sei $(X,\mathcal{T})$ ein topologischer Raum. 
\begin{align*}
	U \subset X \text{ offen } &\iff \forall x \in U : \exists x \in V \subset U, V \in \mathcal{T} \\
	U \subset X \text{ nicht offen } &\iff \exists x \in U : \text{ für jedes } x \in V \underset{\mathclap{\text{offen}}}{\subset} X \text{ gilt } V \cap (X \setminus U) 
	\not= \emptyset \\
	&\stackrel{\text{(\#)}}{\iff} \exists x \in U, \exists \text{ Netz } (x_\lambda)_{\lambda \in \Lambda} \subset X \setminus U \text{ mit } x_\lambda \to x
\end{align*}
$\stackrel{\text{(\#)}}{\Longleftarrow}$ ist trivial. $\stackrel{\text{(\#)}}{\Longrightarrow}$: $\Lambda := \set[V]{x \in V \subset X \text{ offen}}$ ist 
gerichtet bezüglich "$\subset$". Zu jedem $V \in \Lambda$ wähle \marginnote{Auswahlaxiom nötig!} $x_V \in (X \setminus U) \cap V$. Dann gilt 
$(x_v)_{v \in \Lambda} \subset X \setminus U$ und $x_v \to x$.

% subsection nachtrag_512 (end)

\subsection[Definition: Reflexiver normierter Raum]{Definition} % (fold)
\label{sub:513}
Der normierte Raum $X$ heißt \Index{reflexiv}, falls die kanonische Abbildung $\iota_X : X \to X^{**}$ ein Isomorphismus ist. 
($\iota_X$ ist Isomorphismus genau dann, wenn $\iota_X$ surjektiv ist.)
% subsection 513 (end)

\subsection[Proposition: Abgeschlossene Unterräume von reflexiven Banachräumen sind reflexiv]{Proposition} % (fold)
\label{sub:514}
Sei $X$ ein reflexiver Banachraum und $Y \subset X$ ein abgeschlossener Teilraum. Dann ist $Y$ reflexiv.
\minisec{Beweis}
Sei $\kappa : Y \hookrightarrow X$ die Inklusion. Dann ist $\kappa \in \mathcal{L}(Y,X)$ und nach \ref{sub:34} kommutiert das Diagramm
\[
	\begin{tikzcd}
		Y \rar[hook, "\iota_Y"] \dar["\kappa"', hook]& Y^{**}  \dar["\kappa^{\tr \,\tr}"] \\
		X \rar["\iota_X","\cong"', hook] & X^{**}
	\end{tikzcd}
\]
Sei $f\in Y^{**}$, dann existiert $x \in X$ mit $\iota_X(x)=\kappa^{\tr\,\tr}(f)$. Für $\varphi \in X^*$ gilt 
\[
	\varphi(x)=\iota_X(x)(\varphi)=\kappa^{\tr\,\tr}(f)(\varphi).
\]
Aber $(\kappa^\tr)^\tr (f) = f \circ \kappa^\tr$ und $\kappa^\tr(\varphi) = \varphi \circ \kappa = \varphi|_{Y}$, also $\varphi(x)= f \enbrace*{\varphi|_Y}$ für jedes 
$\varphi \in X^*$. Falls $x \not\in Y$, so existiert nach Hahn-Banach ein $\psi \in X^*$ mit $\psi|_{Y}=0$, $\psi(x)\not= 0$. Damit folgt 
\[
	0 \not= \varphi(x) = f \enbrace*{\psi|_{Y}}=0 \enspace\light
\]
also gilt $x \in Y$. Weiter gilt $\iota_Y(x)=f$, denn $\kappa^{\tr\,\tr}$ ist injektiv. Andernfalls existiert $0 \not= g \in Y^{**}$ mit $\kappa^{\tr\,\tr}(g)=0$, also
existiert $\sigma \in Y^*$ mit $g(\sigma) \not= 0$. Nach Hahn-Banach existiert $\rho \in X^*$ mit $\rho|_{Y}=\sigma$. Nun gilt 
\[
	0 = \kappa^{\tr\,\tr}(g) = g \enbrace*{\rho|_Y} = \rho(g) \not= 0 \enspace \light 
\]
$\Rightarrow \kappa^{\tr\,\tr}$ ist injektiv. \bewende
% subsection 514 (end)

\subsection[Proposition: Für Banachräume gilt: $X$ reflexiv $\iff$ $X^*$ reflexiv]{Proposition} % (fold)
\label{sub:515}
Für einen Banachraum $X$ gilt: $X$ reflexiv $\iff$ $X^*$ reflexiv.
\minisec{Beweis}
\begin{description}
	\item["$\Rightarrow$":] $\iota_X : X \xrightarrow{\cong} X^{**}$ ist ein isometrischer Isomorphismus. Nach Blatt 6, Aufgabe 2 ist $\iota_X^{\tr} : X^{***} \to X^*$
	auch ein isometrischer Isomorphismus. Nach \ref{sub:34} gilt  $\iota_X^{\tr} \circ \iota_{X^*}= \id_{X^{*}}$, also ist $\enbrace*{\iota_X^\tr}^{-1} = \iota_{X^*}$
	ebenfalls ein isometrischer Isomorphismus.
	\item["$\Leftarrow$":] genauso \checkmark \bewende
\end{description}
% subsection 515 (end)

\subsection[Satz: Einheitskugel eines Banachraumes $X$ liegt $\w^*$-dicht in Einheitskugel von $X^{**}$]{Satz} % (fold)
\label{sub:516}
Sei $X$ ein Banachraum. Dann gilt 
\[
	\overline{B_X(0,1)} \underset{\text{$\w^*$-dicht}}{\subset} \overline{B_{X^{**}}(0,1)}  
\]
\minisec{Beweis}
Zu zeigen: Zu $f \in \overline{B_{X^{**}}(0,1)}$ und $\varepsilon>0$, $\varphi_1, \ldots , \varphi_n \in X^*$, existiert $x \in \overline{B_X(0,1)}$ mit
$\abs*{f(\varphi_i)- \varphi_i(x)} < \varepsilon$ für $i=1,\ldots ,n$. Behauptung: Es gilt (gegeben $f,\varphi_1, \ldots , \varphi_n$)
\[
	\inf_{x \in \overline{B_X(0,1)} } \sum_{i=1}^{n} \abs[\big]{f(\varphi_i)- \varphi_i(x)}^2 = 0
\]
Sei $h(x) := \sum_{i=1}^{n} \abs*{f(\varphi_i)- \varphi_i(x)}^2 $ und $\inf_{x\in \overline{B_X(0,1)}} h(x) =: \mu \ge 0$. Wähle 
$(x_k)_{k \in \mathds{N}} \subset \overline{B_X(0,1)}$ mit $\lim_{k \to \infty} h(x_k)= \mu$. Nach Übergang zu einer Teilfolge dürfen wir annehmen, dass 
$\enbrace[\big]{\varphi_i(x_k)}_{k \in \mathds{N}}$ gegen  $\gamma_i \in \mathds{K}$ konvergiert für $i=1,\ldots ,n$. Setze $\delta_i := f(\varphi_i) -\gamma_i$, dann
gilt $\mu = \sum_{i=1}^{n} \abs*{\delta_i}^2$. Für $y \in \overline{B_X(0,1)}$, $t \in [0,1]$ gilt 
\begin{align*}
	\mu \le h\enbrace[\Big]{\underbrace{(1-t) \cdot x_k + t \cdot y}_{\in \overline{B_X(0,1)}}} &= 
	\sum_{i=1}^{n} \abs[\Big]{f(\varphi_i)- (1-t) \cdot \varphi_i(x_k) - t \cdot \varphi_i(y)}^2 \\
	&= \sum_{i=1}^{n} \abs[\Big]{f(\varphi_i)- \varphi(x_k) - t \cdot  \enbrace[\big]{\varphi_i(y)- \varphi_i(x_k)} }^2 \\
	&= \sum_{i=1}^{n} \abs[\big]{f(\varphi_i)- \varphi_i(x_k)}^2 + t^2 \cdot \sum_{i=1}^{n} \abs[\big]{\varphi_i(y)- \varphi_i(x_k)}^2 \\
	& \quad - 2 t \cdot \re \sum_{i=1}^{n} \enbrace[\big]{\varphi_i(y)- \varphi_i(x_k)} \cdot \overline{\enbrace[\big]{f(\varphi_i)- \varphi_i(x_k)} } \\
	\xrightarrow{k \to \infty} &\underbrace{\sum_{i=1}^{n} \abs[\big]{f(\varphi_i)- \gamma_i}^2}_{=\mu} + t^2 \cdot \sum_{i=1}^{n} \abs[\big]{\varphi_i(y)- \gamma_i}^2  
	-2 t \cdot \re \sum_{i=1}^{n} \enbrace[\big]{\varphi_i(y)- \gamma_i} \cdot \overline{\delta_i }
\end{align*}
Nach Abziehen von $\mu$ folgt für $t >0$
\begin{align*}
	&0 \le t \cdot \sum_{i=1}^{n} \abs[\big]{\varphi_i(y)- \gamma_i}^2 - 2 \cdot \re \sum_{i=1}^{n} \enbrace[\big]{\varphi_i(y)- \gamma_i} \cdot \overline{\delta_i}  \\
	\xRightarrow{t \searrow 0} \enspace & 2 \cdot \re \sum_{i=1}^{n} \enbrace[\big]{\varphi_i(y)- \gamma_i} \cdot \overline{\delta_i} \le 0 \text{ für alle } y \in \overline{B_X(0,1) }
\end{align*}
Für $y \in \overline{B_X(0,1)}$ gilt $\re \sum_{i=1}^{n} \enbrace*{\varphi_i(y)- \gamma_i} \cdot \overline{\delta _i} \le 0$.
Definiere $\varphi := \sum_{i=1}^{n} \overline{\delta _i} \cdot \varphi_i \in X^*$, dann gilt für $y \in \overline{B_X(0,1)}$
\[
	\re \varphi(y) \le \re \sum_{i=1}^{n} \gamma_i \overline{\delta_i} 
\]
also $\norm[X^*]{\varphi} \le \re \sum_{i=1}^{n} \gamma_i \cdot \overline{\delta_i}$. Aber 
\[
	\varphi(x_k) = \sum_{i=1}^{n} \overline{\delta_i} \cdot \varphi_i(x_k) \xrightarrow{k \to \infty} \sum_{i=1}^{n} \overline{\delta_i} \cdot \gamma_i, 
\]
das heißt $\abs*{\sum_{i=1}^{n} \overline{\delta_i} \cdot \gamma_i} \le \norm{\varphi} \le \re \sum_{i=1}^{n} \gamma_i \cdot \overline{\delta_i}$ und somit
$\sum_{i=1}^{n} \overline{\delta_i} \cdot \gamma_i = \norm[X^*]{\varphi}$. Wir erhalten 
\begin{align*}
	\mu = \sum_{i=1}^{n} \overline{\delta_i} \delta_i = \sum_{i=1}^{n}\overline{\delta_i} \enbrace[\big]{f(\varphi_i)- \gamma_i} = f(\varphi)-\norm[X^*]{\varphi} \le 0 \bewende
\end{align*}
% subsection 516 (end)

\subsection{Corollar} % (fold)
\label{sub:517}
Sei $X$ ein Banachraum, $f \in \overline{B_{X^{**}}(0,1)}$, $\varphi_1, \ldots , \varphi_n \in X^*$. Falls $\overline{B_X(0,1)}$ $\w$-kompakt ist, so existiert 
$x \in \overline{B-X(0,1)}$ mit $f(\varphi_i)= \varphi_i(x)$, $i=1,\ldots ,n$.
\minisec{Beweis}
Definiere $h : X \to \mathds{R}$ durch $h(x):= \sum_{i=1}^{n} \abs*{f(\varphi_i)- \varphi_i(x)}^2$ ist stetig auf $X$ bezüglich $\weakT{X}$, nimmt also auf
$\overline{B_X(0,1)}$ ihr Minimum an. \bewende
% subsection 517 (end)


\subsection{Satz} % (fold)
\label{sub:518}
Sei $X$ ein Banachraum. Dann sind äquivalent: 
\begin{enumerate}[(i)]
	\item $X$ ist reflexiv.
	\item $\overline{B_X(0,1)}$ ist $\w^*$-kompakt.
	\item $\weakT{X^*}= \weakTstar{X^*}$.
\end{enumerate}
\minisec{Beweis}
\begin{description}
	\item[(i)$\Rightarrow$(ii):] $\overline{B}_{X^{**}}(0,1)$ ist $\w^*$-kompakt nach \ref{sub:510}. Nach \ref{sub:515} ist $X^*$ reflexiv, daher stimmen auf $X^{**}$
	die $\w^*$- und die $\w$-Topologie überein. $\Rightarrow \overline{B_{X^{**}}(0,1)}$ ist $\w$-kompakt. Aber unter $X \cong X^{**}$ stimmen die jeweiligen $\w$-Topologien
	überein; ebenso $\overline{B_{X}(0,1)} \cong \overline{B_{X^{**}}(0,1)}$.
	\item[(ii)$\Rightarrow$(i):] Es gilt $(X,\weakT{X}) \subset (X^{**}, \weakTstar{X^{**}})$ in der Relativtopologie (warum?).
	$\overline{B_X(0,1)}$ kompakt $\Rightarrow \overline{B_X(0,1)} \subset \overline{B_{X^{**}}(0,1)}$ $\w^*$-abgeschlossen und $\w^*$-dicht. Also folgt
	$\overline{B_X(0,1)} = \overline{B_{X^{**}}(0,1)}$ und somit $X=X^{**}$.
	\item[(iii)$\Rightarrow$(i):] $\overline{B_{X^*}(0,1)}$ ist $\w^*$-kompakt nach Banach-Alaoglu (\ref{sub:510}). Mit (iii) folgt, dass $\overline{B_{X^*}(0,1)}$
	$\w$-kompakt ist. Da (ii)$\Rightarrow$(i) folgt, dass $X^*$ reflexiv ist und somit ist nach \ref{sub:516} auch $X$ reflexiv.
	\item[(i)$\Rightarrow$(iii):] klar. \bewende
\end{description}
% subsection 518 (end)

\subsection{Bemerkung} % (fold)
\label{sub:519}
\begin{enumerate}[a)]
	\item Sei $X$ ein normierter Raum. Falls $X^*$ separabel (bezüglich $\mathcal{T}_{X^*}^{\norm{\cdot}}$) ist, so ist $X$ separabel (bezüglich 
	$\mathcal{T}_X^{\norm{\cdot}}$)
	\item Für einen Banachraum $X$ gilt: $X$ reflexiv $\iff$ $\overline{B}_X(0,1)$ ist folgenkompakt bezüglich $\weakT{X}$.
\end{enumerate}
\minisec{Beweis}
\begin{enumerate}[a)]
	\item Übung.
	\item \begin{description}
		\item["$\Rightarrow$":] leicht.
		\item["$\Leftarrow$":] nicht so leicht. 
	\end{description}
\end{enumerate}
% subsection 519 (end)
% section 5 (end)
\newpage

\section{Gleichmäßig konvexe Räume, noch einmal $L^p$} % (fold)
\label{sec:6}

\subsection[Definition und Proposition: Gleichmäßig konvexer Raum]{Definition und Proposition} % (fold)
\label{sub:61}
Ein normierter Raum $X$ heißt \Index{gleichmäßig konvex}, falls eine der folgenden äquivalenten Bedingungen erfüllt ist:\todo{RevChap6}
\begin{enumerate}[(i)]
	\item Für jedes $\varepsilon>0$ existiert $\delta >0$, sodass für $x,y \in X$ mit $\norm{x}, \norm{y} \le 1$ gilt:
	\[
		\norm{x-y} \ge \varepsilon \enspace\Longrightarrow \enspace\norm{\frac{1}{2} (x+y) } \le 1 -\delta 
	\]
	\item \label{61:enum:2} Für jedes $\varepsilon>0$ existiert $\delta >0$, sodass für $x,y \in X$ mit $\norm{x}, \norm{y}= 1$ gilt:
	\[
		\norm{x-y} \ge \varepsilon \enspace\Longrightarrow \enspace\norm{\frac{1}{2} (x+y)} \le 1-\delta   
	\]
	\item Für $(x_n)_\mathds{N}, (y_n)_\mathds{N} \subset X$ mit $\norm{x_n} = \norm{y_n}=1$ für alle $n \in \mathds{N}$ gilt: 
	\[
		\lim_{ n \to \infty} \norm{\frac{1}{2} (x_n+y_n)}=1 \enspace\Longrightarrow \enspace\lim_{ n \to \infty} \norm{x_n- y_n}=0  
	\]
	\item Für $(x_n)_\mathds{N}, (y_n)_\mathds{N} \subset X$ mit $\limsup_n \norm{x_n} = \limsup_n \norm{y_n}=1$ gilt:
	\[
		\lim_{ n \to \infty} \norm{\frac{1}{2} (x_n+y_n)}=1 \enspace\Longrightarrow \enspace\lim_{ n \to \infty} \norm{x_n- y_n}=0  
	\]
\end{enumerate}
\minisec{Beweis}
\begin{description}
	\item[(i)$\Rightarrow$(ii)$\Rightarrow$(iii):] trivial.
	\item[(iii)$\Rightarrow$(iv):] einfach. Zeige zunächst: Aus $\limsup \norm{x_n}, \limsup \norm{y_n}\le 1$ und $\lim_{ n \to \infty} \norm{\frac{1}{2} (x_n+y_n)}=1$
	folgt $\lim_{ n \to \infty} \norm{x_n}=\lim_{ n \to \infty} \norm{y_n}=1$.
	\item[(iv)$\Rightarrow$(i):] Angenommen (i) gilt nicht, d.h. $\exists \varepsilon>0$ und für jedes $\delta_n=\frac{1}{n+1}$ existieren $x_n,y_n \in X$ mit 
	$\norm{x_n}, \norm{y_n}\le 1$, $\norm{x_n-y_n}\ge \varepsilon$ und $\norm{\frac{1}{2}(x_n+y_n)} > 1- \frac{1}{n+1}$. \light zu (iv). \bewende
\end{description}
% subsection 61 (end)

\subsection[Beispiel für gleichmäßig konvexe Räume]{Beispiel} % (fold)
\label{sub:62}
$(\mathds{R}^n, \norm[p]{\cdot})$ ist gleichmäßig konvex für $1 < p < \infty$ aber nicht für $p=1, p=\infty$
\begin{figure}[H]
	\centering{
	\begin{tikzpicture}[scale=0.8]
		\coordinate (x1) at (5:2);
		\coordinate (x2) at (80:2);
		\draw[lightgray] (-2.4,0) -- (2.4,0);
		\draw[lightgray] (0,-2.4) -- (0,2.4);
		\draw[very thick, Chartreuse3] (0,0) circle(2);
		\draw[thick, dashed] (x1) -- (x2) node[midway, name=mid]{ };
		\draw[fill=DodgerBlue2] (x1) circle(0.05) node[right]{$x$};
		\draw[fill=DodgerBlue2] (x2) circle(0.05) node[above]{$y$};
		\draw[fill=DodgerBlue2] (mid) circle(0.05) node[below left]{$\frac{1}{2}(x+y)$};
	\end{tikzpicture}
	\qquad 
	\begin{tikzpicture}[scale=0.8]
		\coordinate (x1) at (1.8,0.2);
		\coordinate (x2) at (0.4,1.6);
		\draw[lightgray] (-2.4,0) -- (2.4,0);
		\draw[lightgray] (0,-2.4) -- (0,2.4);
		\draw[very thick,DarkGoldenrod2] (2,0) -- (0,2) -- (-2,0) -- (0,-2) -- cycle;
		\draw[thick, dashed] (x1) -- (x2) node[midway, name=mid]{ };
		\draw[fill=DodgerBlue2] (x1) circle(0.05) node[above right]{$x$};
		\draw[fill=DodgerBlue2] (x2) circle(0.05) node[above right]{$y$};
		\draw[fill=DodgerBlue2] (mid) circle(0.05) node[above right]{$\frac{1}{2}(x+y)$};
	\end{tikzpicture}
	\qquad 
	\begin{tikzpicture}[scale=0.8]
		\coordinate (x1) at (2,-1.1);
		\coordinate (x2) at (2,1.8);
		\draw[lightgray] (-2.4,0) -- (2.4,0);
		\draw[lightgray] (0,-2.4) -- (0,2.4);
		\draw[very thick,Firebrick3] (-2,-2) rectangle (2,2) ;
		\draw[thick, dashed] (x1) -- (x2) node[midway, name=mid]{ };
		\draw[fill=DodgerBlue2] (x1) circle(0.05) node[right]{$x$};
		\draw[fill=DodgerBlue2] (x2) circle(0.05) node[right]{$y$};
		\draw[fill=DodgerBlue2] (mid) circle(0.05) node[right]{$\frac{1}{2}(x+y)$};
	\end{tikzpicture}
	\caption[Einheitskugel für $p=2$, ${\color{DarkGoldenrod2}p=1}$ und ${\color{Firebrick3}p=\infty}$]{Einheitskugel für $p=2$, ${\color{DarkGoldenrod2}p=1}$ und 
	${\color{Firebrick3}p=\infty}$. Nur für $p \not\in \set{1,\infty}$ ist (\ref{61:enum:2}) aus 
	\ref{sub:61} erfüllt. Anschaulich gesprochen bedeutet dies, dass nur in gleichmäßig konvexen Räumen Bälle auch tatsächlich \enquote{rund} sind.}}
\end{figure}
% subsection 62 (end)

\subsection[Definition und Proposition: Strikt konvexer Raum]{Definition und Proposition} % (fold)
\label{sub:63}
Ein gleichmäßig konvexer normierter $\mathds{R}$-Vektorraum ist \Index{strikt konvex}, d.h. für alle $0\not= x,y \in X$ mit $\norm{x+y}=\norm{x} + \norm{y}$ existiert 
$\lambda >0$ mit $x=\lambda \cdot y$. 
\minisec{Beweis}
Seien $x,y$ wie angegeben; \oE $\norm{x}=1$. Setze $x' := \frac{x+y}{\norm{x+y}}$, $a := \frac{x+x'}{2}$, $\alpha := \norm{a}$, $\beta := \norm{x+y-a}$
\begin{figure}{H}
	\centering{
	\begin{tikzpicture}[scale=3, my brace/.style={decorate, decoration={brace, amplitude=5pt, mirror, raise=3pt}, gray}]
		\coordinate (0) at (0,0);
		\coordinate (a) at (2,0.5);
		\coordinate (x) at (2,1);
		\coordinate (x') at (2,0);
		\coordinate (x+y) at (2.8,0);
		\draw[thick] (0,0)  -- (x+y);
		\draw[my brace] (0) -- (x') node[midway, yshift=-0.6cm]{$\norm{x'}=1$};
		\draw[my brace] (x') --(x+y) node[midway, yshift=-0.6cm]{$\norm{y}$};
		
		\draw[thick] (0) -- (x);
		\draw[my brace] (x) -- (0) node[midway, above left, yshift=0.16cm]{$\norm{x}=1$};
		
		\draw[thick] (x) -- (x');
		
		\draw[thick] (x+y) -- (x);
		\draw[my brace] (x+y) -- (x) node[midway, above right, xshift=0.1cm]{$\norm{y}$};
		
		\draw[thick,dashed] (0) -- (a) -- (x+y);
		\draw[my brace] (0) -- (a) node[midway, below,yshift=-0.2cm,xshift=0.17cm]{$\alpha$};
		\draw[my brace] (a) -- (x+y) node[midway, below,yshift=-0.2cm,xshift=-0.22cm]{$\beta$};
		
		\foreach \t/\pos in {0/left,x/above,x'/below,x+y/right,a/above right}{
			\draw[fill=DodgerBlue2,thick] (\t) circle (0.02) node[\pos]{$\t$};
		}
	\end{tikzpicture}
	\caption[Zeichnung zum Beweis von \ref{sub:63}]{Die Tatsache, dass $\norm{x}$ und $\norm{x'}$ in dieser Zeichnung nicht gleich sind, weißt schon darauf hin, dass diese 
	Zeichnung in Wahrheit \enquote{degeneriert}.}}
\end{figure}
Es gilt 
\[
	\norm{x+y-x'} = \enbrace*{1- \frac{1}{\norm{x+y}}} \cdot \norm{x+y} = \norm{x+y}-1 = 1 + \norm{y} -1 = \norm{y}      
\]
also $\beta = \norm{x+y-a} \le \doppelstrich[\big]{\frac{x+y-x}{2}} + \doppelstrich[\big]{\frac{x+y-x'}{2}}\le \frac{\norm{y}}{2} + \frac{\norm{y}}{2} = \norm{y}$. Weiter gilt
$\alpha = \norm{a} \le \frac{\doppelstrich{x}}{2} + \frac{\doppelstrich{x'}}{2} =1$ und
\begin{align*}
	1 + \norm{y} = \norm{x+y} \le \norm{a} + \norm{x+y-a} = \alpha+ \beta\le 1 + \norm{y} 
\end{align*}
also $1 \le \alpha \le 1$. Wegen $\norm{x}=1$, $\norm{x'}=1$ und $\norm{\frac{1}{2} (x+x')}= \norm{a}=\alpha=1 > 1 -\delta$ für jedes $\delta >0$ erhalten wir aus 
gleichmäßiger Konvexität $\norm{x-x'} \le \varepsilon$ für jedes $\varepsilon>0 \Rightarrow x=x'$ und $y=\norm{y} \cdot x$. \bewende
% subsection 63 (end)

\subsection{Proposition} % (fold)
\label{sub:64}
Sei $X$ ein normierter $\mathds{R}$-Vektorraum, $a \in X$, $W \subset X$ konvex und abgeschlossen. 
\begin{enumerate}[(i)]
	\item Ist $W$ strikt konvex, so existiert höchstens ein $x \in W$ mit $\norm{a -x}= \inf_{y \in W} \norm{a-y}$.
	\item Ist $X$ vollständig und gleichmäßig konvex, so existiert genau ein $x \in W$ mit $\norm{a-x}= \inf_{y \in W} \norm{a-y}$
\end{enumerate}
\minisec{Beweis}
Wir dürfen ohne Einschränkungen $a=0$ annehmen. Für $a \in W$ ist nichts zu beweisen; für $a \not\in W$ gilt $\inf_{y \in W} \norm{a-y} >0$ (denn $W$ ist abgeschlossen).
Ohne Einschränkungen dürfen wir außerdem $\inf_{y \in W} \norm{a-y}=1$ annehmen. 
\begin{enumerate}[(i)]
	\item Sei $x,x' \in W$ mit $\norm{a-y}= \norm{a-x'}= \inf_{y \in W} \norm{a-y}$, also $\norm{x}=\norm{x'}=1$. Dann gilt 
	\[
		1= \inf_{y \in W} \norm{a-y} \le \doppelstrich[\bigg]{\underbrace{\frac{1}{2} (x+x')}_{\in W}} \le \frac{1}{2} \norm{x} + \frac{1}{2} \norm{x'} =1   
	\]
	Es folgt aus strikter Konvexität, dass $x=\lambda \cdot x'$ für ein $\lambda \ge 0$. Wegen $\norm{x}=\norm{x'}=1$ folgt $\lambda=1$ und $x=x'$. 
	\item Wähle eine Folge $(x_n)_{\mathds{N}} \subset W$ mit \marginnote{\oE wieder $a=0$}
	\[
		\lim_{ n \to \infty} \norm{x_n} = \lim_{ n \to \infty} \norm{a-x_n} = \inf_{y \in W} \norm{a-y} = 1 
	\]
	Behauptung: Diese Folge ist Cauchy. $W$ ist konvex, also folgt $\frac{1}{2} (x_n+x_m) \in W$. Es gilt dann 
	\begin{equation*}
		1 \le \norm{\frac{1}{2}(x_n+x_m)} \le \frac{1}{2} \enbrace[\big]{\norm{x_n}+ \norm{x_m}} \enspace\text{ für } n,m \in \mathds{N} \label{eq:64} \tag{\#}
	\end{equation*}
	Angenommen $(x_n)_\mathds{N}$ sei nicht Cauchy. Dann existiert $\varepsilon>0$, so dass gilt: Für jedes $N \in \mathds{N}$ existieren $n,m >N$ mit 
	$\norm{x_n +x_m} >\varepsilon$. Wir können daher Teilfolgen $(x_{n_k})_{k \in \mathds{N}}$, $(x_{m_k})_{k \in \mathds{N}}$ bilden mit 
	$\liminf_k \norm{x_{n_k} - x_{m_k}} \ge \varepsilon$. Es gilt aber $\lim_{ k \to \infty} \norm{x_{n_k}} = \lim_{ k \to \infty} \norm{x_{n_k}}=1$ und wegen \eqref{eq:64}
	gilt $\lim_{ k \to \infty} \norm{\frac{1}{2} (x_{n_k} + x_{m_k})} = 1$ \light zu \ref{sub:61} (iv). \bewende
\end{enumerate}
% subsection 64 (end)

\subsection{Bemerkung} % (fold)
\label{sub:65}
Der Beweis von (ii) zeigt auch: Sei $X$ gleichmäßig konvex und $(x_n)_\mathds{N} \subset X$ eine Folge mit $\limsup_n \norm{x_n}\le 1$ und 
$\lim_{ n,m \to \infty} \norm{\frac{1}{2} (x_n+x_m)}=1$. Dann gilt $\lim_{ n \to \infty} \norm{x_n}=1$ und $(x_n)_{\mathds{N}}$ ist Cauchy.
% subsection 65 (end)

\subsection{Proposition: Jensensche Ungleichung} % (fold)
\label{sub:66}
Für $\alpha,\beta \ge 0$, $p \ge r >0$ gilt $(\alpha^p + \beta^p)^{\frac{1}{p}} \le \enbrace*{\alpha^r + \beta^r}^{\frac{1}{r}}$.\index{Jensen-Ungleichung}
\minisec{Beweis}
Setze $A := \enbrace*{\alpha^r +\beta^r}^{\frac{1}{r}}$ ($\not=0$ oBdA). Dann gilt 
\[
	1=\frac{1}{A^r} \cdot \enbrace*{\alpha^r+\beta^r} = \enbrace[\bigg]{\underbrace{\frac{\alpha}{A}}_{\le 1}}^r + \enbrace[\bigg]{\underbrace{\frac{\beta}{A}}_{\le 1}}^r
	\ge \enbrace[\bigg]{\frac{\alpha}{A} }^p + \enbrace[\bigg]{\frac{\beta}{A} }^p    
\]
$\Rightarrow \enbrace*{\frac{1}{A^p} \cdot \enbrace*{\alpha^p + \beta^p}}^{\frac{1}{p} } \le 1 = \enbrace*{\frac{1}{A^r} \enbrace*{\alpha^r+ \beta^r}}^{\frac{1}{r} }$. 
Kürzen von $\frac{1}{A}$ liefert die Behauptung. \bewende
% subsection 66 (end)

\subsection{Proposition} % (fold)
\label{sub:67}
Sei $(X,\Sigma,\mu)$ ein Maßraum, $1 < p,q < \infty$ mit $\frac{1}{p} +\frac{1}{q}=1$ und $f,g \in L^p(\mu)$.
\begin{enumerate}[(i)]
	\item Falls $p\ge 2$, so gilt $\norm[p]{f+g}^p + \norm[p]{f-g}^p \le 2^{p-1} \enbrace*{\norm[p]{f}^p + \norm[p]{g}^p}$.
	\item Falls $p < 2$, so gilt $\norm[p]{f+g}^q + \norm[p]{f-g}^q \le 2 \cdot \enbrace*{\norm[p]{f}^p + \norm[p]{g}^p}^{q-1}$.
\end{enumerate}
\minisec{Beweis}
\begin{enumerate}[(i)]
	\item Nach \ref{sub:66} gilt für $\gamma,\delta  \in \mathds{R}$
	\begin{equation*}
		\enbrace[\big]{\abs*{\gamma+\delta }^p + \abs*{\gamma-\delta }^p}^{\frac{1}{p}} \le \enbrace*{\abs*{\gamma+\delta }^2 + \abs*{\gamma-\delta}^2}^{\frac{1}{2}}
		= 2^{\frac{1}{2}} \cdot \enbrace*{\gamma^2 + \delta^2}^{\frac{1}{2}} \label{eq:67} \tag{*}
	\end{equation*}
	$[p>2]$ Sei $r := \frac{p}{p-2}$, sodass $\frac{1}{\frac{p}{2}} + \frac{1}{r}=1$. Die Höldersche Ungleichung (\ref{sub:45}) für $\mathds{R}^2$ liefert für 
	$(\gamma^2,\delta^2)$ und $(1,1)$
	\[
		\gamma^2 \cdot 1 + \delta^2 \cdot 1 \le \enbrace*{\enbrace{\gamma^2}^{\frac{p}{2} } + \enbrace*{\delta^2}^{\frac{p}{2}} }^{\frac{2}{p}} 
		\cdot (1^r+ 1^r)^{\frac{1}{r}} = \enbrace[\big]{\abs*{\gamma}^p + \abs*{\delta}^p}^{\frac{2}{p}} \cdot 2^{\frac{p-2}{p}} 
	\]
	Es folgt 
	\begin{align*}
		\sqrt{2} \cdot \enbrace*{\gamma^2+ \delta^2}^{\frac{1}{2}} \le 2^{\frac{p}{2p}} \cdot \enbrace[\big]{\abs*{\gamma}^p + \abs*{\delta }^p}^{\frac{1}{p}} \cdot 2^{\frac{p-2}{2p} } = 2^{\frac{p-1}{p}} \cdot \enbrace[\big]{\abs*{\gamma}^p +\abs*{\delta}^p}^{\frac{1}{p} } 
	\end{align*}
	Mit \eqref{eq:67} folgt $\abs*{\gamma+\delta}^p + \abs*{\gamma-\delta}^p \le 2^{p-1} \cdot  \enbrace[\big]{\abs*{\gamma}^p + \abs*{\delta}^p}$. Integration liefert (i).
	
	Für $p=2$ folgt die Aussage mit Pythagoras.
	\item Übung. \bewende
\end{enumerate}
% subsection 67 (end)

\subsection[Satz: Für $1 <p <\infty$ ist $L^p(\mu)$ gleichmäßig konvex]{Satz} % (fold)
\label{sub:68}
Sei $(X,\Sigma,\mu)$ ein Maßraum, $1 <p < \infty$. Dann ist $L^p(\mu)$ gleichmäßig konvex.
\minisec{Beweis}
Seien $(f_n)_\mathds{N}, (g_n)_\mathds{N} \subset L^P(\mu)$ Folgen mit $\norm[p]{f_n}= \norm[p]{g_n}= 1 $ für $n \in \mathds{N}$ und 
$\lim_{ n \to \infty} \norm[p]{\frac{1}{2} (f_n+g_n)}=1$. Zu zeigen: $\norm[p]{f_n-g_n} \xrightarrow{n \to \infty} 0$, dann gilt \ref{sub:61}(iii). Nach \ref{sub:67}
gilt 
\[
	\underbrace{\norm[p]{f_n+g_n}^p}_{\xrightarrow{n \to \infty} 2^p} + \norm[p]{f_n-g_n}^p \le 2^{p} 
\]
Falls $p\ge 2$ bzw. für $p < 2$
\[
	\underbrace{\norm[p]{f_n+g_n}^q }_{\xrightarrow{n \to \infty} 2^q}+ \norm[p]{f_n-g_n}^q \le 2^q
\]
Also folgt $\norm[p]{f_n -g_n}^p \xrightarrow{n \to \infty} 0$ bzw. $\norm[p]{f_n-g_n}^q \xrightarrow{n \to \infty} 0$. Also gilt 
$\norm[p]{f_n -g_n} \xrightarrow{n \to \infty} 0$. \bewende
% subsection 68 (end)

\subsection{Satz} % (fold)
\label{sub:69}
Sei $X$ ein gleichmäßig konvexer Banachraum. Dann ist $X$ reflexiv.
\minisec{Beweis}
Sei $f \in X^{**}$, $\norm[X^{**}]{f}=1$. Zu zeigen: $f \in \iota(X)$. Wir wählen $\varphi_n \in X^*$ mit $\norm[X^*]{\varphi_n}=1$ und $f(\varphi_n) > 1- \frac{1}{n+1}$
für $n \in \mathds{N}$. Nach (dem Beweis von) \ref{sub:516} gilt für $n \in \mathds{N}$
\[
	\inf_{x \in \overline{B}_X(0,1) } \sum_{i=0}^{n} \,\abs[\big]{f(\varphi_i) - \varphi_i(x)}^2 =0
\]
$\Rightarrow \exists (x_n)_{n \in \mathds{N}} \subset X$, sodass $\norm[X]{x_n}=1$ und $\abs*{f(\varphi_m)- \varphi_m(x_n)} < \frac{1}{2(n+1)}$ für $0 \le m \le n$ also auch
$1- \frac{3}{2(m+1)}\le \varphi_m(x_n) \le 1$. Behauptung: $(x_n)_\mathds{N}$ ist Cauchy. Es gilt für $n \ge m$
\begin{align*}
	\underbrace{1- \frac{3}{2(m+1)}  + 1 - \frac{3}{2(m+1)}}_{\xrightarrow{n,m \to \infty} 2} \le \varphi_m(x_m) + \varphi_m(x_n) = \varphi_m(x_m+x_n) &\le \norm{\varphi_m} \cdot \norm[X]{x_m+x_n} \\
	&\le \norm[X]{x_m} + \norm[X]{x_n} = 2  
\end{align*}
$\Rightarrow \norm[X]{\frac{1}{2} (x_m+x_n)} \xrightarrow{n,m \to \infty} 1$. Mit Bemerkung \ref{sub:65} folgt, dass $(x_n)_\mathds{N}$ Cauchy ist. Sei 
$\overline{x} := \lim_{ n \to \infty} x_n$. Es gilt $\norm[X]{\overline{x}}=1$. Behauptung: $\iota_X(\overline{x})=f$. Für $n \ge m$ gilt
\begin{align*}
	\abs[\Big]{f(\varphi_m) - \underbrace{\varphi_m(x_n)}_{\xrightarrow{n \to \infty} \varphi_m(\overline{x} )}} < \frac{1}{2(n+1)}  
\end{align*}
$\varphi_m$ ist stetig, also ist $f(\varphi_m)=\varphi_m(\overline{x})= \iota(\overline{x})(\varphi_m)$ für jedes $m \in\mathds{N}$. Sei $\tilde{x} \in X$ ein weiterer Punkt
mit $\norm{\tilde{x}}=1$ und $f(\varphi_m) = \varphi_m(\tilde{x})$, $m \in \mathds{N}$. Das obige Argument für die Folge 
$(\overline{x}, \tilde{x}, \overline{x}, \tilde{x}, \ldots)$ liefert: $(\overline{x}, \tilde{x}, \overline{x}, \tilde{x}, \ldots)$ ist Cauchy 
$\Rightarrow \overline{x}=\tilde{x}$.

Noch zu zeigen: $f(\varphi)= \varphi(\overline{x})$ für alle $\varphi \in X^*$. Das obige Argument für $(\varphi, \varphi_1, \varphi_2, \ldots )$ an Stelle von
$(\varphi_0, \varphi_1, \varphi_2,\ldots )$ liefert eine Folge $(y_n)_\mathds{N} \subset X$ mit $\norm{y_n}=1$, $(y_n)_\mathds{N}$ ist Cauchy; 
Dann erfüllt $\overline{y} := \lim_{ n \to \infty} y_n$ die Gleichung $f(\varphi_m)=\varphi_m(\overline{y})$, $m \ge 1$. $f(\varphi)=\varphi(\overline{y})$.
$\Rightarrow \overline{y}=\overline{x}$. Damit folgt $\iota(\overline{x})(\varphi)= \varphi(\overline{x})= \varphi(\overline{y})= f(\varphi)$ für jedes $\varphi \in X^*$,
also ist $\iota(\overline{x})=f$. \bewende
% subsection 69 (end)

\subsection[Corollar: $L^p(\mu)$ ist für $1<p<\infty$ reflexiv.]{Corollar} % (fold)
\label{sub:610}
Für einen Maßraum $(X,\Sigma, \mu)$ und $1 < p< \infty$ ist $L^p(\mu)$ reflexiv.
\minisec{Beweis}
Kombiniere die Aussagen aus \ref{sub:68} und \ref{sub:69}. \bewende
% subsection 610 (end)

\subsection[Satz: Isometrischer Isomorphismus $\iota_p \colon L^p(\mu) \to L^q(\mu)^*$]{Satz} % (fold)
\label{sub:611}
Sei $(X,\Sigma,\mu)$ ein Maßraum, $\frac{1}{p} + \frac{1}{q}=1$. Dann ist die Abbildung $\iota_p \colon L^p(\mu) \to L^q(\mu)^*$ gegeben durch 
$\iota_p(f)(g) := \int \!f g \,\mathd\mu$ ein isometrischer Isomorphismus.
\minisec{Beweis}
$\iota_p$ ist eine lineare Isometrie nach Proposition \ref{sub:410}, wir müssen also nur Surjektivität zeigen. $\iota_p \enbrace*{L^p(\mu)}$ ist vollständig, also 
abgeschlossen in $L^q(\mu)^*$. Falls $\iota_p(L^p(\mu)) \subsetneq L^q(\mu)$, so existiert nach Hahn-Banach (\ref{sub:227}) ein $0 \not= \psi \in L^q(\mu)^{**}$ mit
$\psi \enbrace*{\iota_p(L^p(\mu))} =0$. $L^q(\mu)$ ist reflexiv, also existiert $g \in L^q(\mu)$ mit $\iota(g)=\psi$, d.h. $\varphi(g)= \iota_p(g)(\varphi)=\psi(\varphi)$,
$\varphi \in L^q(\mu)^*$. Für $f\in L^p(\mu)$ gilt dann
\[
	0 = \psi \enbrace*{\iota_p(f)} = \iota_p(f)(g) = \int\! f g \, \mathd\mu \enspace \Longrightarrow g=0 \light \text{ zu } \psi \not= 0 \bewende
\]
% subsection 611 (end)

\subsection{Bemerkung} % (fold)
\label{sub:612}
Weiter kann man zeigen, dass die kanonische Abbildung $\iota_\infty : L^\infty(\mu) \to L^1(\mu)^*$, $\iota_\infty(f)(g) = \int\! fg \, \mathd\mu$, ein isometrischer
Isomorphismus ist. \hfill \emph{(ohne Beweis)}\\
Die Isometrie $\iota_1 : L^1(\mu) \to L^\infty(\mu)^*$ ist jedoch \emph{nicht} surjektiv! \hfill \emph{(Übung)}
% subsection 612 (end)
% section 6 (end)
\newpage
\section{Hilberträume und selbstadjungierte Operatoren} % (fold)
\label{sec:7}

\subsection[Definition: Hermitesche Form, Skalarprodukt]{Definition} % (fold)
\label{sub:71}
Sei $X$ ein $\mathds{K}$-Vektorraum. Eine \Index{hermitesche Form} auf $X$ ist eine Abbildung $\skal*{\cdot}{\cdot} \colon X \times X \to \mathds{K}$ mit\todo{RevChap7}
\begin{align*}
	\skal*{x}{y+y'} &= \skal*{x}{y} + \skal*{x}{y'} \\
	\skal*{x}{\alpha \cdot y} &= \alpha \cdot \skal*{x}{y} \\
	\skal*{x}{y} &= \overline{\skal*{y}{x}} \qquad \text{ für } x,y,y' \in X, \alpha \in \mathds{K}       
\end{align*}
$\skal*{\cdot}{\cdot } $ heißt \bet{positiv semidefinit}\index{hermitesche Form!positiv semidefinit}, falls $\skal*{x}{x}\ge 0$ für $x \in X$ und 
\bet{positiv definit}\index{hermitesche Form!positiv definit}, falls 
$\skal*{x}{x} >0$ für $0 \not= x \in X$. In diesem Fall heißt $\skal*{\cdot }{\cdot }$ auch \Index{Skalarprodukt}.  
% subsection 71 (end)

\subsection[Bemerkung: Einfache Eigenschaften von hermiteschen Formen]{Bemerkung} % (fold)
\label{sub:72}
Für eine hermitesche Form $\skal*{\cdot}{\cdot}$ und $x,x',y \in X$, $\alpha \in \mathds{K}$ gilt 
\begin{align*}
	\skal*{x+x'}{y} = \skal*{x}{y} + \skal*{x'}{y} \quad , \quad \qquad 
	\skal*{\alpha \cdot x}{y} = \overline{\alpha} \cdot \skal*{x}{y} \qquad \text{und} \qquad 
	\skal*{x}{x} \in \mathds{R} 
\end{align*}
% subsection 72 (end)

\subsection[Beispiele für hermitesche Formen]{Beispiele} % (fold)
\label{sub:73}
\begin{enumerate}[(i)]
	\item Auf $\mathds{C}^n$ definiert $\skal*{\underline{x}}{\underline{y}} := \sum_{i=1}^{n} \overline{x}_i \cdot y_i$ ein Skalarprodukt.
	\item Auf $L^2(\mathds{R})$ definiert $\skal*{f}{g} := \int\! fg \,\mathd\lambda$ ein Skalarprodukt.
	Auf 
	\[
		L^2_\mathds{C}(\mathds{R}) := \set[f \colon \mathds{R} \to \mathds{C}]{\re f \text{ und } \im f \text{ messbar}, \int \! \bar{f} f \mathd\mu < \infty } / \mathcal{N}
	\]
	definiert $\skal*{f}{g} := \int\! \bar{f} g \, \mathd\lambda$ ein Skalarprodukt.
	\item Auf $\ell^2_{\color{gray}\mathds{C}}(\mathds{N})$ definieren wir ein Skalarprodukt durch $\skal*{(a_n)}{(b_n)} := \sum_{n \in \mathds{N}} \overline{a}_n b_n$.
\end{enumerate}
% subsection 73 (end)

\subsection{Proposition (Cauchy-Schwarz-Ungleichung)} % (fold)
\label{sub:74}
Sei $X$ ein $\mathds{K}$-Vektorraum mit einer positiv semidefiniten hermiteschen Form. Dann gilt für $x,y \in X$ \index{Cauchy-Schwarz-Ungleichung}
\[
	\abs[\big]{\skal*{x}{y}}^2 \le \skal*{x}{x} \cdot \skal*{y}{y}
\]
\minisec{Beweis}
Für $\lambda \in \mathds{K}$, $x,y \in X$ gilt: 
\begin{align*}
0 \le \skal*{x+\lambda y}{x + \lambda y} = \skal*{x}{x} + \lambda \cdot \skal*{x}{y} + \overline{\lambda} \cdot \skal*{y}{x} + \lambda \overline{\lambda} \skal*{y}{y}  
\end{align*}
Falls $\skal*{y}{y} \not=0$, so setze $\lambda := - \frac{\skal*{y}{x}}{\skal*{y}{y}}$. Es folgt 
\begin{align*}
	0 \le \skal*{x}{x} - \frac{\skal*{y}{x} \cdot \skal*{x}{y}}{\skal*{y}{y}} + \overline{\lambda} \enbrace*{\skal*{y}{x}  - \frac{\skal*{y}{x}}{\skal*{y}{y}} \skal*{y}{y}   } = \skal*{x}{x} - \frac{\abs*{\skal*{x}{y}}^2 }{\skal*{y}{y}} 
\end{align*}
Falls $\skal*{x}{x} \not= 0$, so vertausche die Rollen von $x$ und $y$. Falls $\skal*{x}{x}=\skal*{y}{y}=0$ gilt, so setzen wir $\lambda := - \skal*{y}{x}$. \bewende
% subsection 74 (end)

\subsection[Proposition: Ein Skalarprodukt definiert eine Norm]{Proposition} % (fold)
\label{sub:75}
Sei $X$ ein $\mathds{K}$-Vektorraum mit Skalarprodukt. Dann definiert $x \mapsto \norm{x} := \skal*{x}{x}^{\frac{1}{2}}$ eine Norm auf $X$.
\minisec{Beweis}
$\norm{\cdot }$ ist positiv definit, also ist auch $\skal*{\cdot}{\cdot}$ positiv definit. Weiter gilt
\begin{align*}
	\norm{\lambda \cdot x} &= \skal*{\lambda \cdot x}{\lambda \cdot x}^{\frac{1}{2}} = \enbrace*{\overline{\lambda} \lambda \cdot \skal*{x}{x}}^{\frac{1}{2}} 
	= \abs*{\lambda} \cdot \norm{x}    \\
	\shortintertext{Die Dreieckungleichung gilt, da}
	\norm{x+y}^2 &= \skal*{x+y}{x+y} = \skal*{x}{x} + \skal*{x}{y} + \skal*{y}{x} + \skal*{y}{y} \le \skal*{x}{x} +  \enbrace*{\skal*{x}{x} \skal*{y}{y}}^{\frac{1}{2}}
	+ \skal*{y}{y}  \\
	&= \norm{x}^2 + 2 \cdot \norm{x} \cdot \norm{y} + \norm{y}^2 = \enbrace*{\norm{x} + \norm{y}}^2    
\end{align*}
$\Rightarrow \norm{x+y} \le \norm{x} + \norm{y}$. \bewende
% subsection 75 (end)

\subsection[Bemerkungen über Cauchy-Schwarz und die Stetigkeit von $\skal*{\cdot}{\cdot}$]{Bemerkung} % (fold)
\label{sub:76}
\begin{enumerate}[(i)]
	\item Die Cauchy-Schwarz Ungleichung schreibt sich auch als $\abs*{\skal*{x}{y}} \le \norm{x} \cdot \norm{y}$.
	\item Die Abbildung $\skal*{\cdot}{\cdot} \colon X \times X \to \mathds{K}$ ist stetig:
	\begin{align*}
		\abs*{\skal*{x}{y}- \skal*{x'}{y'}} = \ldots \le \norm{x-x'} \cdot \norm{y} + \norm{x'} \cdot \norm{y'-y}.     
	\end{align*}
\end{enumerate}
% subsection 76 (end)

\subsection[Proposition über Zusammenhang von Norm und Skalarprodukt mit dem Parallelogrammgesetz]{Proposition} % (fold)
\label{sub:77}
Für einen normierten Raum $X$ sind äquivalent: \todo{eventuell Bild von Parallelogramm hinzufügen}
\begin{enumerate}[(i)]
	\item Es gilt das \Index{Parallelogrammgesetz}: $\norm{x+y}^2 + \norm{x-y}^2 = 2 \enbrace*{\norm{x}^2 + \norm{y}^2}$ für alle $x,y \in X$.
	\item $\skal*{x}{y} := \frac{1}{4} \enbrace*{\norm{x+y}^2 - \norm{-x+y}^2 + i \cdot \norm{i \cdot x +y}^2 - i \cdot \norm{-i \cdot x+ y}^2}$ für $\mathds{K}=\mathds{C}$ bzw.
	
	$\skal*{x}{y} := \frac{1}{4} \enbrace*{\norm{x+y}^2 - \norm{x-y}^2}$ für $\mathds{K}=\mathds{R}$ definiert ein Skalarprodukt auf $X$. In diesem Fall gilt
	$\skal*{x}{x} = \norm{x}^2$ für alle $x \in X$; insbesondere ist $\skal*{\cdot}{\cdot}$ stetig auf $X \times X$.
\end{enumerate}
\minisec{Beweis}
\begin{description}
	\item[(ii)$\Rightarrow$(i):] Es gilt $\norm{x}^2 = \skal*{x}{x}$ für $x \in X$. und (i) folgt durch einfache Rechnung.
	\item[(i)$\Rightarrow$(ii):] $\skal*{\cdot}{\cdot} \colon X \times X \to X$ ist stetig (klar). Linearität (für $\mathds{K}=\mathds{C}$):
	\begin{gather*}
		\skal*{x}{y+y'} = \frac{1}{4} \enbrace*{\norm{x+y+y'}^2 - \norm{-x+y+y'}^2 + i  \norm{i x + y+ y'}^2 - i \norm{- ix +y+y'}^2} \label{eq:77:1}\tag{*} \\ 
		\begin{split}
			\skal*{x}{y} + \skal*{x}{y'} = \frac{1}{4} \Big( &\norm{x+y}^2 - \norm{-x+y}^2 + i  \norm{ix +y}^2 - i  \norm{-ix +y}^2 + \\
			&\norm{x+y'}^2 + \norm{-x+y'} + i  \norm{ix + y'} - i  \norm{-ix +y'} \Big)
		\end{split}
		\label{eq:77:2}\tag{**}
	\end{gather*}
	Vergleich der Imaginärteile: Es gilt
	\begin{align*}
		\norm{x+y+y'}^2 &= - \norm{x+y-y'}^2 + 2 \norm{x+y}^2 + 2 \norm{y'}  \\
		\shortintertext{und}
		\norm{x+y+y'}^2 &= - \norm{x-y+y'}^2 + 2 \norm{x+y'} + 2 \norm{y}^2 \\
		&= - \frac{1}{2} \enbrace*{\norm{x+y-y'}^2 + \norm{x-y+y'}^2} + \norm{x+y}^2 + \norm{x+y'}^2 + \norm{y'}^2 + \norm{y}^2         
	\end{align*}
	$\Rightarrow$ Realteile von \eqref{eq:77:1} und \eqref{eq:77:2} stimmen überein. $\ldots \Rightarrow \skal*{x}{y} + \skal*{x}{y'}= \skal*{x}{y+y'}$.
	$\Rightarrow \skal*{x}{\alpha \cdot y} = \alpha \cdot \skal*{x}{y}$ für $\alpha \in \mathds{N} \setminus \set{0}$ ($\skal*{x}{0}=0$ ist trivial). Damit folgt auch
	$\skal*{x}{\alpha \cdot y} = \alpha \cdot \skal*{x}{y}$ für $\alpha \in \mathds{Z}$. Es gilt nun für $\alpha=\frac{m}{n} \in \mathds{Q}$
	\[
		\skal*{x}{\alpha \cdot y} = \frac{n}{n} \cdot \skal*{x}{\frac{m}{n} \cdot y } = \frac{m}{n} \cdot \skal*{x}{\frac{n}{n}\cdot y } = \alpha \cdot \skal*{x}{y}  
	\]
	Mit der Stetigkeit von $\skal*{\cdot}{\cdot}$ folgt nun $\skal*{x}{\alpha \cdot y} = \alpha \cdot \skal*{x}{y}$ für $\alpha \in \mathds{R}$.
	$\skal*{x}{i \cdot y} = i \cdot \skal*{x}{y}$ ist trivial und damit ist $\skal*{x}{\alpha \cdot y }= \alpha \cdot \skal*{x}{y}$ für $\alpha \in \mathds{C}$. Also ist
	$\skal*{\cdot}{\cdot}$ linear in der zweiten Variablen. $\skal*{x}{y} = \overline{\skal*{y}{x}}$ ist trivial. \bewende
\end{description}
% subsection 77 (end)

\subsection{Bemerkung} % (fold)
\label{sub:78}
Ist $X$ ein $\mathds{C}$-Vektorraum und $q \colon X \times X \to \mathds{C}$ eine \Index{Sesquilinearform} (d.h. $q$ ist linear in der zweiten Variable und konjugiert 
linear in der ersten), dann ist $q$ durch die Abbildung $n_q \colon X \to \mathds{C}$, $n_q(x) := q(x,x)$ eindeutig bestimmt durch die \Index{Polarisierungsidentität}
\[
	q(x,y) = \frac{1}{4} \enbrace[\Big]{n_q(x+y) - n_q(-x+y) + i \cdot n_q(i x +y) - i \cdot n_q(-ix +y)}.  
\]
% subsection 78 (end)

\subsection[Definition: Prä-Hilbertraum und Hilbertraum]{Definition} % (fold)
\label{sub:79}
Ein $\mathds{K}$-Vektorraum $\mathcal{H}$ heißt \Index{Prä-Hilbertraum}, falls $\mathcal{H}$ ein Skalarprodukt besitzt. $\mathcal{H}$ heißt \Index{Hilbertraum}, falls
$\mathcal{H}$ vollständig ist bzgl. der durch $\skal*{\cdot}{\cdot}$ induzierten Norm. 
% subsection 79 (end)

\subsection{Proposition} % (fold)
\label{sub:710}
Sei $\mathcal{H}$ ein Prä-Hilbertraum und $\overline{\mathcal{H}}^{\norm{\cdot}}$ die Vervollständigung bezüglich $\norm{x} = \skal*{x}{x}^{\frac{1}{2}}$. Dann ist
$\overline{\mathcal{H}}^{\norm{\cdot}}$ in kanonischer Weise ein Hilbertraum.
\minisec{Beweis}
$\enbrace*{\mathcal{H}, \norm{\cdot}} \hookrightarrow \enbrace*{\overline{\mathcal{H}}^{\norm{\cdot}}, \norm[\overline{\mathcal{H}}^{\norm{\cdot}}]{\cdot}}$ ist isometrisch
mit dichtem Bild. $\norm{\cdot}$ erfüllt das Parallelogrammgesetz und daher auch $\norm[\overline{\mathcal{H}}^{\norm{\cdot}}]{\cdot}$. Mit \ref{sub:77} folgt, dass 
$\overline{\mathcal{H}}^{\norm{\cdot}}$ ein Skalarprodukt besitzt, welches die Norm $\norm[\overline{\mathcal{H}}^{\norm{\cdot}}]{\cdot}$ induziert und auf $\mathcal{H}$ mit
$\skal*{\cdot}{\cdot}$ übereinstimmt. Das Skalarprodukt auf $\overline{\mathcal{H}}^{\norm{\cdot}}$ ist wegen Stetigkeit hierdurch eindeutig bestimmt. \bewende
% subsection 710 (end)

\subsection{Proposition} % (fold)
\label{sub:711}
Sei $\mathcal{H}$ ein Hilbertraum und $\kappa \colon \mathcal{H} \to \mathcal{H}^*$ definiert durch $\kappa(x)(y) := \skal*{x}{y}$, $x,y \in \mathcal{H}$. Dann gilt
\begin{enumerate}[(i)]
	\item $\kappa$ ist \Index{konjugiert linear}, d.h. $\kappa(\alpha x+ \alpha' x') = \overline{\alpha} \kappa(x) + \overline{\alpha'}\kappa(x')$ mit $x,x' \in \mathcal{H}$, 
	$\alpha, \alpha' \in \mathds{K}$.
	\item $\kappa$ ist isometrisch (insbesondere injektiv).
	\item $\kappa$ ist surjektiv.
\end{enumerate}
\minisec{Beweis}
Es gilt $\kappa(x) \in \mathcal{H}^*$, da $\kappa(x) \colon \mathcal{H} \to \mathds{K}$ offensichtlich linear ist. Weiter ist $\kappa(x)$ stetig, da
\begin{align*}
	\abs*{\kappa(x)(y)}= \abs*{\skal*{x}{y}}  \le \norm{x} \cdot \norm{y}  \Longrightarrow \norm[\mathcal{H}^*]{\kappa(x)} \le \norm{x}  
\end{align*}
\begin{enumerate}[(i)]
	\item klar.
	\item Es gilt $\kappa(x) \enbrace*{\frac{1}{\norm{x}}x }= \frac{1}{\norm{x}} \skal*{x}{x} = \norm{x}$. Daraus folgt $\norm[\mathcal{H}^*]{\kappa(x)} \ge \norm{x}$, also
	auch $\norm[\mathcal{H}^*]{\kappa(x)}= \norm{x}$.
	\item Sei $\varphi \in \mathcal{H}^*$. Ohne Einschränkungen gelte $\norm[\mathcal{H}^*]{\varphi}=1$. 
	
	Behauptung: Es gibt $\overline{x} \in \mathcal{H}$ mit $\norm{\overline{x}}=1 $ und
	$\varphi(\overline{x} )=1$. Wähle $(x_n)_\mathds{N} \subset \mathcal{H}$ mit $\norm{x_n}=1$ und $0\le \varphi(x_n) \xrightarrow{n \to \infty} 1$.
	Zu $1 >\varepsilon>0$ existiert $N \in \mathds{N}$ mit $\varphi(x_n) > 1- \frac{\varepsilon}{8} $ für $n \ge N$, also $\varphi(x_n +x_m) > 2- \frac{\varepsilon}{4}$.
	Aus dem Parallelogrammgesetz folgt für $n,m \ge N$
	\begin{align*}
		\norm{x_n -x_m}^2 = 2 \cdot \norm{x_n}^2 + 2 \cdot \norm{x_m}^2 - \norm{x_n +x_m}^2 \le 4 - \enbrace*{2-\frac{\varepsilon}{4} }^2 \le \varepsilon 
	\end{align*}
	Es folgt, dass $(x_n)_\mathds{N}$ eine Cauchyfolge ist. $\overline{x} := \lim_{ n \to \infty} x_n$ erfüllt die Behauptung.
	
	Behauptung: $\varphi= \kappa(\overline{x})$, d.h. $\varphi(y)= \kappa(\overline{x})(y)$ für $y \in \mathcal{H}$. Sei $\mathds{K}=\mathds{R}$. Für $\lambda>0$ gilt
	wegen der Dreiecksungleichung und $\norm[\mathcal{H}*]{\varphi}\le 1$
	\begin{align*}
		-\frac{1}{\lambda} \enbrace*{\norm{\overline{x} - \lambda \cdot y } - \norm{\overline{x}}} \le -\frac{1}{\lambda} \enbrace*{\varphi \enbrace*{\overline{x}- \lambda 
		\cdot y} - \varphi(\overline{x})} = \varphi(y) \\
		\shortintertext{und}
		\varphi(y) = \frac{1}{\lambda} \enbrace*{\varphi(\overline{x}+\lambda \cdot y) - \underbrace{\varphi(\overline{x})}_{=1}} 
		\le \frac{1}{\lambda} \cdot \enbrace[\Big]{\norm{\overline{x} + \lambda \cdot y } - \underbrace{\norm{\overline{x}}}_{=1} }    \\
		\shortintertext{also insgesamt}
		 \underbrace{-\frac{1}{\lambda} \enbrace[\Big]{\norm{\overline{x} - \lambda \cdot y } - \norm{\overline{x}}} }_{\xrightarrow{\lambda \to 0} \skal*{\overline{x}}{y}}
		 \le \varphi(y) 
		 \le \underbrace{\frac{1}{\lambda} \cdot \enbrace[\Big]{\norm{\overline{x} +\lambda\cdot y} - \norm{\overline{x}}}}_{\xrightarrow{\lambda \to 0} \skal*{\overline{x}}{y}}
	\end{align*}
	nach l'Hospital. Im Detail: \todo{stimmt noch nicht so ganz, aber passt schon …}
	\begin{align*}
		\norm{\overline{x} + \lambda y} - \norm{\overline{x}} = \skal*{\overline{x} + \lambda y }{\overline{x} + \lambda y }^{\frac{1}{2}} - \norm{\overline{x}}= \enbrace*{\skal*{\overline{x}}{\overline{x}} +
		\skal*{\overline{x}}{\lambda y} + \lambda^2 \skal*{y}{y}}^{\frac{1}{2}} - \skal*{\overline{x}}{\overline{x}}^{\frac{1}{2}}
	\end{align*}
	ist in $\lambda=0$ differenzierbar nach $\lambda$ mit 
	\[
		\diffd{}{t}(\cdots) = \enbrace*{2 \skal*{\overline{x}}{y} + 2 \lambda \norm{y}^2}^{-\frac{1}{2} } \frac{1}{2} \enbrace*{\cdots}^{-\frac{1}{2}} 
		\xrightarrow{\lambda \to 0}  \skal*{\overline{x}}{y}
	\]
	Wir erhalten $\varphi(y)= \skal*{\overline{x}}{y} = \kappa(\overline{x})(y)$, also $\varphi=\kappa(\overline{x})$. 
	
	Für $\mathds{K}=\mathds{C}$: $\mathcal{H}$ ist reeller Vektorraum und $\skal*{\cdot}{\cdot}_\mathds{R} := \re \enbrace*{\skal*{\cdot}{\cdot}}$. Es gilt dann 
	$\varphi(y) = \rho(y) - i \cdot \rho(y)$ mit $\rho(\cdot ) := \re \enbrace*{\varphi(\cdot )}$. Es gilt nämlich
	\[
		\im \enbrace*{\varphi(y)} = - \re \enbrace*{i \cdot \varphi(y)} = - \re \enbrace*{\varphi(i \cdot y)} = - \rho (i \cdot y)   
	\]
	Es ist $\rho(y) = \skal*{\overline{x}}{y}_\mathds{R} = \re \enbrace*{\skal*{\overline{x}}{y}}$ für ein $\overline{x} \in \mathcal{H}$. Also gilt
	$\varphi(y)= \rho(y)- i \cdot \rho(i \cdot y)= \ldots = \skal*{\overline{x}}{y}$. \bewende
\end{enumerate}
% subsection 711 (end)

\subsection{Korollar} % (fold)
\label{sub:712}
Hilberträume sind reflexiv.
\minisec{Beweis}
Das Diagramm
\[
	\begin{tikzcd}[row sep=large, column sep=large]
		\mathcal{H} \rar["\iota_{\mathcal{H}}"] \dar["\kappa"] & \mathcal{H}^{**} \dar["\kappa^\tr"]\\
		\mathcal{H}^* \rar["\varphi \mapsto \overline{\varphi}"] & \mathcal{H}^*
	\end{tikzcd}
\]
kommutiert, da gilt
\begin{align*}
	\kappa^\tr \circ \iota_{\mathcal{H}}(x)(y) = \kappa^\tr \enbrace*{\iota_{\mathcal{H}}(x)}(y)  =\iota_{\mathcal{H}}(x) \circ \kappa (y) 
	= \kappa(y)(x) = \skal*{y}{x} = \overline{\skal*{x}{y}} = \overline{\kappa(x)(y)} 
\end{align*}
Die Abbildungen $\kappa, \kappa^\tr$ und $(\varphi \mapsto \overline{\varphi})$ sind surjektive Isometrien, also ist auch 
$\iota_{\mathcal{H}} \colon \mathcal{H} \to \mathcal{H}^{**}$ eine surjektive Isometrie.\bewende
% subsection 712 (end)

\subsection[Definition: Orthogonal, Orthonormalsystem und Hilbertraumbasis]{Definition} % (fold)
\label{sub:713}
Sei $\mathcal{H}$ ein Hilbertraum.
\begin{enumerate}[(i)]
	\item $x,y \in \mathcal{H}$ heißen \Index{orthogonal}, $x\bot y$, falls $\skal*{x}{y}=0$.
	\item $A,B  \subset \mathcal{H}$ heißen \Index{orthogonal}, falls $x \bot y$ für alle $x \in A$, $y \in B$. 
	\item Für $A \subset \mathcal{H}$ heißt $A^\bot := \set[x \in \mathcal{H}]{\skal*{x}{y}=0 \text{ für alle } y \in A}$ \Index{orthogonales Komplement} von $A$.
	\item $\set{x_i}_{I} \subset \mathcal{H}$ heißt \Index{Orthonormalsystem}, falls $\skal*{x_i}{x_j}= \delta_{ij}$, $i,j \in I$.
	\item $\set{x_i}_{I} \subset \mathcal{H}$ heißt \Index{Hilbertraumbasis} (bzw. Basis), falls es ein maximales Orthonormalsystem ist. 
\end{enumerate}
% subsection 713 (end)

\subsection{Bemerkung} % (fold)
\label{sub:714}
\begin{enumerate}[(i)]
	\item Für $x \bot y \in \mathcal{H}$ gilt der \Index{Satz des Pythagoras}: $\norm{x+y}^2 = \norm{x}^2  + \norm{y}^2$.
	\item Wir werden sehen, dass jeder Hilbertraum eine Hilbertraumbasis besitzt.
	\item Im Allgemeinen ist eine Hilbertraumbasis \emph{keine} Vektorraumbasis (\Index{Hamelbasis}).
\end{enumerate}
% subsection 714 (end)

\subsection[Proposition über Schnitt und Summe eines Unterraumes mit seinem orthogonalen Komplement]{Proposition} % (fold)
\label{sub:715}
Sei $\mathcal{H}$ ein Hilbertraum, $\mathcal{H}_0 \subset \mathcal{H}$ ein abgeschlossener Unterraum. Dann gilt $\mathcal{H}_0 \cap \mathcal{H}_0^\bot = \set{0}$ und 
$\mathcal{H}_0 + \mathcal{H}_0^\bot = \mathcal{H}$. Wir schreiben in diesem Fall auch $\mathcal{H} = \mathcal{H}_0 \oplus \mathcal{H}_0^\bot$.
\minisec{Beweis}
$\mathcal{H}_0 \cap \mathcal{H}_0^\bot= \set{0}$ ist klar (warum?). Sei $x \in \mathcal{H}$. Dann ist $\kappa(x) \in \mathcal{H}^*$ gegeben durch 
$\kappa(x)(y)=\skal*{x}{y}$, und $\kappa(x)|_{\mathcal{H}_0} \in \mathcal{H}_0^*$. $\mathcal{H}_0 \subset \mathcal{H}$ abgeschlossen, ist also selbst ein Hilbertraum. 
Sei $\kappa_0 \colon \mathcal{H}_0 \to \mathcal{H}_0^*$ der konjugiert lineare, isometrische Isomorphismus aus \ref{sub:711}. Nach \ref{sub:711} existiert 
$x_0 \in \mathcal{H}_0$ mit $\kappa_0(x_0)= \kappa(x)|_{\mathcal{H}_0}$. Das heißt für $y \in \mathcal{H}_0$ gilt 
\begin{align*}
	\skal*{x_0}{y} = \kappa_0(x_0)(y) = \kappa(x)\big|_{\mathcal{H}_0}(y) = \kappa(x)(y) = \skal*{x}{y}
\end{align*}
Daraus folgt $\skal*{x-x_0}{y}=0$ für jedes $y \in \mathcal{H_0}$, also ist $x_1 := x-x_0 \in \mathcal{H}_0^\bot$ (und $x=x_0+x_1$). \bewende
% subsection 715 (end)

\subsection[Definition: Summierbar]{Definition} % (fold)
\label{sub:716}
Sei $\mathcal{H}$ ein Hilbertraum. Eine Familie $\set{x_i}_{i \in I} \subset \mathcal{H}$ heißt \Index{summierbar} zu $x \in \mathcal{H}$, falls gilt: 
Zu $\varepsilon>0$ existiert $\overline{J} \finSub I$ endlich, sodass für alle $\overline{J} \subset J \finSub I$ endlich gilt $\norm{\sum_{i \in J}x_i - x } < \varepsilon$.

Wir schreiben dann auch $x= \sum_{i \in I} x_i$. (Warum kann es höchstens ein solches $x$ geben?) Formulierung mit Netzen: 
\[
	\enbrace*{\sum_{i \in J} x_i}_{\!\!\set[J]{J \finSub I \text{ endlich}}} \hspace{-4em} \longrightarrow \enspace x \qquad \text{ oder} \qquad \sum_{i \in J}x_i \xrightarrow{J \finSub I} x 
\]
Dabei ist $\set[J]{J \finSub I \text{ endlich}}$ gerichtet bezüglich Inklusion.
% subsection 716 (end)

\subsection{Proposition} % (fold)
\label{sub:717}
Sei $\mathcal{H}$ ein Hilbertraum, $\set{x_i}_{i \in I} \subset \mathcal{H}$. 
\begin{enumerate}[(i)]
	\item $\set{x_i}_{i \in I}$ ist genau dann summierbar, wenn gilt: Zu jedem $\varepsilon>0$ existiert $\overline{J} \finSub I$ endlich, sodass 
	$\norm{\sum_{i \in J'}x_i}< \varepsilon$ für jedes $J' \finSub I$ endlich mit $J' \cap \overline{J} = \emptyset$.
	\item $\set{x_i}_{i \in I}$ ist genau dann summierbar zu $x \in \mathcal{H}$, wenn gilt: Es sind höchstens abzählbar viele der $x_i \not=0$ und für jede Abzählung 
	$x_{i_0}, x_{i_1}, \ldots $ dieser $x_i$ gilt 
	\[
		x= \lim_{n \to \infty} \sum_{k=0}^{n} x_{i_k}.
	\]
\end{enumerate}
\minisec{Beweis}
\begin{enumerate}[(i)]
	\item Angenommen $\sum_{i\in I} x_i \in \mathcal{H}$ existiert. Sei $\varepsilon>0$. Dann existiert $\overline{J}\finSub I$, sodass für $\overline{J}\subset J\finSub I$ 
	gilt $\norm{\sum_{i \in J} x_i - \sum_{i \in I} x_i} < \frac{\varepsilon}{2}$. Für $J' \finSub I$, $J' \cap \overline{J}=\emptyset$ gilt
	\begin{align*}
		\norm{\sum_{i \in J'} x_i} = \norm{\sum_{i \in J' \cup \overline{J}} x_i - \sum_{i \in \overline{J}} x_i}  \le \norm{\sum_{i \in J' \cup \overline{J}}x_i - \sum_{i \in I} x_i} + \norm{\sum_{i \in \overline{J}} x_i - \sum_{i \in I}x_i}< \varepsilon  
	\end{align*}
	
	Zur Rückrichtung: Für $n \in \mathds{N}$ existiert $J_n \finSub I$ endlich, sodass für $J' \finSub I$ mit $J' \cap J_n = \emptyset$ gilt 
	$\norm{\sum_{i \in J'} x_i}< \frac{1}{n+1}$. Aber dann ist 
	\begin{align*}
		\enbrace*{\sum_{i \in J_0 \cup J_1 \cup \ldots \cup J_n} \hspace{-1em} x_i}_{n \in \mathds{N}} \subset \mathcal{H}
	\end{align*}
	eine Cauchyfolge. Für den Limes $x$ gilt $x= \sum_{i \in I} x_i$ (Warum? Zu $\varepsilon>0$ wähle $n$ mit $\frac{1}{n+1}<\frac{\varepsilon}{2}$ und 
	$\norm{x - \sum_{i \in J_0 \cup J_1 \cup \ldots \cup J_n} x_i}< \frac{\varepsilon}{2}$. Setze $\overline{J}:= J_n$)
	\item Angenommen $\set{x_i}_{i \in I}$ ist summierbar. Seien die $J_n$ wie oben. Dann ist $\bigcup_{n \in \mathds{N}} J_n$ abzählbar. Falls 
	$I \ni i \not\in \bigcup_{n \in \mathds{N}} J_n$, so gilt $\norm{x_i} < \frac{1}{n+1}$ für jedes $n \in \mathds{N}$, also $x_i=0$.
	
	Die Rückrichtung ist klar. (sonst noch einmal an Analysis I. erinnern \ldots ) \bewende
\end{enumerate}
% subsection 717 (end)

\subsection{Proposition} % (fold)
\label{sub:718}
Sei $\mathcal{H}$ ein Hilbertraum, $z \in \mathcal{H}$, $\alpha, \beta \in \mathds{K}$, $\set{x_i}_{i \in I}, \set{y_i}_{i \in I}$ summierbare Familien. Dann gilt
\begin{align*}
	\alpha \cdot \sum_{i \in I} x_i + \beta \cdot \sum_{i \in I} y_i = \sum_{i \in I} \enbrace*{\alpha \cdot x_i + \beta \cdot y_i} 
	\qquad \text{ und } \qquad 
	\skal*{\sum_{i \in I} x_i}{z} = \sum_{i \in I} \skal*{x_i}{z}
\end{align*}
\minisec{Beweis}
Übung. \bewende
% subsection 718 (end)

\subsection[Proposition: Kriterium für die Summierbarkeit orthogonaler Vektoren]{Proposition} % (fold)
\label{sub:719}
Sei $\set{x_i}_{i \in I}\subset \mathcal{H}$ eine Familie paarweise orthogonaler Vektoren. Dann ist $\set{x_i}_{i \in I}$ summierbar genau dann, wenn 
$\sum_{i \in I} \norm{x_i}^2$ endlich ist. In diesem Fall gilt
\[
	\norm{\sum_{i \in I}x_i}^2 = \sum_{i \in I} \norm{x_i}^2  
\]
\minisec{Beweis}
Falls $\set{x_i}_{i \in I}$ summierbar ist, so gilt
\begin{align*}
	\norm{\sum_{i \in I} x_i}^2 =  
	\skal*{\sum_{i \in I}x_i}{\sum_{j \in I} x_j} \stackrel{\text{\ref{sub:718}}}{=} \sum_{i \in I} \skal*{x_i}{\sum_{j \in I} x_j}\stackrel{\text{\ref{sub:718}}}{=}
	\sum_{i \in I} \enbrace*{\sum_{j \in I} \skal*{x_i}{x_j}} = \sum_{i \in I} \norm{x_i}^2  
\end{align*}
Äquivalenz: Übung mit 
\[
	\norm{\sum_{i \in I}x_i}^2 \xleftarrow{J \finSub I}\norm{\sum_{i \in J} x_i}^2 = \sum_{i \in J} \norm{x_i}^2 \xrightarrow{J \finSub I} \sum_{i \in I} \norm{x_i}^2  
\]
für $J \finSub I$ endlich und \ref{sub:717}. \bewende
% subsection 719 (end)

\subsection{Proposition} % (fold)
\label{sub:720}
Sei $\set{x_i}_{i \in I} \subset \mathcal{H}$ ein Orthonormalsystem, $x \in \mathcal{H}$. 
\begin{enumerate}[(i)]
	\item Es gilt $\sum_{i \in I} \abs*{\skal*{x_i}{x}}^2 \le \norm{x}^2$, 
	(\Index{Besselsche Ungleichung}\footnote{nach Friedrich Wilhelm Bessel, \url{http://de.wikipedia.org/wiki/Friedrich_Wilhelm_Bessel}}).
	\item Es gilt $\sum_{i \in I} \abs*{\skal*{x_i}{x}}^2 = \norm{x}^2$ 
	(\Index{Parselvalsche Gleichung}\footnote{nach Marc-Antoine Parseval, \url{http://de.wikipedia.org/wiki/Marc-Antoine_Parseval}}) 
	genau dann, wenn 
	\[
		x= \sum_{i \in I} \skal*{x_i}{x} \cdot x_i.
	\]
\end{enumerate}
\minisec{Beweis}
\begin{enumerate}[(i)]
	\item Für $J \finSub I$ endlich gilt
	\begin{align*}
		0 \le\norm{ x - \sum_{i \in J} \skal*{x_i}{x} x_i}^2 &= \skal*{x}{x} - \sum_{i \in J} \skal*{\skal*{x_i}{x}  x_i}{x} - 
		\sum_{i \in J} \skal*{x}{\skal*{x_i}{x} x_i} \\
		&\hphantom{= } + \sum_{i \in J} \sum_{j \in J} \skal*{\skal*{x_i}{x} x_i}{\skal*{x_j}{x} x_j} \\
		&= \skal*{x}{x} - \sum_{i \in J} \overline{\skal*{x_i}{x}} \skal*{x_i}{x} - \sum_{i \in J} \skal*{x_i}{x} \overline{\skal*{x_i}{x}}
		+ \sum_{i \in J} \overline{\skal*{x_i}{x}} \skal*{x_i}{x}   
	\end{align*}
	$\Rightarrow \sum_{i \in J} \abs*{\skal*{x_i}{x}}^2 \le \norm{x}^2 < \infty$, also gilt $\sum_{i \in I} \abs*{\skal*{x_i}{x}}^2 \le \norm{x}^2$.
	\item Die Rechnung aus (i) zeigt auch: 
	\[
		\sum_{i \in J} \abs*{\skal*{x_i}{x}}^2 \xrightarrow{J \finSub I} \norm{x}^2  \iff \sum_{i \in J} \skal*{x_i}{x} \cdot x_i \xrightarrow{J \finSub I}  x 
		\bewende
	\]
\end{enumerate}
% subsection 720 (end)

\subsection[Definition und Proposition: Summe von Hilberträumen]{Definition und Proposition} % (fold)
\label{sub:721}
Seien $\mathcal{H}_i$ für $i \in I$ Hilberträume. Dann ist
\[
	\bigoplus_{i \in I} \mathcal{H}_i := \set[(x_i)_{i \in I}]{x_i \in \mathcal{H}_i, \sum_{i \in I} \norm[\mathcal{H}_i]{x_i}^2 < \infty} 
\]
ein Hilbertraum mit $\skal*{(x_i)_{i \in I}}{(y_i)_{i \in I}} := \sum_{i \in I} \skal*{x_i}{y_i}_{\mathcal{H}_i}$.
\minisec{Beweis}
\emph{Übung!} \bewende
% subsection 721 (end)

\subsection{Satz} % (fold)
\label{sub:722}
Sei $\mathcal{H}$ ein Hilbertraum und $\set{x_i}_{i \in I} \subset \mathcal{H}$ ein Orthonormalsystem. Dann sind äquivalent: 
\begin{enumerate}[(i)]
	\item $\set{x_i}_{i \in I}$ ist eine Hilbertraumbasis. 
	\item Falls $x \bot x_i$ für alle $i \in I$, so gilt $x=0$.
	\item Es gibt einen isometrischen Isomorphismus $\bigoplus_{i \in I} \mathcal{H}_i \cong \mathcal{H}$ 
	\[
		\enbrace*{\lambda_i \cdot x_i}_{i \in I} \longmapsto \sum_{i \in I} \lambda_i \cdot x_i 
	\]
	wobei $\mathcal{H}_i = \mathds{K} \cdot x_i$ der von $x_i$ erzeugte Unterhilbertraum ist.
	\item Für jedes $x \in \mathcal{H}$ gilt $x= \sum_{i \in I} \skal*{x_i}{x} \cdot x$.
	\item Für alle $x,y \in \mathcal{H}$ gilt $\skal*{x}{y} = \sum_{i \in I} \skal*{x}{x_i} \cdot \skal*{x_i}{y}$.
	\item Für alle $x \in \mathcal{H}$ gilt die Parsevalsche Gleichung $\norm{x}^2 = \sum \abs*{\skal*{x_i}{x}}^2$.
\end{enumerate}
\minisec{Beweis}
\begin{description}
	\item[(i)$\Rightarrow$(ii):] Falls $0 \not= x \bot x_i$ für $i \in I$, so ist $\set{x_i}_{i \in I} \cup \set{\frac{1}{\norm{x}} \cdot x}$ ein Orthonormalsystem. \light
	\item[(ii)$\Rightarrow$(iii):] $\gamma \colon \bigoplus_{i \in I} \mathcal{H}_i \to \mathcal{H}$, $(\lambda_i \cdot x_i)_I \mapsto \sum \lambda_i \cdot x_i$ ist eine
	wohldefinierte lineare Isometrie nach \ref{sub:719}. Außerdem erhält $\gamma$ das Skalarprodukt. $\gamma \enbrace*{\bigoplus_{i \in I} \mathcal{H}_i}$ ist vollständig,
	also abgeschlossen in $\mathcal{H}$. Falls $\gamma \enbrace*{\bigoplus_{i\in I} \mathcal{H}_o} \not= \mathcal{H}$, so ist nach \ref{sub:715} 
	$\gamma \enbrace*{\bigoplus_{i \in I} \mathcal{H}_i}^\bot \not \set{0}$. \light zu (ii). 
	\item[(iii)$\Rightarrow$(iv):] Für $x \in \mathcal{H}$ existieren $\lambda_i \in \mathds{K}$,  sodass 
	$x= \gamma \enbrace[\big]{(\lambda_i \cdot x_i)_{i \in I}} = \sum_{i \in I} \lambda_i \cdot x_i$ gilt. Es gilt 
	\begin{align*}
		\skal*{x_j}{x} = \skal*{x_j}{\sum_{i \in I} \lambda_i \cdot x_i} = \sum_{i \in I} \skal*{x_j}{\lambda_i \cdot x_i} = \sum_{i\in I} \lambda_i \cdot \delta_{ij} = \lambda_j
	\end{align*}
	Damit folgt (iv).
	\item[(iv)$\Rightarrow$(v):] Klar.
	\item[(v)$\Rightarrow$(vi):] Klar.
	\item[(vi)$\Rightarrow$(i):] Falls $x \bot x_i$ für alle $i \in I$, so gilt $\norm{x}^2= \sum_{i \in I} \abs*{\skal*{x_i}{x}}^2 =0$. Also ist $x=0$ und damit das 
	Orthonormalsystem maximal. \bewende
\end{description}
% subsection 722 (end)

\subsection{Bemerkung} % (fold)
\label{sub:723}
\begin{enumerate}[(i)]
	\item Jeder Hilbertraum besitzt eine Basis (Lemma von Zorn).
	\item Wegen \ref{sub:722} (iii) ist ein Hilbertraum bis auf isometrischen Isomorphismus durch die Mächtigkeit seiner Basis eindeutig bestimmt.
	\item Jeder separable Hilbertraum ist isometrisch isomorph zu $\mathds{K}^n$ für ein $n \in \mathds{N}$ oder zu $\ell^2(\mathds{N})$.
\end{enumerate}
% subsection 723 (end)

\subsection[Beispiele für Basen von Hilberträumen]{Beispiele} % (fold)
\label{sub:724}
\begin{enumerate}[(i)]
	\item $\set{e_i}_{i \in \mathds{N}}$ bzw. $\set{e_i}_{i \in \mathds{Z}}$ bilden Basen von $\ell^2(\mathds{N})$ bzw. $\ell^2(\mathds{Z})$, wobei 
	$e_i(j)= \delta_{ij}$ ist.
	\item $\enbrace*{e^{2 \pi i n t}}_{n \in \mathds{N}}$ bilden eine Basis von $L^2([0,1])$.
\end{enumerate}
% subsection 724 (end)

\subsection[Definition und Proposition: Der zu $T \in \mathcal{L}(\mathcal{H}_1, \mathcal{H}_2)$ adjungierte Operator $T^*$]{Definition und Proposition} % (fold)
\label{sub:725}
Seien $\mathcal{H}_1$, $\mathcal{H}_2$ Hilberträume und $T \in \mathcal{L}(\mathcal{H}_1, \mathcal{H_2})$. Dann existiert genau ein Operator 
$T^* \in \mathcal{L}(\mathcal{H}_2, \mathcal{H}_1)$, sodass für alle $x \in \mathcal{H}_1$, $y \in \mathcal{H}_2$ gilt 
\begin{equation*}
	\skal*{T^* y}{x}_{\mathcal{H}_1} = \skal*{y}{T x}_{\mathcal{H}_2} \label{eq:725} \tag{\#}
\end{equation*}
Die Abbildung ${ }^*  \colon \mathcal{L}(\mathcal{H}_1, \mathcal{H}_2) \to \mathcal{L}(\mathcal{H}_2, \mathcal{H}_1)$ ist konjugiert linear und isometrisch. Weiter gilt
$T^{**} = T$. Ist $\mathcal{H}_3$ ein weiterer Hilbertraum und $S \in \mathcal{L}(\mathcal{H}_2, \mathcal{H}_3)$, so gilt
\[
	(S \circ T)^* = T^* \circ S^*
\]
\minisec{Beweis}
Sei $\kappa_i \colon \mathcal{H}_i \to \mathcal{H}_i^*$ wie in \ref{sub:711}. Wir haben 
\begin{align*}
	& \skal*{T^* y}{x}_{\mathcal{H}_1} = \skal*{y}{T x}_{\mathcal{H}_2} & \tag*{[\textbf{\#}]} \\
	\iff & \kappa_1(T^* y)(x) =\kappa_2(y)(T x) &\text{ für } x \in \mathcal{H}_1, y \in \mathcal{H}_2 \\
	\iff & \kappa_1(T^* y) =\kappa_2(y)\circ T & \text{ für } y \in \mathcal{H}_2\\
	\iff & \kappa_1(T^* y) = T^\tr \circ \kappa_2(y) &\text{ für } y \in \mathcal{H}_2\\
	\iff & T^* y = \kappa_1 ^{-1} \circ T^\tr \circ \kappa_2(y) &\text{ für } y \in \mathcal{H}_2
\end{align*}
Also definieren wir $T^* := \kappa_1 ^{-1} \circ T^\tr \circ \kappa_2 \in \mathcal{L}(\mathcal{H}_2, \mathcal{H}_1)$. Durch \eqref{eq:725} ist $T^*$ eindeutig bestimmt.
$\kappa_1 ^{-1}$ und $\kappa_2$ sind Isometrien; da außerdem $\norm{T^\tr}= \norm{T}$, gilt $\norm{T^*}=\norm{T}$ und somit ist $T \mapsto T^*$ isometrisch.
$\kappa_1 ^{-1}$ ist konjugiert linear, also auch $T \mapsto T^*$. Weiter gilt
\begin{align*}
	(S T)^* = \kappa_1 ^{-1} (S T)^\tr \kappa_3 = \kappa_1 ^{-1} T^\tr S^\tr \kappa_3 = \kappa_1 ^{-1} T^\tr \kappa_2 \kappa_2 ^{-1} S^\tr \kappa_3 = T^* S^*
\end{align*}
Des weiteren erfüllt $T^{**} = (T^*)^* \in \mathcal{L}(\mathcal{H}_1, \mathcal{H}_2)$ für alle $x \in \mathcal{H}_1, y \in \mathcal{H}_2$
\[
	\skal*{T^{**} x}{y}_{\mathcal{H}_2} = \skal*{x}{T^* y}_{\mathcal{H}_1}
\]
Ebenso gilt $\skal*{T x}{y}_{\mathcal{H}_2}= \overline{\skal*{y}{T_x}_{\mathcal{H}_2}} = \overline{\skal*{T^* y}{x}_{\mathcal{H}_1}} = \skal*{x}{T^* y}_{\mathcal{H}_1}$ für
alle $x \in \mathcal{H}_1$, $y \in \mathcal{H}_2$. Folglich muss $T^{**} x = T x$ gelten.\marginnote{gleiche Werte auf allen Elementen, Hahn-Bannach} \bewende 
% subsection 725 (end)

\subsection{Proposition} % (fold)
\label{sub:726}
Für $T \in \mathcal{L}(\mathcal{H}_1, \mathcal{H}_2)$ gilt: 
\begin{align*}
	\norm{T} &= \inf \set[\lambda \in \mathds{R}]{ \abs*{\skal*{y}{T x}_{\mathcal{H}_2}} \le \lambda \cdot \norm{y} \cdot \norm{x} \forall x \in \mathcal{H}_1, y \in \mathcal{H}_2}  \\
	&= \sup \set[\abs*{\skal*{y}{T x}}_{\mathcal{H}_2}]{ x \in \overline{B}_{\mathcal{H}_1}(0,1), y \in \overline{B}_{\mathcal{H}_2}(0,1)} 
\end{align*}
\minisec{Beweis}
Sei $\gamma$ die rechte Seite der oberen Gleichung.
Es gilt $\abs*{\skal*{y}{Tx}_{\mathcal{H}_2}} \le \norm{T} \cdot \norm{y} \cdot \norm{x}$ für $x \in \mathcal{H}_1$, $y \in \mathcal{H}_2$, also ist 
$\gamma \le \norm{T}$. Weiter gilt 
\[
	\norm{T x}^2 = \abs*{\skal*{T x}{T x}} \le \gamma \norm{T x} \cdot \norm{x} \le \gamma \cdot \norm{T} \cdot \norm{x}^2
\]
also ist $\norm{T x} \le \sqrt{\gamma} \cdot \sqrt{\norm{T}} \cdot \norm{x}$ für $x \in \mathcal{H}_1$ und es folgt $\norm{T} \le \sqrt{\gamma} \cdot \sqrt{\norm{T}}$. Also
gilt $\sqrt{\norm{T}} \le \sqrt{\gamma}$ $\Rightarrow $ $\norm{T} \le \gamma$ und somit insgesamt $\gamma = \norm{T}$.
Die zweite Gleichheit ist klar. \bewende
% subsection 726 (end)

\subsection{Proposition} % (fold)
\label{sub:727}
Sei $\mathcal{H}$ ein $\mathds{K}$-Hilbertraum, $T\in \mathcal{L}(\mathcal{H})$ und $\beta\ge 0$. Falls $\abs*{\skal*{x}{T x}}\le\beta\norm{x}^2$ für alle $x\in\mathcal{H}$,
so gilt für $x,y \in \mathcal{H}$
\begin{align*}
	\abs*{\skal*{y}{T x} + \skal*{T y}{x}} &\le 2 \beta \cdot \norm{x} \cdot \norm{y} \\
	\shortintertext{Falls $\mathds{K}=\mathds{C}$, gilt sogar}
	\abs*{\skal*{y}{T x}} + \abs*{\skal*{T y}{x}} &\le 2 \beta \cdot \norm{x} \cdot \norm{y}
\end{align*}
In diesem Fall gilt also $T=0 \iff \skal*{x}{T x}=0$ für alle $x \in \mathcal{H}$.
\minisec{Beweis}
Für alle $x,y \in \mathcal{H} \setminus \set{0} $ gilt (für $0$ ist die Ungleichung trivial)
\begin{align*}
	2 \abs*{\skal*{y}{T x}+ \skal*{T y}{x}} = \abs*{\skal*{x+y}{T(x+y)} - \skal*{x-y}{T(x-y)}} &\le \beta \norm{x+y}^2 + \beta \norm{x-y}^2 \\
	&\le 2 \beta \enbrace*{\norm{x}^2 + \norm{y}^2}
\end{align*}
Ersetze $x$ durch $\enbrace*{\frac{\norm{y}}{\norm{x}} }^{\frac{1}{2}} \!\cdot x$ und $y$ durch $\enbrace*{\frac{\norm{x}}{\norm{y}}}^{\frac{1}{2}}\! \cdot y$. Dann folgt
\begin{equation*}
	\bcancel{2}\cdot  \abs*{\skal*{y}{T x} + \skal*{T y}{x}} \le \bcancel{2} \cdot \beta \cdot \enbrace[\Big]{\norm{x} \cdot \norm{y} + \norm{x} \cdot \norm{y}} 
	= \bcancel{2} \cdot 2 \cdot \beta \cdot \norm{x} \cdot \norm{y} \label{eq:727} \tag{*}
\end{equation*}
Falls $\mathds{K}=\mathds{C}$, so gilt für $\rho, \sigma \in \mathds{R}_+$ geeignet \todo{das stimmt noch nicht so ganz \ldots}
\begin{align*}
	\abs*{\skal*{y}{T x}} + \abs*{\skal*{T y}{x}} &= \abs*{ e^{i(\rho +\sigma)} \cdot \skal*{y}{T x} + e^{i(\rho-\sigma)} \cdot {\skal*{x}{T y}}} \\
	&= \abs*{e^{i \rho} \cdot \enbrace*{\skal*{y}{T \enbrace*{e^{i \sigma} \cdot x} } + \skal*{e^{i \sigma} \cdot x}{T y} } } \\
	&= \abs*{\skal*{y}{T \enbrace*{e^{i \sigma} \cdot x} } + \overline{\skal*{y}{T \enbrace*{e^{i \sigma} \cdot x}}}} \\
	\ldots & \stackrel{\mathclap{\mbox{\scriptsize\eqref{eq:727}}}}{\le} 2 \cdot \beta \cdot \norm{x} \cdot \norm{y}  \bewende
\end{align*}
% subsection 727 (end)

\subsection[Definition: Selbstadjungierter Operator]{Definition} % (fold)
\label{sub:728}
$T \in \mathcal{L}(\mathcal{H})$ heißt \Index{selbstadjungiert}, falls $T=T^*$. Wir setzen
\[
	\mathcal{L}(\mathcal{H})_\sa := \set[T \in \mathcal{L}(\mathcal{H})]{T \text{ selbstadjungiert}} 
\]
% subsection 728 (end)

\subsection{Bemerkung} % (fold)
\label{sub:729}
Für $T$ selbstadjungiert gilt 
\[
	\norm{T}= \inf \set[\lambda \in \mathds{R}]{\abs*{\skal*{x}{T x}} \le \lambda \norm{x}^2 \text{ für alle } x \in \mathcal{H}}  
\]
\minisec{Beweis}
Sei $\beta$ das Infimum auf der rechten Seite. Dann gilt
\begin{align*}
	2 \cdot \abs*{\skal*{x}{T x}} = \abs*{\skal*{y}{T x} + \skal*{T y}{x}} \stackrel{\mbox{\scriptsize \ref{sub:727}}}{\le} 2 \cdot \beta \norm{x} \cdot \norm{y} 
\end{align*}
Nach \ref{sub:726} folgt $\norm{T} \le \beta \le \norm{T}$. \bewende
% subsection 729 (end)

\subsection{Proposition} % (fold)
\label{sub:730}
Falls $\mathds{K}=\mathds{C}$ und $T \in \mathcal{L}(\mathcal{H})$, so gilt
\[
	T \text{ selbstadjungiert } \iff \skal*{x}{T x} \in \mathds{R} \text{ für alle } x \in \mathcal{H}
\]
\minisec{Beweis}
Um "$\Rightarrow$" zu zeigen, betrachte $\skal*{x}{T x} = \skal*{T x}{x} = \overline{\skal*{x}{T x}}$ für $x \in \mathcal{H}$. Für die Rückrichtung betrachte
$\skal*{x}{T x} = \skal*{T x}{x} = \skal*{x}{T^* x}$. Dann ist $\skal*{x}{(T-T^*)x}=0$ für $x \in \mathcal{H}$. Es folgt dann $T- T^*=0$. \bewende 
% subsection 730 (end)

\subsection{Definition und Proposition} % (fold)
\label{sub:731}
Seien $R,S,T \in \mathcal{L}(\mathcal{H})$ selbstadjungiert. Wir schreiben 
\begin{itemize}
	\item $T \ge 0$, falls $\skal*{x}{T x} \ge 0$ für alle $x \in \mathcal{H}$ gilt.
	\item $T \ge S$, falls $T-S \ge 0$
\end{itemize}
"$\ge$" ist eine partielle Ordnung auf $\mathcal{L}(\mathcal{H})_\sa$, das heißt $T \ge T$; $T \ge S$, $S \ge T \Rightarrow S=T$, $T \ge S, S \ge R \Rightarrow T \ge R$ (folgt mit \ref{sub:729}).
Es gilt $-\norm{T} \id_{\mathcal{H}} \le T  \le \norm{T} \cdot \id_{\mathcal{H}}$. Falls $T \ge 0$ ist, so gilt 
\begin{equation*}
	\abs*{\skal*{y}{T x}}^2 \le  \skal*{x}{T x} \cdot \skal*{y}{T y} \label{eq:731} \tag{*}
\end{equation*}
insbesondere gilt $\skal*{x}{T x}=0 \iff T x=0$
\minisec{Beweis}
\emph{Übung, für die letzte Aussage benutze \ref{sub:74} und: $(y,x) \mapsto \skal*{y}{T x}$ ist positiv semidefinit hermitesch.}
% subsection 731 (end)

\subsection{Satz} % (fold)
\label{sub:732}
Sei $(T_j)_{j \in I} \subset \mathcal{L}(\mathcal{H})_\sa$ ein monoton wachsendes und beschränktes Netz selbstadjungierter Operatoren. Dann konvergiert $(T_j)_{j \in I}$
punktweise gegen ein $T \in \mathcal{L}(\mathcal{H})_\sa$, das heißt $T_j x \to T x$ in Norm für jedes $x \in \mathcal{H}$. Wir schreiben auch $T_j \xrightarrow{\so} T$,
d.h. $(T_j)_{j \in I}$ konvergiert gegen $T$ in der \bet{starken Operatortopologie}\index{starke Operatortopologie}.
\minisec{Beweis}
Nach Polarisierung (vergleiche \hyperref[sub:78]{Bemerkung \ref*{sub:78}}) gilt (für $\mathds{K}=\mathds{C}$)
\[
	\skal[\big]{x}{T_j y} = \frac{1}{4} \enbrace[\Big]{\skal[\big]{x+y}{T_j(x+y)} + \skal[\big]{x-y}{T_j(x-y)} + i \skal[\big]{i x +y}{T_j(ix +y)} 
	- i \skal[\big]{i x -y}{T_j(ix-y)}}
\]
Die rechte Seite konvergiert (warum?), also auch die linke Seite. $\Rightarrow$ Für jedes $x \in \mathcal{H}$ konvergiert 
$\skal*{x}{T_j(\cdot)} \colon \mathcal{H} \to \mathds{C}$ gegen ein $\varphi_x \colon \mathcal{H} \to \mathds{C}$. Aber $\varphi_x \in \mathcal{H}^*$ (warum?). Es gilt
\begin{equation*}
	\norm[\mathcal{H}^*]{\varphi_x} \le \limsup_j \norm{T_j} \norm{x} < \infty. 
\end{equation*}
Die Abbildung $x \mapsto \kappa ^{-1}(\varphi_x)$ ist linear (warum?). Definiere 
$T \colon \mathcal{H} \to \mathcal{H}$ durch $T(x) := \kappa ^{-1}(\varphi_x)$. Es gilt $\norm{T(x)} = \norm{\varphi_x}$, also $\norm{T}\le\limsup_j \norm{T_j} \le \infty$.
Damit ist $T \in \mathcal{L}(\mathcal{H})$. Außerdem gilt für $x,y \in \mathcal{H}$
\begin{align*}
	\skal*{T x}{y} = \kappa(T x)(y) = \varphi_x(y) = \lim_j \skal[\big]{x}{T_j y} = \lim_j \skal*{T_j x}{y} = \overline{\lim_j \skal*{y}{T_j x}} 
	= \overline{\varphi_y(x)} &= \overline{\skal*{T y}{x}}\\  &= \skal*{x}{T y}  
\end{align*}
Also ist $T$ auch selbstadjungiert. Es gilt für $\lim_j \skal*{x}{T_j y} = \skal*{x}{T y}$ für $x,y \in \mathcal{H}$, das heißt $T_j \xrightarrow{\mathrm{w.o.}} T$
($T_j \to T$ in der schwachen Operatortopologie). Für $j \ge j'$ gilt $T_j - T_{j'} \ge 0$, also nach \ref{sub:731} für $x \in \mathcal{H}$
\begin{align*}
	\skal[\Big]{\enbrace*{T_j - T_{j'}}(x)}{\enbrace*{T_j - T_{j'}}(x)} &\le \skal[\Big]{x}{\enbrace*{T_j - T_{j'}}(x) } \cdot 
	\skal*{\enbrace*{T_j - T_{j'}}(x)}{\enbrace*{T_j - T_{j'}}^2(x) } \\
	&\le \skal[\Big]{x}{\enbrace*{T_j - T_{j'}}(x)} \cdot \enbrace*{2 \cdot \limsup_j \norm{T_j} }^3 \cdot \norm{x}^2 <  \varepsilon
\end{align*}
falls $j'$ groß genug ist. Also $\norm{T_j(x-)- T_{j'}(x)}< \varepsilon$ falls $j \ge j'$ groß. Damit folgt dann $T_j(x) \to T(x)$ in Norm (warum?). \bewende
% subsection 732 (end)
% section 7 (end)
\newpage

\section{Kompakte Operatoren und ein Spektralsatz} % (fold)
\label{sec:8}

\subsection[Definition: Spektrum eines Bachnachraumes]{Definition} % (fold)
\label{sub:81}
Sei $X$ ein $\mathds{K}$-Banachraum und $T \in \mathcal{L}(X)$. Wir definieren das \Index{Spektrum} von $T$ als \todo{RevChap8}
\[
	\sigma(T) := \set[\lambda \in \mathds{K}]{T- \lambda \cdot \id_X \text{ ist nicht invertierbar}} \subset \mathds{K} \marginnote{auch $\spec(T)$} 
\]
Die \Index{Resolventenmenge} ist $\res(T) :=  \rho(T) := \mathds{K} \setminus \sigma(T)$. Wir schreiben $\resolv(T,\lambda) := (T- \lambda \cdot \id_X) ^{-1}$ für
die \Index{Resolvente} von $T$ in $\lambda \in \rho(T)$. 
% subsection 81 (end)

\subsection[Bemerkung zum Zusammenhang von Spektralwerten und Eigenwerten]{Bemerkung} % (fold)
\label{sub:82}
\begin{enumerate}[(i)]
	\item Jeder Eigenwert von $T$ ist ein Spektralwert.
	\item Für $X$ endlichdimensional gilt auch die Umkehrung, für beliebige $X$ jedoch nicht: Sei $X=\ell^2(\mathds{N})$, dann ist $S \in \mathcal{L}\enbrace[\big]{\ell^2(\mathds{N})}$
	gegeben durch $S(x_0, x_i, \ldots) := (0,x_0, x_1, \ldots)$ (\Index{unilateraler Shift}) linear, stetig und sogar isometrisch, aber \emph{nicht} surjektiv. Dann ist
	$0 \in \sigma(S)$, aber $0$ ist kein Eigenwert, da $S$ injektiv ist.
\end{enumerate}
% subsection 82 (end)

\subsection{Proposition} % (fold)
\label{sub:83}
Seien $X,Y$ Banachräume und $T \in \mathcal{L}(X,Y)$ bijektiv (also injektiv). Falls für $S \in \mathcal{L}(X,Y)$ gilt 
\[
	\norm[\mathcal{L}(X,Y)]{T-S} < \norm[\mathcal{L}(Y,X)]{T^{-1}}^{-1}, 
\]
so ist auch $S$ bijektiv.
\minisec{Beweis}
Es gilt $S= T \enbrace[\big]{\id_X - T ^{-1}(T-S)}$. Für $k \le l$ gilt mit $C := \norm{T ^{-1}} \cdot \norm{T-S} < 1$ 
\begin{align*}
	\norm{\sum_{n=k}^{l} \enbrace[\big]{T^{-1}(T-S)}^n } \le \sum_{n=k}^{l} \norm{\enbrace[\big]{T ^{-1}(T-S)}^n} \le \sum_{n=k}^{l} \norm{T ^{-1}} \cdot \norm{T-S}^n
	= \sum_{n=k}^{l} C^n      
\end{align*}
also ist $\sum_{n=0}^{\infty} \enbrace[\big]{T ^{-1}(T-S)}^n$ konvergent in $\mathcal{L}(X,X)$. Weiter ist 
$\enbrace*{\sum_{n=0}^{\infty} \enbrace[\big]{T ^{-1}(T-S)}^n} T ^{-1} \in \mathcal{L}(Y,X)$ das Inverse von $S$ (warum?). \bewende
% subsection 83 (end)

\subsection{Corollar} % (fold)
\label{sub:84}
Für $T \in \mathcal{L}(X)$ ist die Resolventenmenge offen in $\mathds{K}$.
% subsection 84 (end)
% section 8 (end)

\cleardoubleoddemptypage
\appendix
\section{Anhang} % (fold)
\label{sec:anhang}

\subsection{Vierecksungleichung} % (fold)
\label{sub:vier_ungl}
Sei $(X,d)$ ein metrischer Raum und $x,y,u,v \in X$. Dann gilt
\[
	\abs[\big]{d(x,y) - d(u,v)} \le d(x,u) + d(y,v) 
\]
\minisec{Beweis}
Einerseits gilt nach der Dreiecksungleichung
\[
	d(x,y) \le d(x,u) + d(u,v) + d(v,y) \quad \Longrightarrow \quad d(x,y) - d(u,v) \le d(x,u) + d(y,v)
\]
Andererseits aber auch
\[
	d(u,v) \le d(u,x) + d(x,y) + d(y,v) \quad \Longrightarrow \quad d(u,v) - d(x,y) \le d(x,u) + d(y,v)
\]
Insgesamt folgt also die Behauptung. \bewende
% subsection vier_ungl (end)

% section anhang (end)
\cleardoubleoddemptypage
\pagenumbering{Alph}
\setcounter{page}{1}
\printindex
\listoffigures
\todototoc
\listoftodos[Todo's und andere Baustellen]
\end{document}